\documentclass[11pt]{article}
\usepackage[margin=1in]{geometry}
\usepackage{amsmath,amssymb,amsthm,bm}
\usepackage{booktabs,tabularx,longtable}
\usepackage{graphicx}
\usepackage{enumitem}
\usepackage{natbib}
\usepackage[hidelinks]{hyperref}
\newtheorem{definition}{Definition}
\newtheorem{theorem}{Theorem}
\newtheorem{proposition}{Proposition}
\newtheorem{corollary}{Corollary}
\title{Conditional Assembly of Gravity, Matter, and Abelian Gauge Terminals\\from One Ordered Tau Descent\\\large Tau Core Foundation Paper VII: First Proof Packet}
\author{Jozsef Olcsak}
\date{August 2026}

\begin{document}
\maketitle

\begin{abstract}
We ask whether the previously constructed Tau gravity, temporal, quantum,
matter and Abelian gauge terminals can descend from one ordered physical
source ledger rather than from independently fitted sector models.  On a
stabilized morphological body we construct one once-counted post-body
effective functional whose effect, metric, gauge and clock-character
derivatives generate probabilities, stress, current and formal clock
response.  Diffeomorphism and $U(1)$ gauge invariance yield the conditional
Einstein--Maxwell--matter equations, Ward/Bianchi identities and their
standard weak-field leaf.  A disjoint occupation ledger prevents duplicate
classical/quantum stress, and exact countermodels show that standard-limit
agreement does not establish common Tau provenance.

On an occupied two-mode restriction, ROOT one-particle coherence and
transverse-traceless stress are proved to arise from one preparation without
an additional cross-terminal gain.  Differential recovery is reduced to a
source-rank problem: after de-duplication the ROOT increment over the common
P1 row space is fixed by one source-owned mixed Gram.  For the P3 sector, a
complete pointed exterior seed produces an orthonormal 256-vector
exterior--record orbit, whereas a faithful full-rank nontracial control proves
that rank and equal central weights alone do not fix its metric.  A positive
orthogonally invariant operator-square functional has exactly the P3 Clifford
law as its zero set, and a three-level theorem combines common-source,
algebra/representation and terminal closure inside the declared enriched
completion.

For two parent-relation-free occupied supports, we additionally prove that
connected-support closure rules which preserve component support and admit no
hidden global acceptance constraint make the joint closure set a Cartesian
product.  Its finite delta-function linearization is therefore the standard
tensor-product algebra.  Complete local marginals with a hidden diagonal
constraint and a bridge-creating rewrite are exact controls, so action
additivity or a vanishing mixed Hessian alone is not sufficient.
The same component-local action can also occupy the complete rank-seven P3
packet without conflict: the P3 grammar remains one connected component,
whereas only primitive-incidence-disconnected supports factorize.  An exact
same-narrow-descriptor control shows that the earlier reduct still does not
select this joint law.

The support--port reconstruction reduces the seven typed directions to one
commuting support flag, two ordered cotangent ports, an independent common
unit and a noncollapsing body decoder.  We then state one class-minimal staged
enriched law whose variational minimum realizes this degree-one packet, the
operator-square relation and the pointed cyclic record orbit.  Its
nonnegative action has an explicit simultaneous zero, so global minimization
occupies that zero without a separate occupation postulate.  This closes
formal law-to-packet ownership inside the proposed enriched physical
completion.  An exact countermodel still proves that the earlier narrow
reduct does not entail that law; whether Nature realizes the proposal is an
empirical question rather than another internal implication.  Absolute
calibration from P3 alone, quantized geometry, unrestricted non-Abelian
occupation, the full Standard-Model spectrum and unrestricted ambient-parent
exhaustivity are not derived here.  As a downstream specialization, the paper
also records a class-relative scalar-15/type-II source reconstruction, the
ratio \(M_*/E_P=1/\sqrt{24\pi}\) with one CODATA calibration, and an exact
functional RG law.  These do not establish physical type-II occupation.
\end{abstract}

\begin{center}
\fbox{\begin{minipage}{0.92\textwidth}\small
\textbf{Inputs from earlier papers.} Paper I supplies the body-first descent;
Paper II the typed source support; Paper III recoverable record transport;
Paper IV the parent-realization boundary; Paper V the observer clock and causal
leaf; Paper VI the conditional quantum terminal.  The theory hub supplies the
P1 gravity--time--radiation ordered ledger, the P3 Abelian handoff and the
pointed exterior-seed closure with its full-rank nontracial control, together
with the component-local no-bridge closure law used below.

\textbf{Results owned here.} The morphology-conditioned joint-terminal
assembly, its finite occupation criterion, the joint terminal ledger, the disjoint
occupation rule, the effective-source derivative ledger, the Ward/Bianchi
consistency theorem, the simultaneous-action and duplicate-stress no-go
results, the ROOT--TT common-preparation restriction, and the same-reduct
parent-realization non-identifiability no-go.  The paper also owns the positive operator-square
source-selector theorem, its equal-Gram commuting control, and the staged
enriched-parent ownership theorem.  It additionally owns the support--port
rank-seven reconstruction, its deletion controls and the class-minimal staged
enriched-law realization theorem, and the self-contained Tau-local
standard-composition theorem with its hidden-relation and crossing-rule
controls.  A joint common-law theorem then proves simultaneous LC1--LC4 and
rank-seven occupation inside that enriched class, together with an exact
narrow-reduct non-entailment control.
The paper finally records, as a downstream branch handoff rather than a result
owned by the P3 spine, the class-relative scalar-15/type-II base--seed
reconstruction, its common Planck-scale ratio and its exact functional RG
evaluation law.
\end{minipage}}
\end{center}

\begin{center}
\fbox{\begin{minipage}{0.92\linewidth}\small
\textbf{Atemporal parent-realization terminology.}
Because the Tau parent is atemporal, phrases such as \emph{Nature selects} do
not denote an external agent or a later choice. Their precise meaning is the
internal physical implication
\[
 (B_\tau^{\rm phys},s_U,\mathcal L_{\rm parent})
 \Longrightarrow \mathfrak P,
\]
together with occupation of the support of the resulting packet
$\mathfrak P$. A theorem inside a declared completion proves only
class-relative realization. It does not prove this unrestricted base--seed
arrow, and agreement of finitely many terminal readouts cannot by itself
establish ambient-source exhaustivity. Current-facing statements therefore
use \emph{parent-law realization} and \emph{physical occupation}; any retained
legacy selection language is shorthand for this same distinction.
\end{minipage}}
\end{center}

\section*{Notation and packet map}

\begin{center}
\small
\begin{tabularx}{0.94\textwidth}{@{}p{0.18\textwidth}>{\raggedright\arraybackslash}X@{}}
\toprule
Symbol or packet & Role in this paper\\
\midrule
$P1$ & Gravity--time--radiation/matter differential packet inherited from the
theory hub.\\
ROOT & Occupied two-mode quantum one-particle descriptor.  Its normalized
shape is $\rho_{\rm ROOT}$, while $C=N_O\rho_{\rm ROOT}$ retains occupation
magnitude.\\
TT & Transverse--traceless metric/stress readout paired with ROOT in the
common-preparation restriction.\\
$P3$ & Seven-role primitive algebra/recovery packet used for the typed
incidence and exterior--record closure.\\
PBSE & Retained six-row base--seed/observer response subsystem used for the
transverse-rank control; it is not the full physical carrier.\\
$J_{\rm phys}^{\rm lin}$ & Full parent-side physical incidence on the
linearized source-frozen quotient.\\
$J_{OS}^{\rm irr}$ & Irreducible observer--source incidence not already
generated by the base--seed columns.\\
$K_*$ & Positive source metric/Hessian on the true-null quotient.\\
GNS & Gelfand--Naimark--Segal representation induced by the occupied positive
source state.\\
MVP & The declared minimum viable physical completion; closure inside it is
conditional and is not unrestricted parent-law realization.\\
$\mathrm{Cl}_\tau(H)$ & Finite parent closure set generated from occupied
support $H$ by the declared component-local Tau closure rules.\\
\bottomrule
\end{tabularx}
\end{center}

\section{Question and claim boundary}

Let $M_\tau$ be the stabilized atemporal morphological body. Paper VII asks:
\begin{equation}
\boxed{\text{Can gravity, matter, gauge/radiation and quantum backreaction}
\atop\text{descend from one once-counted ordered physical source ledger?}}
\end{equation}
The word \emph{one} denotes common provenance and calibration. It does not
mean one emitter, one scalar varied at every generative stage, or one terminal
observable. The generative order remains
\begin{equation}
B_\tau^{\rm phys}+s_U\longrightarrow M_\tau
\longrightarrow \mathfrak L_7
\longrightarrow \operatorname{Desc}_O
\longrightarrow \{g,A,\Phi,\rho_O\}.
\end{equation}
No terminal residual may select $M_\tau$ or the ledger.

Even exact closure of the P3 exterior--record subpacket does not by itself
close five deeper physical questions.  Its dimensionless Clifford and trace
relations are invariant under a common change of the parent action/length
calibration; the same P3 reduct admits both classical-metric and
quantized-metric completions; its seven-role algebra does not select the
separate rank-five $3+2$ internal carrier needed for the conditional
$S(U(3)\times U(2))$ construction; complete anomaly-free generation copies
share the same one-generation reduct; and a finite terminal family admits the
loaded hidden-source extension in Eq.~\eqref{eq:p3loadedhiddenextension}.
Thus absolute calibration from P3 alone, quantized geometry, unrestricted
non-Abelian occupation, the full Standard-Model spectrum and ambient-parent
exhaustivity remain distinct upstream claims rather than consequences silently
bundled into P3 closure.  A downstream scalar-15/type-II specialization fixes
one dimensionless mass/Planck ratio and one SI-calibrated value, but it does
not change this P3 claim boundary or predict \(G\).

\section{Inherited-results ledger}

\begin{table}[ht]
\centering\small
\begin{tabularx}{\textwidth}{@{}p{0.20\textwidth}p{0.19\textwidth}Xp{0.18\textwidth}@{}}
\toprule
Result & Owner & Status and assumptions & Use here\\
\midrule
Body-first descent & Paper I & conditional physical MVP & generative order\\
Typed incidence & Paper II & finite source theorem & occupied domains\\
Record recovery & Paper III & exact/approximate conditional theorem & observer terminal\\
Parent-law realization & Paper IV & present reduct does not entail full packet & claim boundary\\
Temporal descent & Paper V & conditional clock/causal leaf & time and weak limit\\
Quantum terminal & Paper VI & conditional complex/effect/QFT packet & $\Gamma_q$ and $\rho_O$\\
Morphology-conditioned readouts & theory hub & common-source assembly for an
occupied complete signature & typed terminal projections\\
Operation-generated occupation & theory hub & cyclic--surjective/root-positive criterion
inside a fixed source representation & block nonvanishing\\
P1 assembly & theory hub & conditional metric/time/radiation ledger & $g$, null radiation, stress\\
P3 Abel handoff & theory hub & conditional $Z_F=p_*K_Y$ & Maxwell/current sector\\
Pointed P3 exterior seed & theory hub & conditional orthonormal 256-vector
root orbit; full-rank nontracial countermodel & representation boundary\\
Support--port P3 synthesis & theory hub & commuting support flag, two ordered
ports, independent unit and smooth decoder; finite audit & rank-seven
source reconstruction\\
Tau-local composition origin & theory hub; reproved here & exact inside the
component-local no-bridge/no-hidden-global-relation class & two-leg record
and effect tensor algebra\\
Component-local P3 joint occupation & owned here & exact inside one
once-counted enriched common law; narrow-reduct selection disproved & joint
LC1--LC4 and rank-seven status\\
Scalar-15/type-II completion & theory hub P4 branch & class-relative
base--seed reconstruction, common-scale ratio and exact functional RG law;
unrestricted branch occupation open & downstream particle-sector handoff\\
\bottomrule
\end{tabularx}
\caption{Dependency and ownership ledger. No inherited conditional result is
promoted to an unrestricted physical parent-law theorem.}
\label{tab:dependencyledger}
\end{table}

For logical self-containment, Appendix~\ref{app:dependencycertificate}
restates every imported premise actually used below.  The paper labels in
Table~\ref{tab:dependencyledger} are provenance locators, not hidden
hypotheses: the conditional
arguments can be checked from the restated premises without accepting an
unquoted result from a companion manuscript.  Several listed papers provide
architectural context only and are not logical dependencies of the main
theorems.

\section{Morphology-conditioned assembly and occupation}

The inherited terminals are not treated as independent information channels.
Let the stabilized body and observer support determine one finite physical
source representation $\mathcal E_{*,O}$.  No direct-sum decomposition of
this source is assumed.  Instead, for the typed role set
\begin{equation}
\mathcal P=\{T,G,R,M,Q,A\},
\qquad
\mathcal Y_{\rm role}:=\bigoplus_{a\in\mathcal P}\mathcal Y_a,
\label{eq:rolecodomain}
\end{equation}
let $R_a:\mathcal E_{*,O}\to\mathcal Y_a$ be terminal-role maps and
$R=(R_a)_a:\mathcal E_{*,O}\to\mathcal Y_{\rm role}$ their stacked map.
The direct sum is therefore a de-duplicated \emph{role codomain}, not a claim
that the source information separates.  Common source directions may affect
several $R_a$ simultaneously; their overlap is measured by the joint kernel
and stacked rank of the $R_a$.

Let $A_O$ be the body- and observer-conditioned access map and let
$Q_O^{(a)}$ be a stable observer-relative resolution map.  The operational
terminals are
\begin{equation}
Y_O^{(a)}=Q_O^{(a)}R_a A_O\Xi_*,
\label{eq:typedterminals}
\end{equation}
for one occupied source representative $\Xi_*$.  Equality under the specified
$Q_O^{(a)}$, rather than a non-transitive distance threshold, defines the
resolved readout classes.  Thus observer-relative quantization remains part
of the full descent and is not replaced by the continuous source geometry.

\begin{proposition}[Conditional morphology-conditioned representation]
\label{prop:moprassembly}
Assume a stabilized body, one occupied complete source representation
$\mathcal E_{*,O}$, a source-owned typed ledger map, and stable resolution
maps $Q_O^{(a)}$.  Then all terminals in Eq.~\eqref{eq:typedterminals} are
functorial restrictions of one source representative.  Their common-source
ownership and once-counting follow without terminal-specific gains.  Different
terminal names do not by themselves imply independent information.  This is
a representation theorem for an already occupied complete source ledger; it
does not establish uniqueness among all admissible source functionals.
\end{proposition}

\begin{proof}
Body stabilization fixes the domain of $A_O$.  The single map $A_O\Xi_*$ is
then mapped by the typed $R_a$ and resolved only afterwards by
$Q_O^{(a)}$.  Consequently every terminal factors through the same occupied
representative and the same access map.  Any additional terminal gain would
be an extra source datum, contrary to the assumed complete ledger.  The
conclusion is role-completeness of this construction, not uniqueness under
field redefinitions, improvement terms or alternative ultraviolet
completions.
\end{proof}

Let $\mathfrak A_{\rm src}$ be the source-owned operation algebra and
$\Omega_s$ the occupied source vector, and define
$V_s:=\operatorname{span}(\mathfrak A_{\rm src}\Omega_s)$.

\begin{definition}[Activation and full occupation]
A required terminal role $a$ is \emph{activated} when $R_a(V_s)\neq\{0\}$.
It is \emph{fully occupied} when
\begin{equation}
R_a(V_s)=\mathcal Y_a.
\label{eq:fulloccupation}
\end{equation}
Full occupation implies activation; the converse need not hold.
\end{definition}

Assembly does not by itself prove either property.  In finite dimension, full
occupation of every required role is equivalent to
\begin{equation}
R_a\!\left(\operatorname{span}(\mathfrak A_{\rm src}\Omega_s)\right)=\mathcal Y_a
\quad\text{for every required }a.
\label{eq:cyclicoccupation}
\end{equation}

\begin{proposition}[Operation-generated occupation criterion]
\label{prop:occupation}
If $\Omega_s$ is cyclic for $\mathfrak A_{\rm src}$ on
$\mathcal E_{*,O}$ and each required $R_a$ is surjective, then
Eq.~\eqref{eq:cyclicoccupation} holds and every required role is fully
occupied.  More weakly, let $X_0\geq0$ be the occupied root covariance with
$\operatorname{supp}X_0\subseteq V_s$, and let the source-owned path maps
$L_e$ preserve the occupied subspace, $L_e(V_s)\subseteq V_s$.  Define
\begin{equation}
X_a:=\sum_{e\in\mathcal P_a}R_aL_eX_0L_e^\dagger R_a^\dagger\geq0.
\label{eq:positivepaths}
\end{equation}
If $R_aL_eX_0^{1/2}\neq0$ for at least one $e\in\mathcal P_a$, then role $a$
is activated, but full rank is not implied.
\end{proposition}

\begin{proof}
Cyclicity gives $V_s=\mathcal E_{*,O}$; surjectivity then gives
$R_a(V_s)=\mathcal Y_a$.  In the weaker case, every summand in
Eq.~\eqref{eq:positivepaths} is positive semidefinite and
\begin{equation}
\operatorname{Tr}X_a
=\sum_{e\in\mathcal P_a}\|R_aL_eX_0^{1/2}\|_{\rm HS}^2.
\end{equation}
One nonzero path therefore implies $X_a\neq0$.  Because
$\operatorname{Ran}(L_eX_0^{1/2})\subseteq V_s$, its nonzero image lies in
$R_a(V_s)$ and hence proves activation.  It does not establish
Eq.~\eqref{eq:fulloccupation}.
\end{proof}

Propositions~\ref{prop:moprassembly} and \ref{prop:occupation} move the open
problem upward.  Once the common representation and its source operation law
are occupied, separate terminal selection is unnecessary.  What remains
unproved is that $B_\tau^{\rm phys}+s_U$ selects this representation, makes
$\Omega_s$ cyclic with surjective required role maps (or realizes the rooted
positive alternative), and chooses
the observer access and resolution maps.  Algebraic generation, terminal
agreement and post-handoff record completeness do not prove that prior
physical realization and occupation.

\section{The ordered Paper-VII ledger}

Let $y$ denote body coordinates, and write $z=(g,A,\Phi)$ for the variational
post-body geometry, Abelian gauge and occupied classical matter fields.  The
quantum preparation is instead a fixed positive trace-one boundary datum
$\rho_{\rm in,O}\geq0$, $\operatorname{Tr}\rho_{\rm in,O}=1$ for the in-in
contour.  It is not an unconstrained stationary field.

\begin{definition}[Once-counted ordered ledger]
The Paper-VII ledger is
\begin{equation}
\mathfrak L_7=\left(
\mathcal A_{\rm pre}[y;B,s],\ y_*,\
\Gamma_7[g,A,\Phi\mid y_*;\rho_{\rm in,O}],\
\operatorname{Desc}_{O\leftarrow y_*},\ \mathcal U
\right),
\end{equation}
where $D_y\mathcal A_{\rm pre}(y_*)=0$, the physical body Hessian is positive,
and
\begin{equation}
\Gamma_7=\Gamma_{\rm grav}[g\mid y_*]
+\Gamma_{\rm EM}[A,g\mid y_*]
+\Gamma_{\rm occ}[\Phi,A,g\mid y_*]
+\Gamma_q[A,g\mid y_*;\rho_{\rm in,O}]+\Gamma_{\rm ct}[A,g].
\label{eq:jointfunctional}
\end{equation}
Each occupied mode and each counterterm appears in exactly one summand.
Observer restriction is evaluated after stationarity and is not varied back
against $y_*$.  The state $\rho_{\rm in,O}$ fixes the CTP preparation; the
Euler equations vary $g$, $A$ and $\Phi$, not $\rho_{\rm in,O}$.  The unit
packet $\mathcal U$ is shared by every summand.
\end{definition}

\section{One effective source on each stationary body fiber}

The symbol $\Gamma_7$ need not be postulated as five unrelated pieces.  Let
$\varphi$ collect the occupied quantum fields, and for each declared outcome
family let the dimensionless constant probe $\eta_a$ couple to a normalized
POVM-effect insertion satisfying
$E_a\geq0$ and $\sum_aE_a=I$ (hence $0\leq E_a\leq I$).  Let $\vartheta$
denote the occupied clock character.  The source convention is chosen so that
\begin{equation}
\Gamma_{\rm PC}[\eta]
=\Gamma_{\rm PC}[0]+\sum_a\eta_a\,\mathsf E_a[\varphi],
\qquad
\langle\mathsf E_a\rangle=\operatorname{Tr}(\rho_OE_a).
\label{eq:effectcoupling}
\end{equation}
On the already selected stationary body $y_*$ define the doubled
closed-time-path influence functional
\begin{align}
\mathcal I_{\rm CTP}
&:=\Gamma_{\rm PC}[y_*;g^+,A^+,\vartheta^+,\eta^+,\varphi^+]\notag\\
&\quad-\Gamma_{\rm PC}[y_*;g^-,A^-,\vartheta^-,\eta^-,\varphi^-],
\label{eq:ctpaction}\\
\exp\!\left({i\over\hbar}W_O[g^\pm,A^\pm,\vartheta^\pm,\eta^\pm
\mid y_*;\rho_{\rm in,O}]\right)
&=\int_{\rho_{\rm in,O}}\!\mathcal D\varphi^+\mathcal D\varphi^-\,
\exp\!\left({i\over\hbar}\mathcal I_{\rm CTP}\right).
\label{eq:ineffective}
\end{align}
The initial boundary condition joins the two branches through the fixed
preparation $\rho_{\rm in,O}$, and the path integral includes the standard
final branch identification.  The in-in contour is required for real causal expectation
values \citep{Schwinger1961,Keldysh1965}.  For every doubled source write
$X_c=(X^++X^-)/2$ and $X_\Delta=X^+-X^-$.  All derivatives below mean
$\delta/\delta X_\Delta$ followed by the physical limit
$X_\Delta=0$, $X_c=X$; differentiating an already identified common branch
would instead lose the response by CTP normalization.  The construction is a
terminal representation of an atemporal ordered functional fibration, not a
parent-time history.  CTP evolution and observer restriction produce the
normalized output state $\rho_O$ used in Eq.~\eqref{eq:effectcoupling}; neither
$\rho_{\rm in,O}$ nor $\rho_O$ is varied as a local Euler field.
Renormalization gives one $W_O^{\rm ren}$ whose physical-limit difference
derivatives define
\begin{align}
p_a&={\delta W_O^{\rm ren}\over\delta\eta_a}\bigg|_{\eta=0}
=\operatorname{Tr}(\rho_OE_a)\in[0,1],
&\langle T^q_{\mu\nu}\rangle_{\rm ren}
&=-{2\over\sqrt{-g}}{\delta W_O^{\rm ren}\over\delta g^{\mu\nu}},
\label{eq:sourceowners1}\\
\langle J_q^\mu\rangle_{\rm ren}
&=-{1\over\sqrt{-g}}{\delta W_O^{\rm ren}\over\delta A_\mu},
&\mathcal R_O^{\rm clock}&={\delta W_O^{\rm ren}\over\delta\vartheta}.
\label{eq:sourceowners2}
\end{align}
Consequently $\sum_a p_a=1$ within each normalized outcome family.  Without
the POVM and source-coupling conditions, the derivative would be only a
generic operator expectation and could not be identified as a Born
probability.
Here the effect-source renormalization convention is required to preserve the
POVM normalization at $\eta=0$; Eq.~\eqref{eq:effectcoupling} is an admitted
operational insertion and not a claim that every POVM has a local polynomial
field representation.
The final expression is schematic until the clock character and its
normalization are fixed, but its ownership is unambiguous.
The minus sign in the current definition is fixed by the Maxwell convention
$\Gamma_{\rm EM}=-(Z_F/4)\int\sqrt{-g}\,F^2$: stationarity then gives
$\nabla_\mu(Z_FF^{\mu\nu})=J^\nu$.  We use the same sign convention for the
classical current,
$J_{\rm cl}^\mu:=-(\sqrt{-g})^{-1}\delta\Gamma_{\rm occ}/\delta A_\mu$.

\begin{theorem}[Effective common-source construction]
\label{thm:effectivesource}
Assume the ordered pre-body solve, a source-owned probe-complete post-body
functional $\Gamma_{\rm PC}$, an observer-admissible integration class
$\mathcal C_O$ with fixed normalized boundary state
$\rho_{\rm in,O}$, and a local covariant, anomaly-free renormalization
prescription for the symmetries used below.  Then
Eqs.~\eqref{eq:ineffective}--\eqref{eq:sourceowners2}, together with the
classical body and coherent mean-field terms, construct the Paper-VII
$\Gamma_7$ ledger from one post-body source functional.  The construction is
complete at the level of the declared physical roles, modulo relocation of a
finite local counterterm between the geometric and quantum bookkeeping
columns while holding their sum fixed.  No uniqueness claim is made across
field redefinitions, improvement terms or inequivalent microscopic
functionals.
\end{theorem}
\begin{proof}
Functional differentiation of one $W_O^{\rm ren}$ gives the four displayed
terminal owners.  Diffeomorphism and gauge invariance give their common Ward
covector.  The fluctuation Green operator is the inverse of the corresponding
gauge-fixed Hessian on the physical quotient (or its prescribed generalized
inverse when zero modes remain) and therefore computes response; inserting it again as an
action summand would duplicate the fluctuation contribution.  A finite local
counterterm can be written in either the geometric or renormalized quantum
column, but linearity of variation shows that only the sum affects the field
equations.  Including the same term in both columns doubles its variation.
Thus this assumed functional supplies a consistent role ledger, while its
harmless counterterm placement convention is not fixed.  The argument does
not exclude a distinct probe-complete functional with the same terminal
equations.
\end{proof}

\paragraph{Scope.}
Theorem~\ref{thm:effectivesource} is a construction theorem inside the
selected ordered completion.  It does not show that the physical base--seed
law selects $\Gamma_{\rm PC}$, $\mathcal C_O$, its field spectrum or its
renormalized continuum measure.

The ordering is structural. If $H_B,H_P>0$ and the post-body equations have
body derivative $C$, the linear response is lower triangular:
\begin{equation}
\begin{pmatrix}\delta y\\\delta z\end{pmatrix}
=\begin{pmatrix}H_B^{-1}&0\\-H_P^{-1}CH_B^{-1}&H_P^{-1}\end{pmatrix}
\begin{pmatrix}\delta j_B\\\delta j_P\end{pmatrix}.
\label{eq:triangular}
\end{equation}
Thus morphology conditions the physical terminals, while a later terminal
load does not rebuild the body.

\begin{proposition}[Simultaneous-action no-go]
For an invertible positive simultaneous Hessian
\begin{equation}
\mathbb H=\begin{pmatrix}H_B&C^\dagger\\C&H_P\end{pmatrix},
\end{equation}
exact vanishing of body response to every pure post-body load implies $C=0$.
Hence nonzero body conditioning and exact body freeze cannot both be encoded
by one simultaneous unconstrained Hessian solve.
\end{proposition}
\begin{proof}
Schur inversion gives
$\delta y=-H_B^{-1}C^\dagger
(H_P-CH_B^{-1}C^\dagger)^{-1}\delta j_P$.
If this vanishes for every $\delta j_P$, invertibility forces $C=0$.
\end{proof}

\section{Disjoint occupation and stress ownership}

Let the occupied mode labels form a disjoint union
\begin{equation}
\mathcal I_{\rm occ}=\mathcal I_{\rm cl}\sqcup\mathcal I_q.
\label{eq:partition}
\end{equation}
Classical coherent modes represented in $\Gamma_{\rm occ}$ are excluded from
the state-dependent part of $\Gamma_q$; vacuum polarization and local
renormalization terms occur only in $\Gamma_q+\Gamma_{\rm ct}$.
This disjoint partition is a sufficient transparent ledger, not the only
possible field-theoretic organization. A background-field, one-particle-
irreducible or two-particle-irreducible construction is also admissible when
its mean field, fluctuation determinant and counterterms are derived from one
effective action and no full contribution is inserted twice.

\begin{proposition}[Duplicate-stress no-go]
If one nonzero occupied mode $i$ is included in both $\mathcal I_{\rm cl}$ and
$\mathcal I_q$ with the same metric variation, its stress contribution is
counted twice. No terminal redefinition removes this without changing the
metric equation or adding a compensating subtraction.
\end{proposition}
\begin{proof}
Metric variation is linear. The duplicated summands give
$-2(\sqrt{-g})^{-1}\delta(\Gamma_i+\Gamma_i)/\delta g^{\mu\nu}
=2T^{(i)}_{\mu\nu}$. A subtraction merely moves the ownership entry and must
therefore be recorded explicitly in $\Gamma_{\rm ct}$.
\end{proof}

\section{Joint terminal theorem}

\begin{proposition}[Conditional standard-leaf form selector]
\label{prop:standardforms}
Assume that the selected readout leaf is a smooth four-dimensional Lorentzian
manifold.  Suppose the pure geometric action is local, diffeomorphism
invariant, metric-only, and has metric Euler equations of at most second
derivative order.  Suppose also that the Abelian action is local,
$U(1)$-gauge invariant, parity even, quadratic in first derivatives of $A$,
contains no derivatives of $F$, and has a constant positive kinetic
coefficient.  Then, up to boundary and four-dimensional topological terms,
the two actions have the forms
\begin{align}
\Gamma_{\rm grav}&={1\over2\kappa}\int d^4x\sqrt{-g}(R-2\Lambda),\\
\Gamma_{\rm EM}&=-{Z_F\over4}\int d^4x\sqrt{-g}\,F_{\mu\nu}F^{\mu\nu},
\qquad Z_F>0.
\end{align}
The proposition fixes the action shapes inside this infrared class, but not
$\kappa$, $\Lambda$, $Z_F$, or physical realization of the class itself.
\end{proposition}
\begin{proof}
The four-dimensional Lovelock classification leaves only the cosmological and
Einstein tensors in the local metric-only second-order Euler tensor; the
Gauss--Bonnet density is topological in four dimensions
\citep{Lovelock1971}.  Integration gives the Einstein--Hilbert form up to
boundary/topological terms and overall normalization.  Abelian gauge
invariance restricts a local first-derivative action to dependence on
$F=dA$.  In four dimensions the quadratic Lorentz scalars are
$F_{\mu\nu}F^{\mu\nu}$ and
$F_{\mu\nu}\widetilde F^{\mu\nu}$.  Parity removes the second from the
dynamical class (and for constant coefficient it is topological), so the
first is unique up to $Z_F$; positivity fixes its sign.  These classification
arguments neither fix the coefficients nor prove that the parent selects
their hypotheses.
\end{proof}

The selector does not rule out higher-curvature gravity, torsion, additional
geometric fields, nonlinear electrodynamics or parity-odd sectors; those lie
outside its hypotheses.  Its content is uniqueness within the declared
standard infrared class, not uniqueness among all physical completions.

For the joint theorem adopt Proposition~\ref{prop:standardforms}, the
conditional common-source handoff $Z_F=p_*K_Y$, and a matter functional,
quantum measure and renormalization prescription that are diffeomorphism and
$U(1)$ gauge invariant and anomaly-free in those symmetries.  The handoff
removes an independent Maxwell gain inside the declared completion; it does
not evaluate $p_*$ or select the occupied representation trace $K_Y$ from the
physical parent.

\begin{theorem}[Common-descent Einstein--Maxwell--matter assembly]
\label{thm:assembly}
Under the ordered ledger, disjoint occupation rule, common unit packet and the
stated invariances, stationarity of Eq.~\eqref{eq:jointfunctional} with
$\rho_{\rm in,O}$ held fixed gives
\begin{align}
G_{\mu\nu}+\Lambda g_{\mu\nu}
&=\kappa\left(T^{\rm EM}_{\mu\nu}+T^{\rm cl}_{\mu\nu}
+\langle T^q_{\mu\nu}\rangle_{\rm ren}\right),
\label{eq:einstein}\\
\nabla_\mu(Z_FF^{\mu\nu})
&=J_{\rm cl}^{\nu}+\langle J_q^\nu\rangle_{\rm ren}.
\label{eq:maxwell}
\end{align}
On the matter and gauge equations, the Ward identities and Bianchi identity
imply
\begin{equation}
\nabla^\mu T^{\rm tot}_{\mu\nu}=0,
\qquad
\nabla_\mu J_{\rm tot}^{\mu}=0.
\label{eq:ward}
\end{equation}
No additional gravity, gauge or quantum-backreaction gain is introduced at the
terminal stage.
\end{theorem}
\begin{proof}
Equations \eqref{eq:einstein} and \eqref{eq:maxwell} are the metric and gauge
Euler equations of the single post-body functional. Gauge invariance gives
current conservation. Diffeomorphism invariance gives the combined stress
Ward identity, including the Lorentz-force exchange between matter and gauge
sectors. The divergence of the left side of Eq.~\eqref{eq:einstein} vanishes
identically, so the equations are mutually consistent. Disjoint occupation
and the single-entry counterterm ledger ensure that each contribution is varied
once.
\end{proof}

\begin{corollary}[Standard weak-field leaf]
For $g_{00}=-(1+2\Phi_N/c^2)$, stationary weak fields, nonrelativistic
sources and a local patch whose boundary data admit the Newtonian limit,
Eq.~\eqref{eq:einstein} reduces to
\begin{equation}
\nabla^2\Phi_N=4\pi G\rho_{\rm tot}-\Lambda c^2,
\end{equation}
where $\rho_{\rm tot}$ is the same once-counted density appearing in the
ledger.  On scales $L\ll |\Lambda|^{-1/2}$, or for $\Lambda=0$, the final term
is negligible and the ordinary Poisson equation follows. Equation
\eqref{eq:maxwell} reduces to the ordinary Maxwell equations after the
inherited field normalization. This is standard recovery, not a Tau-specific
deviation.
\end{corollary}

\section{Relation to established frameworks}

The use of one renormalized effective action and its metric variation belongs
to semiclassical gravity and curved-spacetime QFT, not uniquely to Tau Core
\citep{Wald1977,BirrellDavies1982,HollandsWald2001}.  The closed-time-path
extension and stress-noise cumulants likewise parallel stochastic gravity
\citep{HuVerdaguer2008}.  Observer-relative descriptions have independent
precedent in quantum-reference-frame constructions
\citep{GiacominiCastroRuizBrukner2019}, while deriving Einstein dynamics from
more primitive thermodynamic data is exemplified by emergent-gravity programs
\citep{Jacobson1995}.

The finite-dimensional algebraic ingredients are also established
mathematics rather than Tau-specific discoveries.  The exterior realization
of Clifford multiplication is standard spin geometry
\citep{LawsonMichelsohn1989}; cyclic representations and their unitary
uniqueness belong to the GNS/operator-algebra framework
\citep{BratteliRobinson1979}; and special dagger--Frobenius structures provide
the finite-dimensional algebraic language for normalized copy/record data
\citep{Vicary2011,CoeckePavlovicVicary2013}.  Paper VII's claim is therefore
not novelty of these ingredients separately, but their conditional placement
in one source-owned, once-counted terminal ledger and the accompanying
non-identifiability controls.

Paper VII's narrower Tau-specific contribution is the body-first ordering,
the typed morphology-conditioned access and resolution maps, the explicit
once-counted cross-terminal ownership ledger, and the finite no-go separating
terminal agreement from common provenance.  The standard effective-action
machinery is used as a consistency leaf; it is not presented as evidence for
the Tau parent ontology.

\section{Local constraint structure}

The covariant equations must reproduce constraints, not only evolution
equations. Let $n^\mu$ be the unit normal to an observer-compatible spacelike
leaf, $h_{\mu\nu}=g_{\mu\nu}+n_\mu n_\nu$ its induced metric, and $K_{ij}$ its
extrinsic curvature. With
$\rho=n^\mu n^\nu T^{\rm tot}_{\mu\nu}$ and
$j_i=-h_i{}^\mu n^\nu T^{\rm tot}_{\mu\nu}$, the normal projections of
Eq.~\eqref{eq:einstein} give, in the stated sign convention,
\begin{align}
{}^{(3)}R+K^2-K_{ij}K^{ij}&=2\kappa\rho+2\Lambda,
\label{eq:hamconstraint}\\
D_j(K^{ij}-h^{ij}K)&=\kappa j^i.
\label{eq:momconstraint}
\end{align}
The normal projection of Eq.~\eqref{eq:maxwell} gives the Abelian Gauss
constraint
\begin{equation}
D_i(Z_FE^i)=\rho_J.
\label{eq:gaussconstraint}
\end{equation}
The Ward identities propagate these constraints when the classical and
quantum matter equations hold. Paper VII therefore imports the standard local
constraint architecture from the one joint functional; it does not derive a
new constraint algebra or quantize it.

\section{Primitive terminal inventory}

The assembled packet contains several typed outputs, not several names for
the same information:
\begin{center}
\small
\begin{tabularx}{0.96\textwidth}{@{}p{0.20\textwidth}Xp{0.25\textwidth}@{}}
\toprule
Family & Primitive terminal content & Derived or composite descendants\\
\midrule
Geometry/gravity & $g$, connection, curvature and constraints & proper time,
tidal response, lensing after a null source\\
Gauge/radiation & $A$, $F$, conserved current, two screen polarizations &
photon propagation, redshift and optical transport\\
Matter occupation & occupied localized or propagating solutions $\Phi$ &
mass, energy, stress and charge records\\
Quantum operation & $\rho$, effects, instruments and correctable composition &
Born statistics and renormalized quantum stress/current\\
Observer access & causal orientation, clock and stable restriction &
observer-relative terminal classes\\
\bottomrule
\end{tabularx}
\end{center}
Stress-energy, current, redshift and lensing are variational or composite
outputs of these families; they are not silently counted as new primitive
source directions. Conversely, the symmetric local gravity descriptor cannot
stand in for internal gauge charge, quantum phase or occupied matter. A global
information-completeness claim must wait for the non-Abelian and fermionic
known-terminal completions.

\section{Why the standard limit does not prove common provenance}

\begin{proposition}[Equal-terminal provenance no-go]
The equations \eqref{eq:einstein}--\eqref{eq:ward} do not imply that their
coefficients or sectors arise from one Tau source.
\end{proposition}
\begin{proof}
Construct two source models with the identical numerical effective functional
$\Gamma_7[g,A,\Phi]$.  In $C_{\rm joint}$ its sector coefficients are images
of one common source line and unit map.  In $C_{\rm sep}$ take four independent
source lines and choose their terminal coefficient maps to have the same
numerical images as the four coefficients in $C_{\rm joint}$.  The two models
therefore have identical functionals as functions of the terminal fields, so
all Euler equations, Ward identities and solution sets coincide.  Their
source ownership graphs are nevertheless non-isomorphic: one has a common
source line and the other has four independent lines.  Hence no observation
or algebraic manipulation of Eqs.~\eqref{eq:einstein}--\eqref{eq:ward} alone
can recover common provenance.  This is an identifiability counterexample,
not a claim that the independently sourced model is physically preferred.
\end{proof}

This no-go fixes the empirical interpretation: terminal data may test a
source-frozen forward ledger, but may not choose its source partition,
susceptibilities or unit map.

\begin{proposition}[Same-reduct selection no-go]
\label{prop:samereduct}
The currently frozen stationary body, common action line, observer quotient
and terminal type signatures do not entail physical occupation of the full
effective-source packet of Theorem~\ref{thm:effectivesource}.
\end{proposition}
\begin{proof}
Hold fixed the body $y_*$, its Hessian, common units, causal orientation and
observer quotient.  Add the same typed internal edge carrier
$\mathcal H_e=\mathbb C^2$ to both completions, with effect
\begin{equation}
E_e=\begin{pmatrix}0&0\\0&1\end{pmatrix},\qquad
\rho_w=\begin{pmatrix}1-w&0\\0&w\end{pmatrix},
\quad 0<w<1,\qquad c_0=0,\quad c_1=1.
\label{eq:counterpair}
\end{equation}
$\rho_w$ is positive, has unit trace and is invariant under the diagonal
$U(1)$ action.  A concrete covariant edge functional is supplied by
a charged scalar doublet $\chi\in\mathbb C^2$ with
\begin{equation}
\Gamma_e[\chi,g,A]
=-\int d^4x\sqrt{-g}\left[
g^{\mu\nu}(D_\mu\chi)^\dagger D_\nu\chi
+m_e^2\chi^\dagger\chi\right],
\qquad D_\mu=\partial_\mu-iqA_\mu I_2 .
\label{eq:counteredgeaction}
\end{equation}
It is diffeomorphism invariant, diagonal-$U(1)$ gauge invariant and has the
same positive internal state normalization in both completions.  Define
\begin{equation}
\Gamma_{\rm PC}^{(j)}[\eta_e]
=\Gamma_{\rm common}+\Gamma_e+c_j\eta_e\mathsf E_e,
\qquad j\in\{0,1\}.
\label{eq:counterprobes}
\end{equation}
At the physical source value $\eta_e=0$, both completions have exactly the
same action, state, body, Hessian, observer quotient, typed carrier,
covariance and normalization.  The frozen upstream reduct does not include a
law selecting the probe coefficient $c_j$.  Nevertheless,
\begin{equation}
{\delta W_O^{(0)}\over\delta\eta_{e,\Delta}}\bigg|_{\rm phys}=0,
\qquad
{\delta W_O^{(1)}\over\delta\eta_{e,\Delta}}\bigg|_{\rm phys}
=\operatorname{Tr}(\rho_wE_e)=w>0.
\end{equation}
Only $C_1$ contains the declared nonzero effect-source owner.  Since the full
effective-source packet includes that owner, the frozen upstream reduct does
not entail the complete packet.  The result concerns source ownership, not
the logically different statement that a realized POVM outcome must have
nonzero probability in every state.  No stress- or current-difference is
claimed because the physical unprobed action and field state are identical.
Additional source-owned occupation or an empirically discriminating
source-frozen prediction is required.
\end{proof}

This finite counterpair is the stopping result for the origin question at the
present source level.  It forbids both an unrestricted parent-law promotion of
Theorem~\ref{thm:effectivesource} and another same-level search for a hidden
normalization.  It does not weaken the conditional forward construction.

Figure~\ref{fig:jointledger} summarizes the resulting generative and
variational order.

\begin{figure}[t]
\centering
\includegraphics[width=0.92\textwidth]{figures/fig_joint_ledger.pdf}
\caption{Ordered common-source construction. The body is stabilized before the
post-body functional is solved; observer restriction is terminal.}
\label{fig:jointledger}
\end{figure}

\section{Quantum backreaction and the quantum-gravity boundary}

The quantum term in Eq.~\eqref{eq:einstein} is a renormalized expectation
value on a classical body-derived metric, in the standard semiclassical sense
\citep{BirrellDavies1982,Wald1977}. Local covariance and admissible
renormalization follow the usual curved-spacetime QFT framework
\citep{HollandsWald2001,BrunettiFredenhagenVerch2003}. The metric can also be
driven by a stress-noise second cumulant in a stochastic extension, but this
does not quantize geometry.

Paper VII therefore establishes a conditional quantum--gravity
\emph{interface}, not quantum gravity. A path integral or operator algebra for
the metric, gravitational gauge fixing, a physical inner product and a
nonperturbative ultraviolet completion remain absent.

\subsection{Two-mode common-preparation theorem}

The preceding interface can be made more specific without quantizing the
metric.  Fix an oriented observer Sachs screen and an occupied two-mode QFT
restriction with annihilation operators $a_m,a_{m+2}$.  Let
\begin{equation}
C_{ij}:=\langle a_i^\dagger a_j\rangle,
\qquad
N_O:=\operatorname{Tr}C>0,
\qquad
\rho_{\rm ROOT}:=\frac{C}{N_O}.
\label{eq:rootdensity}
\end{equation}
The two screen weights differ by two, so their coherence transforms in the
same spin-two representation as the complex screen stress
$T_+ + iT_\times$.  Denote the once-counted normal-ordered STF overlap of the
two occupied modes by $K_{m,m+2}$.  It is not a free terminal gain.  Let
$u_i$ be the source-normalized mode functions, let
$e_+^\mu=(e_1^\mu+i e_2^\mu)/\sqrt{2}$ be the oriented complex screen vector,
and let $f_O$ be a normalized observer smearing window.  If
$\mathcal T_{\mu\nu}[u_i^*,u_j]$ denotes the polarized bilinear appearing in
the once-renormalized Hilbert stress derived from the same QFT action, define
\begin{equation}
K_{ij}[u,f_O]
:=\int d\mu_O\,f_O\,e_+^\mu e_+^\nu
\mathcal T_{\mu\nu}[u_i^*,u_j].
\label{eq:overlapformula}
\end{equation}
Thus the action, mode normalization, screen, smearing and subtraction scheme
determine $K_{ij}$.  The common-preparation identity is
\begin{equation}
\boxed{
\left\langle :T_+ + iT_\times: \right\rangle
=K_{m,m+2}C_{m,m+2}
=N_O K_{m,m+2}(\rho_{\rm ROOT})_{m,m+2}.}
\label{eq:rootttidentity}
\end{equation}

\begin{proposition}[Conditional ROOT--TT common preparation]
\label{prop:roottt}
Assume (i) the physical restriction occupies a common, oriented screen pair
whose weights differ by two; (ii) the relevant state is restricted to a
number-preserving two-mode second-moment sector, so anomalous moments
$\langle a_i a_j\rangle$ and $\langle a_i^\dagger a_j^\dagger\rangle$ vanish
or are independently projected out; (iii) all other mode, vacuum and
polarization contributions to the selected STF component are orthogonal or
included once in the declared subtraction; (iv) the ROOT one-particle reduced
state and the stress expectation are formed from the same source-owned QFT
preparation; (v) normal
ordering or renormalization is entered once in the source ledger; and (vi)
$K_{m,m+2}\neq0$, with $K_{m,m+2}$ calculated by
Eq.~\eqref{eq:overlapformula} from the declared source preparation, screen and
subtraction prescription.  Then Eq.~\eqref{eq:rootttidentity} follows and,
after $K_{m,m+2}$ is source-frozen, admits no additional independent
terminal-specific gain within this restricted sector.
\end{proposition}

\begin{proof}
The ROOT one-particle reduced state is the normalized one-particle coherence matrix by
Eq.~\eqref{eq:rootdensity}.  Under assumptions (ii)--(iii), projection of the
normal-ordered quadratic stress onto the selected spin-two component leaves
$K_{m,m+2}\langle a_m^\dagger a_{m+2}\rangle$.  Substitution of
$C=N_O\rho_{\rm ROOT}$ gives Eq.~\eqref{eq:rootttidentity}.  Introducing an
additional gain after $K_{m,m+2}$ is fixed would assign two calibrations to one
occupied preparation and would violate the once-counted source ledger.  The
formula does not select the mode pair, screen or subtraction prescription and
does not prove that the resulting overlap is nonzero.  Changing those
source-side inputs changes the calculated $K_{m,m+2}$; it does not authorize a
post-hoc terminal gain.
\end{proof}

The use of the unnormalized matrix $C$, rather than only
$\rho_{\rm ROOT}$, also closes the quantum-induced part of the P1 gluing.
Let $J_C:\delta s\mapsto\delta C$, let $\pi_{12}$ extract the real and
imaginary parts of $C_{m,m+2}$, let $M_K$ denote real multiplication by the
source-frozen complex number $K_{m,m+2}$, and let $L_{P1}^{\rm TT}$ be the
retarded linear response from the screen-TT stress plane to the P1 geometry
rows.  Then
\begin{equation}
\boxed{
J_{P1}^{(Q)}=L_{P1}^{\rm TT}M_K\pi_{12}J_C,}
\qquad
\ker J_C\subseteq\ker J_{P1}^{(Q)}.
\label{eq:rootp1factor}
\end{equation}
If $K_{m,m+2}\neq0$ and $L_{P1}^{\rm TT}$ is injective on the occupied TT
plane, this restriction has real rank two.  This is a factorization theorem
for the quantum-induced stress row, not for the complete P1 terminal.  The
distinction is essential: $J_Q(x,y,z)=(x,y)$ and
$J_{P1}(x,y,z)=(x,y,z)$ have a rank-two common submap although
$(0,0,1)\in\ker J_Q\setminus\ker J_{P1}$.  Thus local ROOT--TT rank two does
not prove the full-carrier zero-cokernel condition.  Conversely, replacing
$C$ by normalized $\rho_{\rm ROOT}$ alone loses this factorization under
occupation-scale variations: $\delta C=\epsilon C$ gives
$\delta\rho_{\rm ROOT}=0$ but a nonzero TT-stress variation whenever the
coherence and overlap are nonzero.

On the physical pure-state ROOT tangent the existing body-first staging makes
this statement stronger.  With observer and radiative boundary data frozen,
the direct Green row, stabilized body stress and direct
ROOT-to-occupied-matter term have zero variation.  Therefore
\begin{equation}
\boxed{
J_{P1}|_{\mathcal V_Q}
=J_{P1}^{(Q)}|_{\mathcal V_Q}
=L_{P1}^{\rm TT}M_K\pi_{12}J_C|_{\mathcal V_Q}.}
\label{eq:rootp1tangentfactor}
\end{equation}
Thus Eq.~\eqref{eq:rootp1factor} is the full local P1 factorization on the
two-real-dimensional post-body ROOT tangent, while the counterexample above
continues to exclude a global conclusion on the complete source domain.

The local result can nevertheless be removed exactly from the remaining
global rank calculation.  Choose the source-metric orthogonal splitting
$\mathcal E_{\rm phys}=\mathcal E_R\oplus\mathcal V_Q$ and write
$J_Q=(C_R,D)$ and $J_{P1}=(A_R,\iota D)$, where $D$ is injective on the
two-dimensional ROOT tangent and $\iota=L_{P1}^{\rm TT}M_K\pi_{12}$.
If $P_D$ projects onto $\operatorname{im}D$, define

\begin{equation}
\mathcal R_{\rm rest}x
=\bigl((I-P_D)C_Rx,(A_R-\iota C_R)x\bigr),
\qquad x\in\mathcal E_R.
\label{eq:stagedglue}
\end{equation}

Then projection to $\mathcal E_R$ gives

\begin{equation}
\boxed{
\ker\begin{pmatrix}J_Q\\J_{P1}\end{pmatrix}
\simeq\ker\mathcal R_{\rm rest},
\qquad
\operatorname{rank}\begin{pmatrix}J_Q\\J_{P1}\end{pmatrix}
=2+\operatorname{rank}\mathcal R_{\rm rest}.}
\label{eq:stagedgluerank}
\end{equation}

Indeed, $C_Rx+Dy=0$ is soluble exactly when
$(I-P_D)C_Rx=0$ and then has the unique solution
$y=-D^{-1}C_Rx$; the P1 equation becomes
$(A_R-\iota C_R)x=0$.  Thus the already closed ROOT tangent contributes
exactly two ranks.  The unresolved differential content is confined to the
rest-domain transverse quantum component and any independent P1 component.
This is an exact staged reduction, not physical realization of the complete
base--seed packet or a replacement for finite-support and algebraic
faithfulness tests.

Nor do the current marginal data determine
$\operatorname{rank}\mathcal R_{\rm rest}$. Take
$\mathcal E_R=\mathbb R^4$, $D=\iota=I_2$,
$C_R=(I_2\;0)$ and respectively

\begin{equation}
A_0=(I_2\;0),\qquad
A_1=\begin{pmatrix}1&0&0&0\\0&-1&0&0\end{pmatrix},\qquad
A_2=(0\;I_2).
\label{eq:restcontrols}
\end{equation}

All three controls have the same faithful local ROOT block, raw ranks and
marginal Grams $J_QJ_Q^\dagger=J_{P1,k}J_{P1,k}^\dagger=2I_2$, but
$\operatorname{rank}\mathcal R_{{\rm rest},k}=k$. Thus the remaining rank
requires the rest-domain mixed P1--quantum incidence. It is calculated by
the common-source mixed-Gram formula once the full packet is physically
selected; separate marginal closure does not select it.

The inherited common rank-four body theorem removes a larger local block.
Let the local source split be $X_B\oplus H_{\rm hid}$ and write the retained
rank-four observer descriptor and P1 differential as
\begin{equation}
J_4=(A\;B),\qquad J_{P1}=(P\;Q),\qquad A:X_B\overset{\sim}{\longrightarrow}Y_4.
\end{equation}
Full local qubit and Lorentz retention makes $A$ invertible.  Defining
\begin{equation}
\Delta_{\rm hid}=Q-PA^{-1}B
\label{eq:hiddenfiberdefect}
\end{equation}
and applying exact block elimination gives
\begin{equation}
\boxed{
\ker\binom{J_4}{J_{P1}}\simeq\ker\Delta_{\rm hid},\qquad
\operatorname{rank}\binom{J_4}{J_{P1}}
=4+\operatorname{rank}\Delta_{\rm hid}.}
\label{eq:hiddenfiberrank}
\end{equation}
Hence the P1 map factors uniquely through the observer/ROOT descriptor on
$X_B$; no selector is missing on that common core.  If the descent is
body-block preserving, then $B=0$ and the remaining defect is simply $Q$.
This result does not determine the hidden-fibre block, identify the local
rank-four packet with the complete 64-real-dimensional P3 descriptor, or
prove physical occupation of either completion.

The complete P3 comparison can be reduced on the same source split.  Write
$J_{P3}=(R\;S)$ and remove its already counted common-core component:
\begin{equation}
\Gamma_{\rm hid}:=S-RA^{-1}B.
\label{eq:p3hiddenreduction}
\end{equation}
The same invertible block operations give
\begin{equation}
\operatorname{rank}
\begin{pmatrix}J_4\\J_{P3}\\J_{P1}\end{pmatrix}
=4+\operatorname{rank}
\begin{pmatrix}\Gamma_{\rm hid}\\\Delta_{\rm hid}\end{pmatrix}.
\label{eq:p3hiddenrank}
\end{equation}
Consequently the hidden P1 defect is represented by P3 exactly when
\begin{equation}
\boxed{\ker\Gamma_{\rm hid}\subseteq\ker\Delta_{\rm hid}}
\quad\Longleftrightarrow\quad
\Delta_{\rm hid}=L_{P1\leftarrow P3}\Gamma_{\rm hid}
\quad\hbox{on }\operatorname{im}\Gamma_{\rm hid}.
\label{eq:p3hiddenfactor}
\end{equation}
This is a source-fibre condition, not a target-rank condition.  Indeed,
$\Gamma_{\rm hid}(x,y)=x$ is onto its one-dimensional P3 target while
$\Delta_{\rm hid}(x,y)=y$ detects a P3-null source direction.  Thus complete
P3 operator-state separation does not by itself close the P1--P3 gluing.
Source-fibre injectivity of $\Gamma_{\rm hid}$ would be sufficient, while the
weaker kernel inclusion in Eq.~\eqref{eq:p3hiddenfactor} is necessary and
sufficient.  Neither condition is presently selected by the physical
base--seed law.

The remaining branch choice is represented by one canonical source-fibre
obstruction,
\begin{equation}
\mathcal O_{P1|P3}
:=\Delta_{\rm hid}\big|_{\ker\Gamma_{\rm hid}} .
\label{eq:p3obstruction}
\end{equation}
Rank--nullity and Eq.~\eqref{eq:p3hiddenrank} give the exact novelty identity
\begin{equation}
\boxed{
\operatorname{rank}
\begin{pmatrix}J_4\\J_{P3}\\J_{P1}\end{pmatrix}
-
\operatorname{rank}
\begin{pmatrix}J_4\\J_{P3}\end{pmatrix}
=\operatorname{rank}\mathcal O_{P1|P3}.}
\label{eq:p3obstructionrank}
\end{equation}
Thus $\mathcal O_{P1|P3}=0$ is equivalent to complete P1 factorization through
P3 after the common core has been counted.  If the obstruction is nonzero,
its source direction is P1-visible, so it is neither a joint readout null nor
an unloaded Schur-neutral presentation auxiliary; source-minimal quotienting
cannot remove it.  The current parent laws admit both branches and do not
select the value of this obstruction.

The common source Hessian also gives a direct evaluation formula.  Let
$K_{\rm hid}>0$ be its Schur complement after the common rank-four block is
eliminated and define
\begin{equation}
\widehat\Gamma=\Gamma_{\rm hid}K_{\rm hid}^{-1/2},\qquad
\widehat\Delta=\Delta_{\rm hid}K_{\rm hid}^{-1/2},\qquad
\Pi_{\ker\widehat\Gamma}=I-\widehat\Gamma^+\widehat\Gamma .
\end{equation}
Then
\begin{equation}
\boxed{\operatorname{rank}\mathcal O_{P1|P3}
=\operatorname{rank}
\left(\widehat\Delta\Pi_{\ker\widehat\Gamma}\right).}
\label{eq:p3obstructionprojector}
\end{equation}
Equivalently, form the marginal and mixed Grams of $\widehat\Gamma$ and
$\widehat\Delta$.  If $\mathcal C_{\Gamma\Delta}$ is their support-normalized
mixed Gram, the principal-angle theorem gives
\begin{equation}
\operatorname{rank}\mathcal O_{P1|P3}
=\operatorname{rank}\Delta_{\rm hid}
-\#\{j:\sigma_j(\mathcal C_{\Gamma\Delta})=1\}.
\label{eq:p3obstructiongram}
\end{equation}
Thus one physically selected common incidence packet computes the obstruction
without an additional terminal coefficient.  The Hessian alone does not fix
it: distinct typed incidences over the same positive metric can realize zero
or nonzero obstruction.

This obstruction also separates two notions of completeness that must not be
conflated.  Let $N_{\rm phys}\subset H_{\rm hid}$ be the source null
annihilated by every retained terminal, so in particular
$N_{\rm phys}\subseteq\ker\Delta_{\rm hid}$.  If P3 is complete on physical
source fibres modulo this joint null, then
\begin{equation}
\ker\Gamma_{\rm hid}\subseteq N_{\rm phys}
\subseteq\ker\Delta_{\rm hid},
\end{equation}
and therefore
\begin{equation}
\boxed{
\text{P3 physical-source completeness}
\Longrightarrow \mathcal O_{P1|P3}=0
\Longrightarrow J_{P1}=L_{P1\leftarrow P3}J_{P3}}
\label{eq:p3physicalcompleteness}
\end{equation}
after the common core and joint null have been removed.  Conversely, a
nonzero obstruction provides a P3-null source direction that P1 still sees;
it therefore certifies failure of P3 source-fibre completeness even when P3
is onto and informationally complete on its quantum operator-state target.
Equation~\eqref{eq:p3physicalcompleteness} is a conditional completeness
criterion, not a proof that the physical base--seed parent selects the
source-complete branch.

The existing seed-chain/Hodge construction fixes a genuine subpacket but does
not close this selection problem.  It supplies a nonzero nilpotent incidence
on the occupied chain image.  On a source split
$E_{\rm phys}=E_{\rm chain}\oplus H_{\rm ext}$, however, fix the same positive
metric and take
\begin{equation}
\Gamma_{\rm hid}=(1,0,0),\qquad
\Delta_0=(2,0,0),\qquad
\Delta_1=(2,1,0).
\label{eq:seedchainextensions}
\end{equation}
The two completions have the same nonzero chain restriction and the same onto
P3 map, while
$\Delta_0|_{\ker\Gamma_{\rm hid}}=0$ and
$\operatorname{rank}(\Delta_1|_{\ker\Gamma_{\rm hid}})=1$.  Thus the proved
chain incidence must be retained and may not be fitted again, but it does not
select the complete terminal-typed extension $J_*$.  That remaining extension
is the physical representation of the primitive relation presentation on
$H_{\rm ext}$, not another Hodge coefficient or terminal gain.

One selected chain block would extend uniquely if a source-owned
typed-isometric transport groupoid made its orbit span $E_{P3}$ and respected
all stabilizer relations.  This sufficient route is unavailable in the
current typed ledger: the two-role observer--instrument defect and the
three-role photon--matter defect cannot belong to one type-preserving source
orbit.  Nor does cyclicity of the represented P3--record GNS carrier replace
this condition.  For example, the defining $M_2(\mathbb C)$ representation
is cyclic from $e_1$ because $M_2(\mathbb C)e_1=\mathbb C^2$, whereas a typed
diagonal source group preserves the line $\mathbb C e_1$.  Thus target-algebra
cyclicity does not imply typed source-orbit generation.  At least one
source-owned inter-orbit relation or separately occupied orbit representative
is still required for the full extension.

This remaining relation has an exact finite closure criterion.  Decompose the
typed source carrier into its source orbits,
$V_r=\bigoplus_\alpha V_\alpha$, let $\mathcal A_0$ be the star-algebra
generated by the fixed intra-orbit representations, and let $B$ contain the
proposed source-owned inter-orbit bridge blocks.  Define
\begin{equation}
\mathcal A_{\rm br}=\operatorname{alg}^*(\mathcal A_0,B).
\label{eq:bridgealgebra}
\end{equation}
For fixed orbit metrics, two relative gluings preserve all generated source
relations exactly when their relative unitary lies in
$U(\mathcal A_{\rm br}')$.  Assuming
$G_{\rm decl}\subseteq U(\mathcal A_{\rm br}')$, the gluing is unique modulo
the declared gauge group precisely when
\begin{equation}
\boxed{U(\mathcal A_{\rm br}')=G_{\rm decl}.}
\label{eq:bridgecommutant}
\end{equation}
This is an internal sufficiency-and-uniqueness theorem, not an existence
theorem for $B$.  In the two-scalar-orbit control, the two orbit projectors
alone have a two-dimensional commutant and leave an independent relative
phase.  Adding one nonzero off-diagonal bridge and its dagger generates
$M_2(\mathbb C)$ and reduces the commutant to the common scalar line.  This
bridge is required only when the declared gauge is the common scalar.  Under
the inclusion assumed above, the actual obstruction is the excess commutant
\begin{equation}
d_{\rm excess}=\dim U(\mathcal A_{\rm br}')-\dim G_{\rm decl}.
\label{eq:bridgeexcess}
\end{equation}
If the two diagonal phases are instead the declared two-sector
gauge-superselection center, the unbridged model already has
$d_{\rm excess}=0$; an off-diagonal bridge would then erase a declared
superselection direction and is inadmissible without an upstream refinement
of the gauge quotient.

This qualification does not close the P3 source problem.  The existing
superselection theorem concerns the quark/lepton multiplet split
$\{Q,u^c,d^c\}$ versus $\{L,e^c\}$, whereas the present no-go separates
relation defects with typed supports $\{Q,O\}$ and $\{R,M,G\}$.  These are not
the same decomposition.  Moreover, $Q$ denotes the ROOT/quantum role in the
P3 packet but the quark-doublet module in the multiplet theorem; equality of
the glyph is not a typed intertwiner.  No current theorem transfers the former
gauge center to the latter defect carrier.  The current base--seed packet therefore
fixes neither the relevant $G_{\rm decl}$ nor, if a positive excess remains,
the required bridge.  The adopted P3 handoff is post-body data and cannot be
used backwards to generate this upstream source relation.

More precisely, a target gauge action $U_g$ on a carrier $H$ induces a gauge
action on the relation-defect carrier $V_r$ only after a source-owned
injective intertwiner $\iota_r:V_r\to H$ is supplied and
$U_g\operatorname{im}\iota_r\subseteq\operatorname{im}\iota_r$.  Then and
only then the pullback is unique:
\begin{equation}
\iota_r\gamma_r(g)=U_g\iota_r,
\qquad
G_{\rm decl}=\gamma_r(G_{\rm phys})/\ker\gamma_r.
\label{eq:gaugepullback}
\end{equation}
The odd-Clifford central projectors act on $\Lambda^\bullet E_{P3}$, the
two-point record projectors on $R_2$, and the quark/lepton center on multiplet
modules.  None acts on $V_r$ through a currently derived $\iota_r$.  Their
common two-sector count and equal trace weights therefore do not determine
$G_{\rm decl}$ on the typed defect orbits.

There is nevertheless a source-internal positive route.  The primitive
faithful GNS completion already contains a typed injection
$\iota_r:V_r\to\mathcal A_{\rm src}/N_{\omega_s}$.  Let
$G_{\rm GNS}^{\rm typed}$ be its state-preserving star-automorphisms that
preserve the image and role typing.  Their GNS implementation is unitary, so
the gauge-null subgroup relevant to the gluing theorem is determined by
\begin{equation}
\boxed{
G_{\rm decl}^{\rm GNS}
=\gamma_r(G_{\rm GNS}^{\rm typed})\cap U(\mathcal A_{\rm br}').}
\label{eq:gnsgaugecandidate}
\end{equation}
Consequently no independent gauge quotient is needed after that complete GNS
source packet is selected.  What remains open is precisely its physical
base--seed selection and whether its typed stabilizer exhausts the physical
gauge group.

The canonical exterior--record construction does not prove that selection.
For its already supplied positive role metric $K_*$, the
creation--contraction operator
$C_{K_*}(u)=\epsilon(u)+\iota_{K_*(u,-)}$ obeys the Clifford
anticommutator.  Cyclicity of the normalized exterior trace therefore gives
\begin{equation}
\boxed{
\operatorname{Re}\tau_{\rm ext}
 \left(C_{K_*}(u)^*C_{K_*}(v)\right)=K_*(u,v).}
\label{eq:gnsmetricpullback}
\end{equation}
This proves exact compatibility and the absence of a free GNS metric gain,
but it does not derive $K_*$: the contraction leg already contains that
metric.  A reverse use of Eq.~\eqref{eq:gnsmetricpullback} would be circular.
The noncircular source target is instead a premetric typed presentation
\begin{equation}
g_r:V_r\longrightarrow\mathcal A_{P3},\qquad
j_0:\mathcal A_{P3}\hookrightarrow\mathcal C_{16},
\label{eq:premetricpresentation}
\end{equation}
fixed without $K_*$.  If $j_0$ is unital and injective and $j_0g_r$ is
injective on the physical defect quotient, then
\begin{equation}
G_0(u,v)=\operatorname{Re}\tau_{16}
 \left[(j_0g_r(u))^*(j_0g_r(v))\right]
\label{eq:premetricpullback}
\end{equation}
is positive definite.  The once-counted relation action can then derive its
Hessian, GNS carrier, trace weights and conditional gauge pullback from this
one packet.  The present manuscript proves the implication, not physical
base--seed occupation of Eq.~\eqref{eq:premetricpresentation}.

The frozen P3 Clifford candidate gives one noncircular realization of this
target.  Take the typed basis $(B,O,R,U,Q,M,G)$, let $g_r$ map it to the
seven self-adjoint generators of $\operatorname{Cl}_7(\mathbb C)$, and let
$j_0$ be their faithful normalized Jordan--Wigner embedding as the first
seven generators of $\operatorname{Cl}_{16}(\mathbb C)$.  Their defining
anticommutator and the normalized trace give
\begin{equation}
\operatorname{Re}\tau_{16}
 \left[(j_0g_r(u))^*(j_0g_r(v))\right]
=\langle u,v\rangle_{\mathbb R^7}.
\label{eq:premetricclifford}
\end{equation}
Hence this declared completion jointly supplies the faithful GNS packet,
identity defect metric, equal odd-Clifford central weights and the
conditional gauge pullback without importing a Hessian.  It proves
nonemptiness, not unrestricted parent-law realization.  Unequal rescalings
$z_a\mapsto\lambda_a z_a$ retain the envelope, dagger and role labels while
changing the trace metric to $\operatorname{diag}(\lambda_a^2)$.  The current
base--seed laws contain no typed full-role operator-square incidence that
selects the normalized Clifford relation over this counterfamily.

There is, however, a strictly weaker conditional selector than a transitive
symmetry among the physically different roles.  Suppose all seven typed P3
lines have a source-owned degree-one lift into one primitive exterior
calculus.  Let $\epsilon_a$ and $\iota_a$ be its creation and contraction
maps, with
\begin{equation}
\{\epsilon_a,\epsilon_b\}=0,\qquad
\{\iota_a,\iota_b\}=0,\qquad
\{\iota_a,\epsilon_b\}=\delta_{ab}I,\qquad
\epsilon_a^\dagger=\iota_a.
\label{eq:p3primitivedagger}
\end{equation}
If the self-adjoint response is role-local and has no independent
degree-one bypass, then
$C_a=\lambda_a(\epsilon_a+\iota_a)$.  If, in addition, every occupied
primitive slot is counted once in the same source action unit,
\begin{equation}
\tau(C_a^2)=1\qquad\text{for every }a,
\label{eq:p3primitiveunitcost}
\end{equation}
then $\lambda_a^2=1$ and Eq.~\eqref{eq:p3primitivedagger} gives
\begin{equation}
\boxed{\{C_a,C_b\}=2\delta_{ab}I.}
\label{eq:p3primitiveselector}
\end{equation}
The remaining signs are primitive-frame reflections and do not alter the
metric or Clifford algebra.  Thus common primitive cost removes the unequal
role-local rescalings without identifying the terminal roles.  Role-locality
is essential: two unit-cost generators related by an oblique mixture can
have a nonzero cross anticommutator.

This reduces but does not close unrestricted parent-law realization and
occupation.  The existing PUSC3 law
owns exterior closure and common unit cost on its declared finite local
primitive source algebra.  The current base--seed derivation does not yet
construct a complete typed lift of $(B,O,R,U,Q,M,G)$ into that one primitive
slot class or exclude a separate degree-one bypass on its image.  The open
physical statement is therefore this typed primitive lift, rather than an
arbitrary metric coefficient or a new target carrier.

Nor can the closest existing off-diagonal constructions be substituted by
matrix shape alone.  The rooted $R_{\rm root}$/TC--CVP1 Hessian acts between
pair-difference and pair-common $W_2$ occurrence modules, not between the P3
defect orbits.  The wall--observer/internal separator candidate leaves
distinct multiplet blocks disconnected, while the existing EACER,
seed-normal and BCCT squared-Hodge carriers have no typed full-role incidence
from $E_{P3}$.  These are domain failures, not missing numerical gains.

The later P3 reconstruction reduces the premise without changing this no-go.
On the declared enriched completion, the canonical exterior--record
realization fixes the P3 matrix carrier, its joint record occupation, the
minimal cyclic/GNS representation and the central trace weights. Contractible
unloaded Schur-neutral presentation directions are quotient-equivalent. These
facts remove four apparent selectors, but they begin after the typed positive
P3 quotient and its metric have been supplied. They do not derive the
base--seed tangent map $J_*$, its P1 and observer restrictions, the exact
closure of a loaded or typed-visible complementary internal sector, or an
isolated rank-two ROOT cluster of the full observer-shorted Hessian. Thus the
remaining source packet is finite and explicit:

\begin{equation}
\boxed{
\begin{gathered}
J_*,\quad \text{exact complement closure},\quad
\operatorname{gap}_{\rm ROOT}(K_{*,O})>0,\\
\text{physical realization of the enriched completion}
\end{gathered}}
\label{eq:p3premisereduction}
\end{equation}

Once this packet is source-owned, Eq.~\eqref{eq:mixedgram} evaluates the mixed
incidence and Eq.~\eqref{eq:stagedgluerank} evaluates the stacked rank. A new
P3 matrix realization, gain or terminal-specific selector cannot supply the
missing value.

In particular, the later quantum factor-designation result does not close
this source problem.  Its observer-event incidence is defined only on an
already adopted quantum event quotient: under exact recovery it selects one
logical $M_2$ embedding, after which canonical descent fixes the complementary
$M_4$ realization algebra.  That map is neither $J_*$ nor the complete
body--observer incidence into the joint P1--P3 recovery ideal.  The domains
and codomains differ, so factor designation cannot imply the missing
rest-domain rank or full typed surjectivity.

The mode-pair premise and the nonzero-overlap premise admit a sharper finite
certificate.  Let $P_O$ be the occupied rank-two projector on the selected
one-particle sector and let $J_{\rm scr}$ be the self-adjoint generator of
rotations of the oriented observer screen.  In a $J_{\rm scr}$ eigenbasis
write
\begin{equation}
P_O C P_O=
\begin{pmatrix}p_m&c\\c^*&p_{m+2}\end{pmatrix}.
\label{eq:rootttcluster}
\end{equation}

\begin{proposition}[Finite ROOT--TT activation certificate]
\label{prop:rootttcertificate}
Within the declared number-preserving quadratic matter sector, the selected
rank-two cluster is a covariantly closed source of a nonzero two-component TT
mean of real rank two if and only if
\begin{equation}
 [P_O,J_{\rm scr}]=0,\qquad
 \operatorname{spec}(P_OJ_{\rm scr}P_O)=\{m,m+2\},\qquad
 c\neq0,\qquad K_{m,m+2}\neq0.
\label{eq:rootttcertificate}
\end{equation}
For the opposite helicity the spectral gap is $-2$.
\end{proposition}

\begin{proof}
Global phase invariance leaves the quadratic coherences
$C_{mn}=\langle a_m^\dagger a_n\rangle$.  Under a screen rotation they carry
character $n-m$, whereas $T_++iT_\times$ carries character $+2$.
Orthogonality of $SO(2)$ characters therefore removes every off-diagonal
coherence except $n-m=2$ (or $-2$ for the conjugate helicity).  Projector
commutation makes these weights well defined inside the occupied cluster.
Equation~\eqref{eq:rootttidentity} then reduces the TT mean to the complex
product $K_{m,m+2}c$.  Multiplication by a nonzero complex number has real
matrix
$\left(\begin{smallmatrix}\Re K&-\Im K\\\Im K&\Re K\end{smallmatrix}\right)$,
whose determinant is $|K|^2$; hence the response has real rank two exactly
when both factors are nonzero.  If projector invariance fails, the selected
subspace is not a covariantly closed two-mode source.  Within an invariant
cluster, failure of the weight, coherence or overlap condition makes the
selected quadratic TT mean vanish.
\end{proof}

Four controls delimit the statement.  First, a single geometric-optics null
ray has stress $I k\otimes k$ and zero STF projection on its own Sachs screen;
non-collinearity relative to a common screen is therefore essential.  For a
ray with transverse components $(p_1,p_2)$,
\begin{equation}
T_+ + iT_\times=\frac{I}{2}(p_1+ip_2)^2,
\end{equation}
whose real Jacobian has rank two away from $p_1=p_2=0$.  Second,
$\rho_{\rm ROOT}$ alone cannot determine the stress magnitude: preparations
with the same normalized state and different $N_O$ have different stress.
Third, an incoherent classical anisotropic stress can be physically relevant,
but it does not establish the coherence-controlled identity above.
Fourth, outside the declared number-preserving sector the full covariance
packet, including anomalous moments, replaces $C$;
Eq.~\eqref{eq:rootttidentity} must not be extrapolated to a generic QFT state.

There is also an explicit nonzero witness for the final condition in
Eq.~\eqref{eq:rootttcertificate}.  On a patch where an isolated ray has fixed
nonzero transverse direction $p_1+ip_2$ relative to the common screen,
positive intensity $I$ and a nonnegative smearing function $f_O$ with
nonzero integral, its smeared STF coefficient is proportional to
\begin{equation}
 K_{\rm ray}={1\over2}(p_1+ip_2)^2
 \int d\mu_O\,f_O I\neq0.
\label{eq:nonzerokwitness}
\end{equation}
This proves that the admissible nonzero-overlap class is nonempty.  It does
not show that an arbitrary multimode packet avoids destructive cancellation,
or that the physical parent selects this witness.

Proposition~\ref{prop:roottt} is a conditional common-source bridge, not an
unrestricted parent-realization theorem.  This paper does not prove that the physical
base--seed system occupies the required non-collinear $\Delta m=2$ pair, that
every ROOT carrier has this two-mode realization, or that geometry is
quantized.

\subsection{Minimal invariant activation potential}

The four activation conditions need not be independent switches.  Let $p$ be
a weight-one complex transverse coordinate on the common screen and $c$ the
weight-two relative coherence.  Under screen rotations
$p\mapsto e^{i\theta}p$ and $c\mapsto e^{2i\theta}c$, the
lowest-total-degree analytic real mixed invariant is unique up to one
coefficient and complex conjugation, $\operatorname{Re}(c^*p^2)$.
Higher-order invariants are not excluded.  Consider
\begin{equation}
\mathcal V(c,p)=r_c|c|^2+u_c|c|^4+r_p|p|^2+u_p|p|^4
-2\lambda\operatorname{Re}(c^*p^2),
\label{eq:ttpotential}
\end{equation}
where $u_c,u_p,r_c>0$, $r_p<0$ and
$\lambda\in\mathbb R\setminus\{0\}$.

\begin{proposition}[Conditional minimal activation]
\label{prop:ttpotential}
Every global minimizer of Eq.~\eqref{eq:ttpotential} has $p_*\neq0$ and
$c_*\neq0$, with the relative phase locked by
$\arg c_*=2\arg p_*$ up to the sign of $\lambda$.  If the once-counted STF
overlap is $K(p)=\kappa p^2$ with $\kappa\neq0$, the full ROOT--TT activation
certificate holds.
\end{proposition}

\begin{proof}
The positive quartics make $\mathcal V$ coercive.  At $p=0$, its unique
minimum is $c=0$ with value zero.  Along $c=0$, however,
$|p|^2=-r_p/(2u_p)$ gives $\mathcal V=-r_p^2/(4u_p)<0$, so a global minimizer
cannot have $p=0$.  At $p\neq0$, the $c^*$ variation evaluated at $c=0$ is
$-\lambda p^2\neq0$, so $c=0$ is not stationary.  Phase minimization gives
the doubled-angle lock.  The square $p^2$ has screen weight two and therefore
selects the $\Delta m=2$ coherence, while $p_*\neq0$ and $\kappa\neq0$ give
nonzero STF overlap.
\end{proof}

This proposition closes activation only inside the declared minimal
completion.  The current physical reduct does not determine $r_p<0$,
$\lambda\neq0$ or $\kappa\neq0$.  In particular, $\lambda=0$ is an exact
same-sector control with $c_*=0$.

\subsection{Source-loaded activation from an occupied Abelian field}

The spontaneous branch above is not the only realization.  If a nonzero
transverse coordinate $p_{\rm src}$ is already supplied by the occupied
source ledger, minimize only
\begin{equation}
\mathcal V_{p_{\rm src}}(c)=r_c|c|^2+u_c|c|^4
-2\lambda\operatorname{Re}(c^*p_{\rm src}^2),
\qquad r_c,u_c>0.
\label{eq:loadedpotential}
\end{equation}

\begin{proposition}[Source-loaded coherence selector]
If $p_{\rm src}\neq0$ and $\lambda\neq0$, Eq.~\eqref{eq:loadedpotential}
has a unique minimizing amplitude and it obeys
\begin{equation}
(r_c+2u_c|c_*|^2)c_*=\lambda p_{\rm src}^2,
\qquad c_*\neq0.
\label{eq:loadedstationary}
\end{equation}
Its phase has screen weight two.  Therefore this branch closes the coherence
condition of Proposition~\ref{prop:rootttcertificate} without assuming
$r_p<0$.
\end{proposition}

\begin{proof}
For fixed phase, the radial factor on the left side of
Eq.~\eqref{eq:loadedstationary} is strictly increasing from zero to infinity.
It consequently has exactly one positive solution for every nonzero right
side, and phase minimization aligns $c_*$ with
$\lambda p_{\rm src}^2$.  Since $p_{\rm src}$ has screen weight one, its
square has weight two.
\end{proof}

The mixer in Eq.~\eqref{eq:loadedpotential} is also fixed once geometry and
stress are varied in the same functional.  With the complex STF convention
$q_h=h_++ih_\times$ and $q_T=T_++iT_\times$, write
\begin{equation}
q_h=\beta p_{\rm src}^2,
\qquad q_T=Kc.
\label{eq:stfpair}
\end{equation}

\begin{proposition}[Metric-variation ownership of the mixer]
\label{prop:mixownership}
Suppose the matter stress is defined by variation of the same once-counted
matter functional used for the ROOT preparation,
\begin{equation}
\delta_g\Gamma_{\rm matter}
=-{1\over2}\int d^4x\sqrt{-g}\,T^{\mu\nu}h_{\mu\nu}.
\label{eq:metricvariation}
\end{equation}
On the common screen, Eq.~\eqref{eq:stfpair} generates a mixed term
\begin{equation}
-g_{\rm TT}\operatorname{Re}(\beta K^*c^*p_{\rm src}^2),
\qquad g_{\rm TT}>0,
\label{eq:ownedmixer}
\end{equation}
where $g_{\rm TT}$ is fixed by the declared screen and smearing
normalization.  After the common source phase is fixed,
\begin{equation}
\lambda={g_{\rm TT}\over2}|\beta K|.
\label{eq:lambdafixed}
\end{equation}
Thus the mixer is nonzero exactly when the geometric TT pullback and the
once-counted stress overlap are both nonzero; it is not an independent
cross-terminal gain.
\end{proposition}

\begin{proof}
For two real STF screen tensors, their contraction is proportional to
$\operatorname{Re}(q_T^*q_h)$.  Substitution of
Eq.~\eqref{eq:stfpair} into Eq.~\eqref{eq:metricvariation} gives
Eq.~\eqref{eq:ownedmixer}; the positive proportionality factor contains only
the frozen screen, measure and normalization conventions.  A common phase
choice puts its coefficient in the real form used in
Eq.~\eqref{eq:loadedpotential}, yielding Eq.~\eqref{eq:lambdafixed}.
Adding another multiplier would count the same metric variation twice.
\end{proof}

The Abelian field supplies a canonical candidate for $p_{\rm src}$.  Define
\begin{equation}
 I_F={1\over2}F_{ab}F^{ab},\qquad
 J_F={1\over2}F_{ab}({}^*F)^{ab},\qquad
 \Delta_F=I_F^2+J_F^2.
\label{eq:faradayinvariants}
\end{equation}
For $\Delta_F>0$, the non-null Faraday field has an unordered pair of
distinct principal null directions $\{[k_+],[k_-]\}$
\citep{WyllemanCostaNatario2020}.  Given observer $u$, line-of-sight
$q=u+n$ and its oriented screen projector $S_{(u,q)}$, set
\begin{equation}
p_\pm={S_{(u,q)}k_\pm\over-u\cdot k_\pm},\qquad
\chi_{\rm dir}=|\varepsilon_S(p_+,p_-)|.
\label{eq:directionallift}
\end{equation}
These expressions are invariant under rescaling or exchange of the null
representatives.

\begin{proposition}[Conditional Faraday directional selector]
On the non-null Abelian branch, $\chi_{\rm dir}>0$ supplies two linearly
independent common-screen directions and hence a nonzero transverse source
for the loaded coherence selector.  For a non-principal observer, writing
$k_\pm/(-u\cdot k_\pm)=u+s_\pm$ gives
\begin{equation}
\chi_{\rm dir}=|n\cdot(s_+\times s_-)|,
\label{eq:directionaltriple}
\end{equation}
so the directional condition fails only for sight lines on the great circle
spanned by $s_+$ and $s_-$.  A principal observer, a null field
$\Delta_F=0$, and an aligned sight line are exact zero controls.
\end{proposition}

\begin{proof}
The non-null electromagnetic classification gives two distinct principal
null lines.  Screen projection and normalization produce
Eq.~\eqref{eq:directionallift}; its two-dimensional determinant is nonzero
exactly for linearly independent projections.  Direct projection into
$n^\perp$ gives the scalar triple product
Eq.~\eqref{eq:directionaltriple}.  For a non-principal observer
$s_+\times s_-\neq0$, and its orthogonal sight lines form one great circle.
For a principal observer $s_-=-s_+$, while a null field has one repeated
principal null line, so both controls give zero.
\end{proof}

This removes an arbitrary directional-pair choice \emph{after} occupation of
the non-null Abelian branch.  It does not prove that the Tau parent selects
that branch, that the QFT preparation lift identifies its support with the
two-mode ROOT cluster, or that every observer lies outside the exact control
sets.

\subsection{Faithful common-completion occupation}

The hub's common-completion criterion can be stated directly in the present
ledger.  Let $K_*>0$ be the true-null-quotiented body Hessian and let
$R_{\rm rec}^{\rm dedup}$ be the de-duplicated primitive recovery incidence
into the required finite recovery ideal.  Define
\begin{equation}
G_*^{\rm rec}=R_{\rm rec}^{\rm dedup}K_*^{-1}
(R_{\rm rec}^{\rm dedup})^\dagger,
\qquad
\rho_*^{\rm rec}={G_*^{\rm rec}\over\operatorname{Tr}G_*^{\rm rec}}.
\label{eq:faithfulcompletion}
\end{equation}

\begin{proposition}[Conditional faithful occupation reduction]
\label{prop:faithfuloccupation}
If $R_{\rm rec}^{\rm dedup}$ is surjective, then
$\rho_*^{\rm rec}$ is faithful on the required recovery ideal and its cyclic
GNS representation is faithful and unique up to unitary equivalence.
In particular, for
every nonzero positive central projector $P_A$ of the required Abelian ideal,
\begin{equation}
w_A:=\operatorname{Tr}(\rho_*^{\rm rec}P_A)>0.
\label{eq:abelianweight}
\end{equation}
Thus Abelian block occupation requires no independent terminal occupation
bit.  If a stable ROOT projector $P_O$ has nonzero weight in the same
recovered QFT preparation $\Gamma_O$, then
\begin{equation}
\rho_{\rm ROOT}
={P_O\Gamma_OP_O\over\operatorname{Tr}(P_O\Gamma_OP_O)}
\label{eq:samestaterestriction}
\end{equation}
and the ROOT state and Hilbert stress are automatically restrictions of one
preparation.
\end{proposition}

\begin{proof}
Surjectivity makes $G_*^{\rm rec}$ positive definite on the recovery ideal,
so its normalized state is faithful.  A faithful state is strictly positive
on every nonzero positive element, giving Eq.~\eqref{eq:abelianweight}.
Faithfulness of the state makes the corresponding finite-dimensional GNS
representation faithful, while the standard GNS uniqueness theorem fixes its
cyclic realization up to a unitary intertwiner.  Finally,
Eq.~\eqref{eq:samestaterestriction} is the normalized compression of the same
$\Gamma_O$ whose metric variation defines the stress; hence no second ROOT
preparation is introduced.
\end{proof}

This closes sector occupation and common-state ownership \emph{conditional}
on one surjective recovered branch.  It does not imply
$\Delta_F>0$ at every spacetime point: faithfulness of the Abelian ideal is a
support statement, whereas null versus non-null Faraday type is a local
state/solution property.  The remaining unrestricted physical problem is
therefore finite: the enriched parent law must realize a unique stable branch with
surjective $R_{\rm rec}^{\rm dedup}$, and the selected observer context must
meet the non-null and nondegenerate-screen conditions above.  The current
reduct does not entail that surjectivity; Proposition~\ref{prop:samereduct}
supplies its deficient control.

\subsection{Exact source-rank reduction}

The hub also supplies an exact factorization of the recovery incidence.  Let
$X_{\rm src}^{\rm lin}$ contain the body-owned and irreducible
observer--source inputs after the true-null quotient, let
$E_{\rm req}^{\rm lin}$ be the declared linearized recovery carrier, and write
\begin{equation}
J_{\rm phys}^{\rm lin}:X_{\rm src}^{\rm lin}
\longrightarrow E_{\rm MVP}^{\rm lin},
\qquad
R_{\rm rec}^{\rm lin}=U_{\rm rec}D_\Sigma J_{\rm phys}^{\rm lin}.
\label{eq:recoveryfactorization}
\end{equation}
Here $D_\Sigma$ is the ordered lower-triangular response map and $U_{\rm rec}$
implements source-owned de-duplication.

\begin{proposition}[Exact source-rank criterion]
\label{prop:sourcerank}
Suppose $D_\Sigma$ and $U_{\rm rec}$ are isomorphisms on their declared
physical quotients and $H_{\rm src}>0$ is the once-counted source Hessian.
Then
\begin{equation}
\operatorname{rank}R_{\rm rec}^{\rm lin}
=\operatorname{rank}J_{\rm phys}^{\rm lin},
\qquad
\ker R_{\rm rec}^{\rm lin}=\ker J_{\rm phys}^{\rm lin},
\quad
H_{\rm phys}=(J_{\rm phys}^{\rm lin})^\dagger
H_{\rm src}J_{\rm phys}^{\rm lin}.
\label{eq:sourcerankcriterion}
\end{equation}
Consequently the realized branch is strictly stable and its differential
recovery has zero cokernel if and only if $J_{\rm phys}^{\rm lin}$ is an
isomorphism. Nonzero projection into every typed role is not sufficient.
\end{proposition}

\begin{proof}
Left composition by the two isomorphisms preserves rank and kernel, proving
the first two identities.  Moreover
$\langle x,H_{\rm phys}x\rangle
=\langle J_{\rm phys}^{\rm lin}x,H_{\rm src}
J_{\rm phys}^{\rm lin}x\rangle$, so positivity of
$H_{\rm src}$ makes strict stability equivalent to injectivity of
$J_{\rm phys}^{\rm lin}$. Surjectivity of the linearized recovery map is
equivalent to surjectivity of $J_{\rm phys}^{\rm lin}$. Both properties hold
exactly when the latter map is an isomorphism.
\end{proof}

A full-column-rank map from a lower-dimensional source quotient may have a
nonzero component in every required role and yield $H_{\rm phys}>0$, while
still having nonzero cokernel.  This is the exact stable-but-incomplete
control.  Conversely, the hub's occupied-P3 result supplies a nonempty
faithful ROOT-edge subrepresentation inside the declared enriched MVP, but a
faithful subrepresentation is not global surjectivity of
$J_{\rm phys}^{\rm lin}$.
Thus the open physical statement is no longer a list of terminal choices: it
is whether the enriched base--seed/body plus irreducible observer--source law
selects one complete differential incidence with zero kernel and zero
cokernel, while the same packet covers the finite record support and acts
faithfully on the physical algebra quotient.

Write
\begin{equation}
J_{\rm phys}^{\rm lin}=[J_{BS},J_{OS}^{\rm irr}],
\qquad
P_{BS}=J_{BS}J_{BS}^{+},
\qquad
C_{OS}=(I-P_{BS})J_{OS}^{\rm irr}.
\label{eq:oscomplement}
\end{equation}

\begin{proposition}[Observer--source complement rank]
\label{prop:oscomplement}
On a source-frozen metric chart,
\begin{equation}
\operatorname{rank}J_{\rm phys}^{\rm lin}
=\operatorname{rank}J_{BS}+\operatorname{rank}C_{OS}.
\label{eq:osrank}
\end{equation}
If $N_{\rm lin}=\dim E_{\rm MVP}^{\rm lin}$, differential recovery has zero
cokernel exactly when
\begin{equation}
\operatorname{rank}C_{OS}
=N_{\rm lin}-\operatorname{rank}J_{BS}.
\label{eq:osfill}
\end{equation}
\end{proposition}

\begin{proof}
Split the target orthogonally as $\operatorname{im}J_{BS}$ and its
orthogonal complement. Column operations remove the component of
$J_{OS}^{\rm irr}$ inside the first summand; its remaining component is
$C_{OS}$ and has orthogonal image. The ranks therefore add, and
Eq.~\eqref{eq:osfill} is the zero-cokernel condition.
\end{proof}

The hub's WR-T113/WR-T117 construction evaluates this complement on the
retained six-row PBSE subsystem.  Its base plus four access rows have rank
five, while the irreducible observer--source row is outside their span and
raises the rank to six.  Thus, whenever this subsystem embeds faithfully into
the Paper-VII linearized carrier,
\begin{equation}
\operatorname{rank}C_{OS}\ge1.
\label{eq:pbselowerbound}
\end{equation}
This rules out a purely base--seed factorization on that retained subsystem.
It does not determine $N_{\rm lin}$ or prove Eq.~\eqref{eq:osfill} for the
complete carrier.  Moreover differential rank cannot replace finite record
coverage or algebraic faithfulness; the three conditions are jointly required
for full typed recovery.

\subsection{De-duplicated primitive recovery dimension}

The broad P1 incidence $R_{P1}$ already includes the common body--matter and
causal/time lanes, together with metric/gravity, radiation and observer-screen
differentials.
On the source-minimal nonzero P1 branch its rank is four.  Let $R_Q$ be the
real linearization of the occupied two-mode ROOT incidence on the same frozen
source quotient and define the source-space row projector
\begin{equation}
P_{P1}=R_{P1}^{+}R_{P1}.
\label{eq:p1rowprojector}
\end{equation}

\begin{proposition}[De-duplicated P1--ROOT rank]
\label{prop:deduprank}
For the adopted P1--ROOT--gauge packet,
\begin{equation}
R_{\rm rec}^{\rm lin}
=\begin{pmatrix}R_{P1}\\ R_Q(I-P_{P1})\end{pmatrix},
\qquad
\operatorname{rank}R_{\rm rec}^{\rm lin}=4+r_{Q\perp},
\label{eq:deduprank}
\end{equation}
where
\begin{equation}
r_{Q\perp}
=\operatorname{rank}_{\mathbb R}R_Q(I-P_{P1})
=\operatorname{rank}_{\mathbb R}R_Q
-\dim_{\mathbb R}(\operatorname{row}R_Q
\cap\operatorname{row}R_{P1}),
\qquad 0\le r_{Q\perp}\le4.
\label{eq:qperprank}
\end{equation}
An invertible gauge transformation on the retained charged/polarization
carrier changes the presentation but adds no recovery rank.
\end{proposition}

\begin{proof}
The second block in Eq.~\eqref{eq:deduprank} lies in the orthogonal complement
of the row space of $R_{P1}$, so the two ranks add.  Restricting the orthogonal
projection $I-P_{P1}$ to $\operatorname{row}R_Q$ gives
Eq.~\eqref{eq:qperprank} by rank--nullity.  The real dimension of the two-mode
complex ROOT packet is at most four.  Left composition by an invertible gauge
map preserves both row-space dimension and stacked rank.
\end{proof}

Separate marginal P1 and ROOT metrics do not determine $r_{Q\perp}$: equal
identity marginal Grams admit row-space intersections of real dimension four,
two or zero, and hence increments $0,2,4$.  The missing datum is not another
gain but the common-source mixed Gram.  With a positive source Hessian $K_*$,
observer access $A_O$, isolated occupied ROOT projector $P_Q$ and faithful
ROOT restriction $C_Q$, it is
\begin{equation}
G_{P1,Q}=R_{P1}K_*^{-1}A_O^\dagger P_QC_Q^\dagger.
\label{eq:mixedgram}
\end{equation}
After normalizing by the two marginal Grams, the singular values are the
cosines of the principal angles between the whitened row spaces.  Therefore
\begin{equation}
r_{Q\perp}=\operatorname{rank}R_Q
-\#\{j:\sigma_j(G_{P1}^{-1/2}G_{P1,Q}G_Q^{-1/2})=1\}.
\label{eq:principalanglerank}
\end{equation}
Thus the differential recovery rank is confined to $4\le
\operatorname{rank}R_{\rm rec}^{\rm lin}\le8$, but no unique value follows
until the enriched physical source selects the common packet
$(K_*,R_{P1},A_O,P_Q,C_Q)$.  The PBSE transverse row proves genuine
observer--source novelty on its retained subsystem; it must not be added again
to Eq.~\eqref{eq:deduprank} unless an independent embedding outside the broad
P1 row space is proved.

\begin{proposition}[Conditional seed-presentation selector]
\label{prop:seedselector}
Let the universal seed own a complete finite involutive presentation
$\{r_\alpha\}$ of the de-duplicated primitive role algebra and let
\begin{equation}
\mathcal A_{\rm rep}(j)
={\mathcal A_*\over2}\sum_\alpha
\|r_\alpha(j)\|_{G_\alpha}^2.
\label{eq:seedrepaction}
\end{equation}
Assume: (i) its zero set is one typed-unitary orbit represented by
$\Pi_{\rm phys}$; (ii) the relation derivative is injective on the
gauge/true-null quotient, so $K_*=D_r^*GD_r>0$; (iii) $R_{P1}$ and $A_O$ are
fixed typed projections of the same represented tangent
$J_*=D\mathcal E_{\Pi_{\rm phys}}$; and (iv) the shorted operator
\begin{equation}
K_{*,O}=(A_OK_*^{-1}A_O^*)^{-1}
\label{eq:shortedsource}
\end{equation}
has an isolated stable occupied complex rank-two cluster with Riesz projector
$P_Q$ and faithful represented ROOT restriction $C_Q$.  Then, up to the same
typed-unitary gauge, the seed representation selects
\begin{equation}
(K_*,R_{P1},A_O,P_Q,C_Q),
\qquad R_Q=C_QP_QA_O,
\label{eq:selectedpacket}
\end{equation}
and Eqs.~\eqref{eq:mixedgram}--\eqref{eq:principalanglerank} follow without
an additional terminal coefficient.
\end{proposition}

\begin{proof}
The single zero orbit fixes $\Pi_{\rm phys}$ up to typed unitary equivalence.
Differentiation of its represented generator evaluation and the fixed role
projections give $R_{P1}$ and $A_O$, while the Hessian of
Eq.~\eqref{eq:seedrepaction} at the zero gives $K_*$.  Functional calculus of
the isolated cluster fixes $P_Q$ covariantly, and faithfulness fixes $C_Q$ up
to a terminal frame.  Such a frame cancels from the ranks and mixed-Gram
principal angles, proving Eq.~\eqref{eq:selectedpacket} and the claim.
\end{proof}

This reduces five apparent inputs to two physical premises, but does not
prove them.  The current seed grammar does not yet supply a complete uniquely
represented P3 (or full eight-role) involutive presentation, and the physical
observer-shorted Hessian has not been proved to contain the required occupied
isolated cluster.  An occupied complete ROOT edge alone is insufficient: its
internal restriction can be held fixed while the complementary spectrum is
chosen either separated or interleaved.  The physical spectral proposition is
therefore undetermined, not false.

\section{P3 closure spine}
\label{sec:p3closurespine}

The detailed P3 typed-role extraction, O/Q and M/G stabilizer analysis and
cotangent-port construction are preserved in
Appendix~\ref{app:p3technical}.  The first theorem below reconstructs the
positive quotient
\[
 (E_{P3}^{\rm phys},K_*),\qquad \dim E_{P3}^{\rm phys}=7,
\]
from a smaller support--port packet.  The subsequent representation argument
then uses its canonical exterior action
\[
 \mathsf C_{\rm ext}(v)=\varepsilon(v)
 +\iota_{K_*(v,\cdot)},
\]
and the normalized two-leg record factor.  The appendix proves exactly which
complete first-jet, metric, typing and operational-port premises are required;
it does not import their physical base--seed occupation into the theorems
below.

The proof spine should be read in five steps.  First, the source-owned support
structure isolates two singleton roles and two unresolved role planes.
Second, ordered ports and the independent unit reconstruct a typed rank-seven
carrier, while the decoder preserves that rank in the body.  Third, the
class-minimal variational law occupies the reconstructed packet rather than
adding occupation as a separate postulate.  Fourth, the operator-square and
pointed-record terms produce the Clifford--record realization.  Finally, the
public pair test asks whether its source-whitened anticommutator relation can
be evaluated without terminal fitting.  Thus the sequence is
\[
 \begin{gathered}
 \text{support structure}
 \longrightarrow \text{rank-seven reconstruction}
 \longrightarrow \text{variational occupation}
 \\
 \longrightarrow \text{Clifford--record packet}
 \longrightarrow \text{empirical pair falsifier}.
 \end{gathered}
\]
The first four arrows are conditional internal results; the last is the
empirical decision frontier and currently remains data-ineligible.

\subsection{Tau-local origin of standard subsystem composition}
\label{sec:taustandardcomposition}

The two-leg record carrier used later should not be inserted merely because
standard four-dimensional notation uses a tensor product.  Let $H_0$ and
$H_1$ be finite occupied subhypergraphs of the one universal seed--body
realization.  Call them \emph{parent relation-free} when no occupied primitive
incidence meets both supports.  For this scope impose the following finite
parent-law clauses:
\begin{enumerate}[label=(LC\arabic*),nosep]
\item every primitive closure rule has connected source support contained in
  one occupied component;
\item a rule supported in $H_i$ has output support inside the component
  closure of $H_i$, so it creates neither a bridge to $H_{1-i}$ nor a shared
  external mediator;
\item the primitive component rules exhaust joint admissibility: there is no
  additional global accept/reject predicate on otherwise lawful endpoint
  pairs; and
\item the finite record linearization is source faithful and inherits the
  source-rigid unrescaled monoidal metric.
\end{enumerate}
The first three clauses concern closure-set factorization.  The fourth fixes
the common action normalization after factorization; it is not needed for the
abstract tensor shape.

\begin{theorem}[Tau-local standard subsystem composition]
\label{thm:taustandardcomposition}
Under LC1--LC3,
\begin{equation}
 \boxed{
 \mathrm{Cl}_\tau(H_0\sqcup H_1)
 \simeq
 \mathrm{Cl}_\tau(H_0)\times\mathrm{Cl}_\tau(H_1).}
 \label{eq:tauclosureproduct}
\end{equation}
Consequently the finite record algebras obey the canonical isomorphism
\begin{equation}
 \boxed{
 \mathbb C^{\mathrm{Cl}_\tau(H_0\sqcup H_1)}
 \simeq
 \mathbb C^{\mathrm{Cl}_\tau(H_0)}
 \otimes
 \mathbb C^{\mathrm{Cl}_\tau(H_1)}.}
 \label{eq:tautensorlinearization}
\end{equation}
Local parent operations have the standard commuting form
\begin{equation}
 T_0\longmapsto T_0\otimes I,\qquad
 T_1\longmapsto I\otimes T_1,\qquad
 [T_0\otimes I,I\otimes T_1]=0.
 \label{eq:taulocaloperators}
\end{equation}
Under LC4 the isomorphism in
Eq.~\eqref{eq:tautensorlinearization} is isometric in the inherited common
action unit, with no independent central tensor gain.
\end{theorem}

\begin{proof}
Take any finite global closure sequence starting from $H_0\sqcup H_1$.
LC1 assigns every primitive step to one component, and LC2 keeps its output
inside that component.  Projecting the sequence onto the two supports
therefore gives a pair of lawful local closure sequences and hence a map from
the left side of Eq.~\eqref{eq:tauclosureproduct} to the right side.

Conversely, rules on disjoint supports commute: a rule on $H_0$ changes
neither the match nor the output of a rule on $H_1$, and conversely.
Arbitrary local closure sequences can therefore be interleaved into a lawful
global sequence.  LC3 forbids a second global condition from deleting the
resulting endpoint pair.  The two constructions are inverse, proving
Eq.~\eqref{eq:tauclosureproduct}.

For finite sets $A$ and $B$, the delta-basis correspondence
\begin{equation}
 \delta_{(a,b)}\longleftrightarrow\delta_a\otimes\delta_b
 \label{eq:taudeltatensor}
\end{equation}
extends uniquely to the algebra isomorphism
$\mathbb C^{A\times B}\simeq\mathbb C^A\otimes\mathbb C^B$, proving
Eq.~\eqref{eq:tautensorlinearization}.  Component-local rules act on only one
delta factor, which gives Eq.~\eqref{eq:taulocaloperators}.  Finally LC4 makes
the primitive delta bases orthonormal in the inherited source metric and the
source-rigid monoidal continuation excludes a separate positive central
rescaling.  This proves the metric statement.
\end{proof}

\paragraph{Sharp controls and scope.}
Let both local alternative sets be $\{0,1\}$ while the global endpoint set is
only
\begin{equation}
 \mathcal D=\{(0,0),(1,1)\}.
 \label{eq:tauhiddendiagonal}
\end{equation}
Both marginal projections are onto, yet
$\dim\mathbb C^{\mathcal D}=2<4$; this hidden global relation violates LC3.
A crossing rewrite $(i,j)\mapsto(1-i,1-j)$ generates exactly the same
diagonal orbit and violates LC1--LC2.  Thus disjoint labels, complete
marginals, action valuation or a vanishing mixed Hessian do not imply
Eq.~\eqref{eq:tauclosureproduct}.  A positive central metric rescaling
preserves the tensor algebra but violates the isometric part of LC4.

The theorem is a Tau-only derivation inside the declared enriched
component-local parent class.  It does not show that the old-bare
admissibility skeleton or the unrestricted physical base--seed law realizes
LC1--LC4 on every claimed isolated scope.  If a crossing hyperedge is
occupied, the relation-free premise fails and this closure-factorization
theorem simply does not apply.  Retaining a tensor carrier as the kinematic
envelope of interacting quantum states then uses the separate strong
dagger-monoidal quantum representation result; it does not follow from
Eq.~\eqref{eq:tauclosureproduct} alone.

The full seven-role carrier need not be named in advance.  Let
$E_{\rm rel}$ be the positive six-dimensional relative source carrier and let
$P_{BMG},P_{OQ},P_{RMG}$ be commuting $K_{\rm rel}$-self-adjoint
projectors.  Write $E_{\epsilon}$ for their joint binary eigenspaces.

\begin{theorem}[Support--port reconstruction and class-minimal enriched-law realization]
\label{thm:p3supportportparent}
Assume
\begin{equation}
 \dim E_{100}=1,\qquad \dim E_{010}=2,\qquad
 \dim E_{001}=1,\qquad \dim E_{101}=2,
 \label{eq:p3supportmultiplicities}
\end{equation}
with every other joint eigenspace zero.  Assume further that the two
two-dimensional blocks possess source-owned quotient-basic ordered
cotangent decompositions
\begin{equation}
 E_{010}=L_O\oplus L_Q,
 \qquad E_{101}=L_M\oplus L_G,
 \label{eq:p3orderedports}
\end{equation}
and that the common-unit vector $z_U$ splits directly from $E_{\rm rel}$.
Then
\begin{equation}
 E_{P3}^{\rm phys}
 =L_B\oplus L_O\oplus L_R\oplus L_U
  \oplus L_Q\oplus L_M\oplus L_G,
 \qquad \dim E_{P3}^{\rm phys}=7,
 \label{eq:p3supportportrankseven}
\end{equation}
where $L_B=E_{100}$, $L_R=E_{001}$ and $L_U=\mathbb Rz_U$.
If a smooth body realization has a local left decoder, this rank is
preserved in the stabilized body.

Normalize the already stabilized body branch by
$\widehat\Gamma_{\rm body}
=\Gamma_{\rm body}-\Gamma_{\rm body}[M_*]\geq0$ on its declared stability
basin.  Write $\Phi_{\rm port}=\Gamma_{\rm port}^{\rm CTC}$ and normalize the
remaining post-body defects by
$\Phi_{P3}=\Gamma_{P3}^{\rm enr}/\mathcal A_*$ and
$\Phi_{\rm rec}=\Gamma_{\rm rec}/\mathcal A_*$.  Define the proposed
class-minimal enriched P3 law within the declared finite,
source-irredundant one-copy completion class by
\begin{equation}
 \boxed{
 \mathcal L_{P3}^{\min}
 =\left(
 \mathcal D_{\rm SF},
 \widehat\Gamma_{P3}^{\min}
 =\widehat\Gamma_{\rm body}
 +\mathcal A_*\bigl(
   \Phi_{\rm port}+\Phi_{P3}+\Phi_{\rm rec}\bigr);
 \operatorname{Real}_B,Q_B,\Omega_\varnothing;
 \mathsf{Phys}=\operatorname*{Argmin}_{\mathcal D_{\rm SF}}
 \widehat\Gamma_{P3}^{\min}
 \right),}
 \label{eq:p3minimalparentlaw}
\end{equation}
where $\mathcal D_{\rm SF}$ owns
Eq.~\eqref{eq:p3supportmultiplicities}, the unit split and the two protected
cotangent source covectors;
$\Gamma_{\rm port}^{\rm CTC}$ is the normalized two-block
Fenchel--Young law; $\Gamma_{P3}^{\rm enr}$ is the once-counted graph plus
operator-square law; $\Gamma_{\rm rec}$ is the normalized two-leg record
law; $Q_B\operatorname{Real}_B=I$; and $\Omega_\varnothing$ is the pointed
cyclic exterior root with the complete physical P3 null and no bypass.  All
three $\Phi$ terms are nonnegative dimensionless defects in the one common
action unit.  Then the physical branch selected by
Eq.~\eqref{eq:p3minimalparentlaw} satisfies
\begin{equation}
 \boxed{
 (B_\tau^{\rm phys},s_U,\mathcal L_{P3}^{\min})
 \Longrightarrow \mathfrak P_{P3}^{\min},}
 \label{eq:p3minimalparentarrow}
\end{equation}
where $\mathfrak P_{P3}^{\min}$ is the class-minimal faithful typed rank-seven
Clifford--record packet, unique up to typed source isometry and typed unitary
equivalence in the declared source-irredundant one-copy class.

The earlier narrow reduct does not entail
$\mathcal L_{P3}^{\min}$: an equal-upstream no-response completion omits a
required ordered port, while dependent-unit and equal-Gram commuting controls
independently remove the rank and operator-square conclusions.
\end{theorem}

\begin{proof}
Simultaneous orthogonal spectral decomposition of the three projectors gives
the intrinsic blocks $L_B\widehat\oplus E_{010}
\widehat\oplus L_R\widehat\oplus E_{101}$.  The two ordered cotangent
decompositions split the only repeated blocks, and the direct unit split adds
one independent line, proving Eq.~\eqref{eq:p3supportportrankseven}.
Differentiating the decoder identity gives
$DQ_BD\operatorname{Real}_B=I$, so body realization cannot collapse a typed
direction.  The two Fenchel--Young graph terms have zero response Schur
remainder.  The faithful separating operator-square term vanishes exactly on
$C(u)C(v)+C(v)C(u)=2K_*(u,v)I$, and eliminating the graph response leaves
that selector once.  The pointed exterior root and normalized record factor
then give $J_{\rm ext}^\dagger J_{\rm ext}=I_{256}$ and faithful cyclic
occupation.  The canonical exterior Clifford incidence, its graph response,
the two metric-dual ordered ports, the stabilized body and the normalized
two-leg record root form an explicit simultaneous zero of every term in
$\widehat\Gamma_{P3}^{\min}$.  Its global minimum is therefore zero.  Since
all summands are nonnegative, every global minimizer makes every summand
vanish; packet occupation follows from the declared variational law and is
not a separate common-zero premise.  Complete physical null and no bypass
give the stated
minimal-class uniqueness.  The three cited controls prove non-entailment from
the narrow reduct.  Appendix~\ref{app:p3technical} gives the componentwise
deletion certificate.  QED.
\end{proof}

Equation~\eqref{eq:p3minimalparentarrow} closes a formal ownership arrow only
after the proposed physical law is specified.  It does not prove that Nature
realizes this proposal, and it does not establish ambient-parent
exhaustivity from finite readouts.

\subsection{One common law for locality and the connected P3 packet}
\label{sec:jointlocalp3}

The locality theorem and the P3 occupation theorem act on different parts of
the source decomposition.  Let $\pi_0(s_U)$ denote the connected components
of occupied primitive incidence and define
\begin{equation}
 \boxed{
 \widehat\Gamma_{\rm CLP3}[X]
 =\sum_{C\in\pi_0(s_U)}\widehat\Gamma_C[X_C],
 \qquad K_{\rm CLP3}=\bigoplus_C K_C.}
 \label{eq:jointlocalp3action}
\end{equation}
Assume that primitive rules exhaust their full connected supports, preserve
component closure, and admit no additional cross-component multiplier or
hidden endpoint predicate; that every component action is nonnegative,
once-counted and has an attainable zero; and that realization is source
faithful in the common unrescaled action unit.  On the occupied connected
component $C_{P3}$ take $\widehat\Gamma_C$ to be the already defined
$\widehat\Gamma_{P3}^{\min}$ of
Eq.~\eqref{eq:p3minimalparentlaw}.

\begin{theorem}[Component-local P3 joint occupation]
\label{thm:jointlocalp3}
Under these hypotheses every global minimizer of
Eq.~\eqref{eq:jointlocalp3action} satisfies LC1--LC4 on every
parent-relation-free support pair and occupies the complete typed rank-seven
packet $\mathfrak P_{P3}^{\min}$ on $C_{P3}$.  The earlier narrow
base--seed/body reduct does not entail this common law.
\end{theorem}

\begin{proof}
Full-support assignment and component preservation give LC1--LC2.  Primitive
exhaustivity and absence of an extra cross-component multiplier give LC3,
while source-faithful unrescaled realization is LC4.  Hence
Theorem~\ref{thm:taustandardcomposition} applies to every relation-free pair.
Nonnegativity, disjoint component variables and attainability give
\begin{equation}
 \operatorname*{Argmin}\widehat\Gamma_{\rm CLP3}
 =\prod_C\operatorname*{Argmin}\widehat\Gamma_C,
 \qquad \min\widehat\Gamma_{\rm CLP3}=0.
 \label{eq:jointlocalp3argmin}
\end{equation}
Restriction to $C_{P3}$ is therefore a zero of
$\widehat\Gamma_{P3}^{\min}$, so
Theorem~\ref{thm:p3supportportparent} forces the complete rank-seven
Clifford--record packet.  Its internal incidences do not violate LC1--LC2:
they belong to one connected component and are never factorized by the
theorem.

For non-entailment, keep the local rules and complete marginals but impose the
hidden diagonal endpoint relation in Eq.~\eqref{eq:tauhiddendiagonal}; LC3
fails.  Independently replace the P3 realization by the equal-total-scalar
rank-six missing-port control used in the support--port deletion certificate.
Their product has the same earlier narrow descriptors while failing both
targets.  Thus that reduct cannot select
Eq.~\eqref{eq:jointlocalp3action}.
\end{proof}

This theorem closes compatibility and common-law occupation inside the
declared enriched completion.  It is not a proof that Nature realizes
$\widehat\Gamma_{\rm CLP3}$ or that the listed components exhaust the ambient
parent.

\subsection{Why the class-minimal law is not a terminal fit}

Here \emph{class-minimal} means componentwise minimal only within the
declared finite, source-irredundant one-copy completion class.  In both the
law and packet symbols, the superscript ``min'' has this restricted meaning;
it does not assert
that the law is Nature's unique or globally minimal parent law.  Nor is the
functional fitted backwards from terminal residuals: its inputs are frozen
on the source side before endpoint evaluation.  The origin of its clauses is
as follows.
\begin{enumerate}[label=(\roman*),nosep]
\item The support multiplicities and projectors are inherited from the
  source-role structure.  Their realization by the unrestricted physical
  base--seed system remains an explicit completion premise.
\item The O/Q and M/G cotangent ports are minimal preimages, within this
  class, of already required ordered dualities.  Deleting either port leaves
  an explicit typed swap symmetry.
\item The independent common unit and the left decoder have separate rank
  roles: the unit raises source rank from six to seven, while the decoder
  prevents collapse in body realization.  Neither is a terminal-specific
  gain.
\item The state-neutral operator-square term is the genuinely new
  Tau-specific cross-role clause.  Without it, a commuting packet can retain
  the same normalized marginal trace Gram.
\item The pointed cyclic exterior root and no-bypass clause are required by
  the faithful full-rank nontracial and extra-record countermodels.  They
  select record occupation and completeness only within the declared class.
\item The two-leg record tensor is not inferred from a terminal residual.
  Theorem~\ref{thm:taustandardcomposition} derives it from the independent
  component-local no-bridge closure law; the diagonal hidden-relation model
  prevents replacing that source law by complete marginals or action
  additivity.
\end{enumerate}
Thus Theorem~\ref{thm:p3supportportparent} is a noncircular internal
realization theorem for a source-frozen proposed law.  It is not a deeper
derivation of that law from the narrow parent, and it does not exclude richer
parent laws with the same physical zero locus.

\begin{theorem}[Pointed exterior-seed closure and rank-only no-go]
\label{thm:p3exteriorseedclosure}
Assume that the occupied typed P3 seed owns the complete oriented word
grammar $\Lambda^\bullet E_{P3}$, explicitly occupies the empty word
$\Omega_\varnothing$, and that the base realizes
$\mathsf C_{\rm ext}(v)$ as in Eq.~\eqref{eq:p3extclifford}.  Assume also that
the normalized two-leg Frobenius record root generates an orthonormal record
basis $r_0,r_1$ with no independent bypass.  Then the vectors
\begin{equation}
 \Psi_{I,a}
 =\mathsf C_{\rm ext}(e_{i_k})\cdots
  \mathsf C_{\rm ext}(e_{i_1})\Omega_\varnothing\otimes r_a,
 \qquad I=\{i_1<\cdots<i_k\}\subseteq\{1,\ldots,7\},
 \quad a\in\{0,1\},
 \label{eq:p3pointedorbit}
\end{equation}
form an orthonormal basis of $\mathcal F_{P3}$.  Equivalently, the root-orbit
map $J_{\rm ext}:\mathbb C^{256}\to\mathcal F_{P3}$ obeys
\begin{equation}
 \boxed{J_{\rm ext}^{\dagger}J_{\rm ext}=I_{256}.}
 \label{eq:p3exteriorgramlock}
\end{equation}
Consequently the relation, source-symmetry and cyclic trace--GNS conditions
are simultaneous kinematic consequences on this declared carrier; no
independent Gram-isometry action term is required there.

Full orbit rank, source nonloss and equal central block weights are not
sufficient substitutes.  On
$\mathcal A=M_8(\mathbb C)^{\oplus4}$ with normalized trace, choose
\begin{equation}
 h_\epsilon
 =\bigoplus_{b=1}^{4}
  \operatorname{diag}(1+\epsilon,1-\epsilon,1,\ldots,1),
 \qquad 0<\epsilon<1,
 \label{eq:p3nontracialdensity}
\end{equation}
and $J_\epsilon(a)=a h_\epsilon^{1/2}$ on the regular carrier.  Then
$J_\epsilon$ has rank 256 and equal weight $1/4$ on every central block, but
\begin{equation}
 J_\epsilon^\dagger J_\epsilon=R_{h_\epsilon}\ne I,
 \qquad
 \omega_\epsilon([E_{01}^{(1)},E_{10}^{(1)}])={\epsilon\over16}\ne0.
 \label{eq:p3rankonlycontrol}
\end{equation}
Thus this faithful full-rank root state is nontracial and retains a nonzero
metric defect.
\end{theorem}

\begin{proof}
Creation followed by $K_*$-adjoint contraction sends the empty exterior word
to the signed wedge basis vector labelled by the same subset $I$.  Distinct
subsets are orthogonal, and every subset occurs once.  Tensoring with the two
orthonormal record vectors therefore gives exactly $2^7\cdot2=256$
orthonormal vectors, proving Eq.~\eqref{eq:p3exteriorgramlock}.  The Clifford
relations give relation generation, the induced exterior action transports
every source isometry, and the full word orbit is cyclic; the normalized
regular vector state is the represented trace.

For the control, positivity of $h_\epsilon$ makes right multiplication by
$h_\epsilon^{1/2}$ invertible, so the orbit map has full rank.  Every block
has trace eight, hence the four central weights are equal.  Direct evaluation
in the normalized regular inner product gives
$J_\epsilon^\dagger J_\epsilon=R_{h_\epsilon}$.  Finally
$[E_{01},E_{10}]=E_{00}-E_{11}$, whose $h_\epsilon$-weighted normalized trace
in one of the four blocks is $2\epsilon/32=\epsilon/16$.  The root state is
therefore not tracial and its Gram cannot be the source identity. QED.
\end{proof}

The positive half of Theorem~\ref{thm:p3exteriorseedclosure} is conditional
on physical ownership of the complete pointed exterior grammar.  The
negative half is an exact countermodel: dimension, rank and the retained
discrete block symmetries cannot replace that source-owned pointed metric.
The unrestricted physical base--seed law has not yet been shown to own this
grammar.


Appendix~\ref{app:p3certificate} packages the same requirement as an exact
complete-incidence characterization: within the declared finite
source-irredundant Clifford--recovery class, a minimal faithful typed P3
subpacket exists uniquely up to typed unitary equivalence if and only if the
parent owns the complete rank-seven incidence.  This is a characterization of
the missing object, not a proof of its unrestricted physical existence.


The certificate also admits one finite positive source selector.  On the
dimensionless normalized quotient $(E_{P3}^{\rm phys},K_*)$, let
$n=\dim E_{P3}^{\rm phys}=7$, let $C:E_{P3}^{\rm phys}\to
\mathcal A_{\rm body}^{\rm sa}$ be linear, and let $\tau$ be a faithful
normalized trace.  Define
\begin{equation}
 \mathcal V_{\rm OS}[C]
 = {\mathcal A_*\over2}
 \int_{S(E_{P3},K_*)}
 \tau\!\left[
 (C(v)^2-I)^\dagger(C(v)^2-I)
 \right]d\mu(v),
 \label{eq:p3operatorsquarefunctional}
\end{equation}
where $d\mu$ is normalized orthogonal-invariant sphere measure.  Before
unit-sphere restriction the defect is $C(v)^2-K_*(v,v)I$, so the common
action unit $\mathcal A_*$ restores dimensions without a role-dependent gain.
For a $K_*$-orthonormal frame, write $A_a=C(e_a)$,
$d_a=A_a^2-I$, and $h_{ab}=A_aA_b+A_bA_a$.  The fourth sphere moments give
the exact finite formula
\begin{equation}
 {2\mathcal V_{\rm OS}[C]\over\mathcal A_*}
 ={1\over n(n+2)}\left[
 2\sum_a\lVert d_a\rVert_\tau^2
 +\left\lVert\sum_a d_a\right\rVert_\tau^2
 +\sum_{a<b}\lVert h_{ab}\rVert_\tau^2
 \right].
 \label{eq:p3operatorsquarefinite}
\end{equation}

\begin{theorem}[Operator-square source selector]
\label{thm:p3operatorsquareselector}
For a faithful normalized trace,
$\mathcal V_{\rm OS}[C]\geq0$, with equality if and only if
\begin{equation}
 C(v)^2=K_*(v,v)I
 \quad\hbox{for every }v\in E_{P3}^{\rm phys}.
 \label{eq:p3operatorsquarezero}
\end{equation}
Equivalently,
$C(u)C(v)+C(v)C(u)=2K_*(u,v)I$.  The zero set is invariant under every
$K_*$-orthogonal source-frame change.
\end{theorem}

\begin{proof}
The integrand in Eq.~\eqref{eq:p3operatorsquarefunctional} is continuous and
nonnegative.  Faithfulness of $\tau$ makes a zero integral equivalent to a
zero operator defect on the unit sphere; homogeneity extends the result to
all $v$.  Polarization gives the anticommutator, and the reverse implication
is immediate.  Orthogonal invariance follows from invariance of $K_*$ and
$d\mu$.  Finally,
$\mathbb E[v_a^4]=3/[n(n+2)]$,
$\mathbb E[v_a^2v_b^2]=1/[n(n+2)]$, and the odd moments vanish, yielding
Eq.~\eqref{eq:p3operatorsquarefinite}.
\end{proof}

This selector contains information absent from the scalar Gram matrix.  The
canonical Clifford family and seven commuting four-qubit Pauli-$Z$ words act
on the same 16-dimensional carrier and both obey
$\tau(A_aA_b)=\delta_{ab}$.  The finite certificate nevertheless gives
\begin{equation}
 {2\mathcal V_{\rm OS}^{\rm Clifford}\over\mathcal A_*}=0,
 \qquad
 {2\mathcal V_{\rm OS}^{\rm commuting}\over\mathcal A_*}={4\over3}.
 \label{eq:p3operatorsquarecontrol}
\end{equation}
Thus scalar Hessian/Gram data cannot select the physical packet, whereas the
operator-valued source term can.

Appendix~\ref{app:p3ownership} gives the staged enriched-parent functional
whose nonnegative zero locus combines this operator-square law with the P3
handoff without double counting.  That theorem proves ownership inside the
declared enriched completion.  It does not prove that the unrestricted
physical base--seed law realizes or occupies the enrichment.

\begin{theorem}[Conditional three-level enriched-MVP closure]
\label{thm:threelevelclosure}
Assume the dependency packet D1--D8 of
Appendix~\ref{app:dependencycertificate}, the occupied complete source
representation of Proposition~\ref{prop:moprassembly}, the occupation
branch of Proposition~\ref{prop:occupation} with cyclic $\Omega_s$ and
surjective required role maps, and the following P3 premises:
\begin{enumerate}[label=(\roman*),nosep]
\item the class-minimal support--port enriched law of
  Theorem~\ref{thm:p3supportportparent} governs the physical branch; and
\item the canonical operational-activation premises collected in
  Appendix~\ref{app:p3technical} hold on the claimed active support.
\end{enumerate}
Then the declared enriched MVP is closed at the following three distinct
levels:
\begin{enumerate}
\item \emph{common-source level:} every primitive terminal factors through
  one body-conditioned observer access and one typed source representative;
\item \emph{algebra/representation level:} the normalized P3 Clifford law,
  exterior--record carrier, faithful cyclic P3--record representation and
  common trace/action bookkeeping form one occupied packet, with no
  independent role gain or variance carrier;
\item \emph{terminal level:} the single once-counted effective functional
  gives the normalized Born effect derivative and the conditional
  Einstein--Maxwell--matter equations, Ward/Bianchi identities and standard
  weak-field leaf without an additional terminal gain.
\end{enumerate}
The conclusion is conditional on the proposed physical law and the inherited
common-source occupation premises stated above; no additional P3
common-zero occupation premise is required.  It does not prove that Nature
realizes that law and does
not derive the full observer duration law: Paper VII uses only the formal
clock response in D7, while the complete temporal descent is owned by Paper V.
\end{theorem}

\begin{proof}
Level 1 is Proposition~\ref{prop:moprassembly}; full occupation, when invoked,
is Proposition~\ref{prop:occupation}.  For Level 2,
Theorem~\ref{thm:p3supportportparent} supplies the typed rank-seven source
packet and its law-to-packet realization.  Appendix~\ref{app:p3technical},
Propositions 19--31, supplies the detailed
typed-role realization: the full-role Clifford algebra, the canonical
complete first-jet quotient and polar normalization, the exterior--record
carrier, the port-variance representation and the deletion certificate.  The
pointed-orbit theorem fixes the cyclic metric and trace weights, while the
occupied enriched-parent zero locus supplies the operator-square relation
rather than assuming it as an unlabelled terminal identity.  Together these
results give the faithful
cyclic representation, sector multiplicities, operational O/Q ports and
nonzero-stress M/G activation on the declared active support without an
additional occupation coefficient.
For Level 3, Theorem~\ref{thm:effectivesource} constructs the single
once-counted terminal functional, Eq.~\eqref{eq:sourceowners1} gives the Born
probabilities, and Theorem~\ref{thm:assembly} with its weak-field corollary
gives the joint gravitational, gauge and matter terminal equations.  All
three levels use the same source ledger and unit packet, so composing them
introduces no new blockwise normalization. QED.
\end{proof}

Theorem~\ref{thm:threelevelclosure} is a dependency-composition theorem.  Its
explicit common-source occupation and variational-law premises prevent the weaker
notions of nonzero activation or rank-seven typing from being mistaken for
physical occupation of the complete packet.


\section{Class-minimal enriched-law closure and empirical frontier}
\label{sec:selectionfrontier}

The earlier narrow reduct has the exact non-entailment
\begin{equation}
 \boxed{
 (B_\tau^{\rm phys},s_U,\mathcal L_{\rm narrow})
 \not\Longrightarrow
 \mathfrak P_{P3}^{\min}.
 }
 \label{eq:narrowparentnogo}
\end{equation}
The no-response port control already proves Eq.~\eqref{eq:narrowparentnogo};
the dependent-unit, equal-Gram commuting and nontracial-root controls locate
three further independent failures.  It is therefore impossible to remove
this no-go by another manipulation of the unchanged reduct.

Theorem~\ref{thm:p3supportportparent} instead supplies an explicit proposed
physical law and proves
\begin{equation}
 \boxed{
 (B_\tau^{\rm phys},s_U,\mathcal L_{P3}^{\min})
 \Longrightarrow \mathfrak P_{P3}^{\min}.
 }
 \label{eq:decisiveparentarrow}
\end{equation}
Thus the formal parent-law realization and occupation arrow is closed inside
the declared enriched physical MVP.  The phrase ``Nature realizes this law''
is not an additional conclusion that can be proved from an unspecified
parent: it is the empirical status of the proposed law.  Likewise,
Eq.~\eqref{eq:decisiveparentarrow} does not establish that the selected
connected packet exhausts every ambient source extension.

Two scientific routes remain, but they have different status:
\begin{enumerate}
\item \emph{Deeper-origin route.}  Derive
  $\mathcal L_{P3}^{\min}$ from a more primitive enriched base--seed law.  This
  could strengthen the theory, but is no longer required for the internal
  consistency of the stated completion; Eq.~\eqref{eq:narrowparentnogo}
  forbids obtaining it from the unchanged narrow reduct.  The standard
  subsystem-composition part of this route is now closed conditionally by
  Theorem~\ref{thm:taustandardcomposition}; the remaining deeper-origin debt
  concerns physical realization of LC1--LC4 and of the full support--port,
  operator-square and pointed-root law.
\item \emph{Discriminator route.}  Freeze a source-side morphological
  descriptor and derive a gain-free, parameter-free relation between at least
  two terminal readouts that is not implied by the corresponding standard
  comparator.  The immediate algebraic target is
  $\{C_a,C_b\}/2=\delta_{ab}I$ in a source-whitened frame; the stronger
  $\mathcal O_{P1|P3}=0$ requires an independently certified global
  source-completeness premise.  Any endpoint prediction must be fixed before
  data are scored.
\end{enumerate}
Failure of a source-frozen discriminator demotes the proposed physical law.
Neither route may be repaired by adding another standard-limit recovery or by
selecting source structure from terminal residuals.

\section{Downstream scalar-15/type-II closure}
\label{sec:scalar15handoff}

The P3 spine fixes the common assembly language but does not itself choose a
non-Abelian particle completion.  A later P4 branch closes three previously
separate questions at three different logical levels.  We record that result
here because it is the first particle-sector specialization of the joint
source architecture; the underlying P4 reconstruction is not reclassified as
a theorem owned by Paper VII.

\subsection{Class-relative base--seed ownership}

Let \(\mathcal C_{15}^{\rm rec}\) denote the recovery-complete,
source-generated and source-irredundant P4/type-II class.  It contains the
recovered one-copy PIS--CAR packet, one occupied Higgs branch, the scalar-15
terminal and no independently sourced terminal-null operation ideal.  The
reverse source reconstruction identifies the pointed exterior \(V_5\) seed
grammar and a positive Hermitian/Hodge base realization.  The one-carrier
type-II selector then fixes the typed edge
\begin{equation}
 e_{HH\Delta}:\operatorname{Sym}^2V_5\longrightarrow S_{15},
 \qquad \dim\operatorname{Sym}^2V_5=15,
 \label{eq:scalar15edge}
\end{equation}
and faithful realization gives rank fifteen.  The unique shortest connected
alternating seed--Higgs word and its positive quartic completion imply
\begin{equation}
 \boxed{
 \mathcal C_{15}^{\rm rec}
 \Longleftrightarrow
 (\mathcal S_{\wedge,5},e_{HH\Delta};h_B,*_B)
 \Longrightarrow h_*\ne0,\quad\Delta_*\ne0 .}
 \label{eq:scalar15ownership}
\end{equation}
Thus charged character, occupation and quartic stabilization require no
independent source datum \emph{inside this recovery class}.  Equation
\eqref{eq:scalar15ownership} is a minimal-preimage reconstruction of the
selected branch.  It is not a theorem that the unrestricted physical parent
occupies type II rather than another neutrino/particle completion.

\subsection{Common spectral scale}

The same completion relates its Higgs length and scalar-15 mass to the common
Planck terminal:
\begin{equation}
 \ell_H=\sqrt{48\pi}\,\ell_P,
 \qquad
 M_*={\sqrt2\,\hbar c\over\ell_H}
 ={E_P\over\sqrt{24\pi}}
 ={\overline M_Pc^2\over\sqrt3},
 \label{eq:scalar15scale}
\end{equation}
where \(E_P=\sqrt{\hbar c^5/G}\) and
\(\overline M_Pc^2=E_P/\sqrt{8\pi}\).  With the 2022 CODATA constants
\citep{CODATA2022},
\begin{equation}
 \boxed{M_*=1.406035(16)\times10^{18}\ {\rm GeV}.}
 \label{eq:scalar15numeric}
\end{equation}
The Tau-specific statement is the dimensionless ratio
\(M_*/E_P=1/\sqrt{24\pi}\) and the absence of an independent scalar-15 length.
Equation~\eqref{eq:scalar15numeric} is a one-anchor metrological calibration,
not a first-principles prediction of the SI value of \(G\).

\subsection{Exact interacting functional law}

Lift the finite particle operator to the full graded four-dimensional carrier
and let \(P\) be the positive background-covariant Laplace-type operator,
including the declared gauge-fixing and ghost blocks.  The no-extra-profile
spectral scheme uses
\begin{equation}
 \boxed{R_k(P)={P\over\exp(P/k^2)-1}}
 \label{eq:scalar15regulator}
\end{equation}
in the exact effective-average-action flow \citep{Wetterich1993}
\begin{equation}
 \boxed{
 \partial_t\Gamma_k={1\over2}\operatorname{STr}
 \left[(\Gamma_k^{(2)}+R_k)^{-1}\partial_tR_k\right],
 \qquad t=\log k .}
 \label{eq:scalar15flow}
\end{equation}
For gauge sectors Eq.~\eqref{eq:scalar15flow} is solved together with the
regulator-modified Slavnov--Taylor/BV identity; the regulator insertion must
vanish as \(k\to0\).  If \(\{\mathcal O_I\}\) is a radiatively and
BRST-complete independent operator basis and
\(G_{IJ}=\langle\mathcal O_I,\mathcal O_J\rangle\) on the source-owned trace
domain, the component beta functions are the unique projection
\begin{equation}
 \boxed{
 \beta^I=(G^{-1})^{IJ}
 \left\langle\mathcal O_J,\partial_t\Gamma_k\right\rangle .}
 \label{eq:scalar15betaprojection}
\end{equation}
Equations~\eqref{eq:scalar15regulator}--
\eqref{eq:scalar15betaprojection} fix one gauge-covariant exact functional
law, threshold prescription and coordinate projection.  Computing the full
component polynomials remains an evaluation problem.  The regulator is an
analysis scheme, not a physical coupling, observer quantizer or additional
law requiring Nature-level selection; exact \(k=0\) observables are regulator
independent, whereas residual dependence of a finite truncation is an
uncertainty diagnostic.

The net closure is therefore deliberately three-tiered:
\begin{center}
\begin{tabular}{@{}ll@{}}
\toprule
Question & Status\\
\midrule
base--seed ownership & proved inside \(\mathcal C_{15}^{\rm rec}\)\\
absolute numerical scale & dimensionless ratio derived; SI value calibrated once\\
full beta law & exact functional flow and projection fixed\\
physical type-II occupation & empirical model-selection question\\
\bottomrule
\end{tabular}
\end{center}


\section{Public-source eligibility audit}

The accompanying reproducibility package freezes a finite source search as of
7 August 2026.  A candidate is eligible for a future independent
action--quantum endpoint score only if all six predeclared gates hold:
\begin{enumerate}[nosep]
\item an independently reconstructed electromagnetic or morphological leg;
\item quantum-state coherence, phase or tomography;
\item same-system joint indexing;
\item an independent positive endpoint action, energy or work measurement;
\item a complete standard comparator; and
\item a candidate-specific operational realization of the source-frozen,
gain-free Tau cross-readout law on independently measured observables.
\end{enumerate}
The script \path{scripts/audit_public_source_eligibility.py} contains all
35 candidate identifiers, public links, gate values and textual verdicts and
regenerates
\path{data/derived/public_source_eligibility_audit.json}.  Zero of 35
candidates pass all six gates.  The gate-pass counts in the listed order are
$(30,25,22,11,35,0)$.  After
Theorem~\ref{thm:p3supportportparent}, the abstract Tau law itself is frozen;
the last zero now records the absence of a candidate-specific independent
operational map from the common physical source into its typed observables and
whitening metric.  Even before that operational gate is applied, zero
candidates pass all five data/comparator gates; the nearest candidates pass
four.  Thus the endpoint-data packet and its common-source realization remain
independently incomplete.  This finite audit is a reproducible
classification of the located sources, not an exhaustive theorem about all
existing experiments; individual gate assignments remain open to correction
when stronger public metadata or raw records appear.  Failure of source
eligibility is an acquisition and identifiability limitation, not a null Tau
signal.  The frozen query families, inclusion rule, source-version policy and
the absence of external payload hashes are recorded in
\path{data/public_source_search_protocol.md}.

Consequently the direct rank-seven Nature test is frozen but not yet
executable: no primary residual is scored on an ineligible packet.  The
available Wen path-phase reconstruction is only an auxiliary branch control;
its morphology increment beyond scalar action is not significant under the
shared-kernel null ($\Delta R^2=0.1758$, $p=0.144$) and does not test or
falsify the complete P3 relation.

\paragraph{Minimal two-role falsifier.}
A first rejection does not require simultaneous reconstruction of all seven
roles.  For any predeclared distinct typed pair $(a,b)$, the proposed minimal
P3 law necessarily restricts to
\begin{equation}
 C_a^2=C_b^2=I,
 \qquad
 \{C_a,C_b\}=0
 \label{eq:minimalpairtest}
\end{equation}
after the same source-frozen whitening used by the full test.  Therefore a
statistically resolved nonzero anticommutator on one eligible pair rejects the
minimal rank-seven completion.  Eligibility still requires two independent
typed source directions and an informationally complete reconstruction of the
symmetrized operator product; a single-state expectation value is not enough.
A passing pair does not establish the remaining roles or rank seven.  The
reproducible focused audit in
\path{scripts/audit_minimal_pair_nature_test.py} finds zero eligible packets
among fourteen targeted public architectures.  This $0/14$ result narrows the
acquisition requirement and supplies no Tau pass or failure.

The already derived ROOT--TT same-preparation identity is not an independent
realization of this pair.  At frozen occupation and overlap it has the form
$T_{\rm TT}=N_OK\rho_{01}$, so its real two-component Jacobian is an
invertible rescaling of the ROOT-coherence Jacobian.  The stacked Jacobian
therefore has rank two, equal to either complex terminal separately, and adds
zero cross-terminal novelty rank.  This factorization is a physical
cross-terminal prediction, but counting it as two primitive source directions
would duplicate one source jet.  The executable test consequently requires
two independently controlled typed perturbations and process or sequential
data that reconstruct their symmetrized cross-response.

Nor can complete tomography confined to the quantum terminal replace the
missing typed leg.  If $L_1=A_1L_Q$ and $L_2=A_2L_Q$, the two maps may describe
independent quantum gates or coordinates, but both still descend through the
same $Q$ role.  Relabelling one as $M$ or $G$ after endpoint inspection would
violate source freeze.  The audited higher-order Quantum SWITCH process matrix
and donor-spin gate-set tomography packages are therefore near candidates,
not eligible cross-role packets: both reconstruct rich quantum operations but
publish no independently sourced metric, stress or morphology perturbation.

A closer hybrid candidate is a public source-data package in which a
superconducting qubit is read through a piezo-opto-mechanical transducer.  It contains transduction
characterization, prepared-qubit IQ records, optical-power scans, and
pump-dependent $T_1/T_2^*$ families on one device.  The release nevertheless
separates these into figure-level aggregate families: it provides neither an
independently reconstructed mechanical/stress state nor shared record keys
for an informationally complete quantum--mechanical symmetrized product.
Optical-pump dependence is therefore a controlled environmental response, not
yet a second primitive typed map.  The package is an acquisition template,
not an executable P3 pair test.

A mechanical carrier does not by itself evade the same-sector no-go.  If
$L_{\rm state}=A_sL_Q$ and $L_{\rm mech}=A_mL_Q$, both reconstructions descend
through one quantum role even when the carrier is a mechanical oscillator.  A
calibrated displacement control is likewise not an independent output map.
The audited ETH motional-grid-qubit release contains rich characteristic,
Wigner, Pauli and process tomography, but all records reconstruct one motional
quantum carrier~\cite{Fluehmann2019}.  The NIST amplified-motion release is closer because it joins
calibrated displacement, squeezing and spin response on common parameter rows;
nevertheless the displacement is a control and the oscillator output is
inferred through the spin/Q readout~\cite{Burd2019}.  Neither release supplies an independently
reconstructed mechanical or stress output whose differential adds rank after
conditioning on $L_Q$, nor the required symmetrized cross-product record.
The fourteenth and physically closest public candidate couples a membrane
oscillator to a distant atomic-spin oscillator and reconstructs a joint
Gaussian EPR witness~\cite{Thomas2021}.  Its released source files expose one
combined optical trajectory or spectrum.  The separately named mechanical,
atomic and backaction spectra are fitted model components rather than
independently measured output maps.  The package therefore does not certify
typed rank two or the informationally complete symmetrized product.

\paragraph{Federated mosaic identifiability no-go.}
Several incomplete experiments can constrain one common source Gram only
after independently certified source-frozen transports place their measured
objects in the same carrier representation and action unit.  Let
$G\in\mathcal G$ be the declared source-whitened Gram, let
$g=\operatorname{vech}G$, and write the constraints from packet $k$ as
$y_k=A_kg$.  For $A=\operatorname{stack}_kA_k$, the Gram is unique exactly
when
\begin{equation}
 \ker A\cap\operatorname{vech}(\mathcal G-\mathcal G)=\{0\}.
 \label{eq:federatedgramcriterion}
\end{equation}
The reproducible audit in
\path{scripts/audit_federated_gram_completion.py} applies this criterion to
the fourteen frozen pair candidates.  Even optimistically granting both
normalized $Q$ and $M$ marginals, the design for
$g=(G_{QQ},G_{MM},G_{QM})$ has rank $2/3$ and leaves $G_{QM}$ free.  All
positive-semidefinite matrices
\begin{equation}
 G(\rho)=\begin{pmatrix}1&\rho\\\rho&1\end{pmatrix},
 \qquad -1\leq\rho\leq1,
 \label{eq:federatedpairfamily}
\end{equation}
therefore reproduce those marginals.  In the seven-role control, the seven
diagonals plus a connected six-edge overlap tree have rank $13/28$ and leave
fifteen symmetric-Gram directions unobserved.  Connectivity alone is not
identifiability; fewer than complete pair coverage can suffice only under
additional algebraic or low-rank identities derived independently of the
endpoint data.

Scalar Gram closure is itself only necessary.  The anticommuting pair
$(Z,X)$ and the commuting pair $(Z\otimes I,I\otimes Z)$ have the same
identity normalized-trace Gram, while their anticommutators differ.  Thus no
synthetic Nature packet or Tau score may be formed by stacking unrelated
marginals.  A valid mosaic still needs common-carrier transport certificates,
at least one independently measured cross-role edge for the scalar pair, and
process or sequential tomography for the operator product.  The finite
verdict is
\texttt{FEDERATED\_MARGINALS\_DO\_NOT\_IDENTIFY\_CROSS\_ROLE\_GRAM}; it is an
identifiability no-go, not a null Tau result.

The reproducibility package is not limited to claim-string checks.  The script
\path{scripts/audit_theorem_certificates.py} independently computes finite
witnesses for the Schur response, the nonzero invariant-potential minimum, the
ROOT--TT Jacobian rank, overlapping terminal maps, the distinction between
activation and full occupation, and the normalized same-reduct counterpair.
These calculations certify representative finite instances; the analytic
proofs remain the source of the universal statements.  The separate script
\path{scripts/audit_rank7_parent_law.py} reconstructs the support
multiplicities, the successive stabilizer reductions, the independent-unit
and decoder rank controls, and the Clifford/equal-Gram counterpair.  It also
verifies a simultaneous zero of the normalized nonnegative action components,
which is the finite certificate for the variational-occupation step in
Theorem~\ref{thm:p3supportportparent}.
The federated audit independently verifies the pair and seven-role design
ranks, positive-semidefinite completion interval, resolution stability and
the equal-Gram commuting operator countermodel; all thirteen checks are
generated from the packaged pair ledger.

\paragraph{Code and data availability.}
The manuscript source, deterministic audit scripts, tests, derived JSON
ledgers and the frozen public-source search protocol are distributed in the
reproducibility archive accompanying this manuscript.  No observational
endpoint data are redistributed or fitted here; external candidate records
are identified by public links in the frozen audit ledger.  Running
\texttt{python3 scripts/reproduce.py} rebuilds the figures, theorem witnesses,
PDF, arXiv source archive and test suite from the packaged files.

\section{Finite decision ledger}

\begin{longtable}{@{}p{0.36\textwidth}p{0.39\textwidth}p{0.18\textwidth}@{}}
\toprule
Question & Verdict & Status\\
\midrule
One simultaneous body/post-body Hessian supplies conditioning and freeze &
Impossible unless the cross derivative vanishes & Exact no-go\\
One ordered ledger supplies both & Yes, Eq.~\eqref{eq:triangular} & Conditional theorem\\
One post-body effective functional generates probability, stress, current
and clock response & Yes, Eqs.~\eqref{eq:ineffective}--\eqref{eq:sourceowners2}
& Conditional theorem\\
Green response is an additional action contribution & No; it is an inverse
Hessian & Exact bookkeeping result\\
Finite local counterterms may be included in two columns & No; relocation is
allowed, duplication is not & Exact bookkeeping result\\
Classical and quantum occupation may overlap silently & No; stress doubles & Exact no-go\\
Einstein, Maxwell and Ward equations share one post-body variation & Yes under the declared packet & Conditional theorem\\
Einstein--Hilbert and Maxwell action shapes follow inside the declared
four-dimensional IR classes & Yes, Proposition~\ref{prop:standardforms};
coefficients and parent-law realization remain open & Conditional classification\\
Newton/Maxwell limits follow & Yes on the standard weak leaf & Conditional recovery\\
Standard-limit agreement proves Tau provenance & No & Exact no-go\\
Geometry is quantized & Not derived & Open\\
Occupied common source implies morphology-conditioned terminal assembly &
Yes, Proposition~\ref{prop:moprassembly} & Conditional theorem\\
Source cyclicity and surjective required role maps fully occupy every required
block & Yes,
Proposition~\ref{prop:occupation} & Conditional theorem\\
Rooted positive propagation activates every required block & Yes, but does
not prove full rank & Conditional theorem\\
The occupied typed P3 packet requires an additional variance representation
for the two ordered ports & No; its normalized exterior--record trace gives
Eq.~\eqref{eq:p3variancegram} & Conditional internal theorem\\
Both ordered ports are operationally active & O/Q is forced by normalized
POVM data; M/G is active exactly on nonzero-stress support,
Proposition~\ref{prop:p3portoccupation} & Conditional theorem with exact
zero-stress control\\
Support flag, two ordered ports and independent unit reconstruct the seven
typed degree-one lines & Yes, Theorem~\ref{thm:p3supportportparent}; a smooth
left decoder preserves their rank in the body & Conditional theorem with
missing-port, dependent-unit and collapsed-decoder controls\\
Relation-free Tau subsystems compose by the standard tensor algebra &
Yes under LC1--LC3 by
Theorem~\ref{thm:taustandardcomposition}; LC4 fixes the unrescaled metric &
Conditional exact theorem with hidden-global-relation, crossing-rewrite and
metric-gain controls\\
The explicit class-minimal enriched P3 law owns the complete P3 packet & Yes,
Eq.~\eqref{eq:decisiveparentarrow}, inside the declared source-irredundant
MVP; nonnegative global minimization supplies occupation & Conditional
parent-law realization and occupation theorem\\
Complete pointed exterior-seed occupation makes the P3 root orbit isometric &
Yes; $J_{\rm ext}^\dagger J_{\rm ext}=I_{256}$ in
Theorem~\ref{thm:p3exteriorseedclosure} & Conditional kinematic theorem\\
Full P3 orbit rank and source nonloss alone force the trace--GNS metric & No;
$J_\epsilon$ is a faithful rank-256 nontracial control & Exact countermodel\\
An occupied quadratic ROOT cluster produces a rank-two TT mean & Exactly
under the four conditions in Eq.~\eqref{eq:rootttcertificate}; the nonzero
overlap class has the witness Eq.~\eqref{eq:nonzerokwitness} & Conditional
iff theorem plus existence witness\\
An occupied non-null Abelian field removes arbitrary directional-pair choice
& Yes for $\chi_{\rm dir}>0$; it supplies the loaded weight-one source and
induces weight-two coherence without $r_p<0$ & Conditional selector\\
The source-loaded ROOT--TT mixer is an independent gain & No; it is the
same-functional metric variation in Proposition~\ref{prop:mixownership} &
Conditional ownership theorem\\
Surjective de-duplicated recovery requires a separate Abelian occupation bit
& No; faithfulness gives positive Abelian support and a faithful cyclic GNS
representation unique up to unitary equivalence & Conditional theorem\\
Strict stability plus nonzero activation of every typed role implies faithful
common recovery & No; differential completeness additionally requires
$\operatorname{coker}J_{\rm phys}^{\rm lin}=0$, while finite and algebraic
lanes have their own typed completeness tests & Exact rank theorem and
countercontrol\\
Base--seed/body data exhaust the observer--source differential content & No
on the retained PBSE subsystem; the observer--source complement contributes
one independent row & Exact subsystem rank result\\
The P1, ROOT and gauge labels contribute independent ranks additively & No;
after de-duplication the rank is $4+r_{Q\perp}$ and invertible gauge action
adds zero & Exact overlap-rank theorem\\
Faithful Abelian support forces $\Delta_F>0$ pointwise & No; block support and
local Faraday algebraic type are distinct & Exact scope separation\\
Current base--seed law selects that source representation and operation
algebra & Not derived & Open\\
ROOT one-particle coherence and TT stress arise from one occupied QFT preparation & Yes,
Eq.~\eqref{eq:rootttidentity}, for the declared two-mode packet & Conditional theorem\\
The quantum-induced P1 stress row factors through the complete ROOT descriptor & Yes,
Eq.~\eqref{eq:rootp1factor}, when the descriptor retains $C=N_O\rho_{\rm ROOT}$; body-first staging makes it the full P1 restriction on $\mathcal V_Q$ & Conditional local theorem; not global factorization\\
A single ray activates its own screen-TT terminal & No; its own-screen STF
projection vanishes & Exact control\\
The physical parent selects the required non-collinear $\Delta m=2$ pair &
Not derived & Open\\
Minimal invariant potential forces the occupied coherent pair & Yes under
$u_c,u_p,r_c>0$, $r_p<0$, $\lambda\kappa\neq0$ & Conditional theorem\\
Hamilton, momentum and Gauss constraints follow & Yes, on the selected
standard leaf & Conditional recovery\\
Non-Abelian gauge and Standard Model spectrum follow & Not derived & Open\\
Recovery-complete scalar-15/type-II branch has a source-minimal preimage &
Yes, Eq.~\eqref{eq:scalar15ownership}, only inside the declared P4/type-II
recovery class & Imported class-relative theorem\\
Scalar-15 mass is absolutely predicted from Tau data & No; the ratio
\(M_*/E_P=1/\sqrt{24\pi}\) is derived in the branch and
Eq.~\eqref{eq:scalar15numeric} uses one CODATA anchor & Derived ratio plus
metrological calibration\\
Full interacting RG law is fixed & Yes at the exact functional and Gram
projection level, Eqs.~\eqref{eq:scalar15flow}--
\eqref{eq:scalar15betaprojection}; component evaluation and truncation errors
remain & Imported exact-law handoff\\
Physical universe occupies the scalar-15/type-II branch & Not established by
class-relative reconstruction & Empirically open\\
The earlier narrow base--seed reduct entails $\mathcal L_{P3}^{\min}$ & No,
Eq.~\eqref{eq:narrowparentnogo} and Proposition~\ref{prop:samereduct} & Exact
current-reduct no-go\\
Nature realizes the proposed $\mathcal L_{P3}^{\min}$ & Not established by
internal closure; requires a source-frozen discriminator & Empirically open\\
Incomplete public marginals form one identifying P3 packet & No; the current
pair mosaic has rank $2/3$, lacks certified cross-experiment transports, and
scalar Gram closure would still not test the operator product & Exact
identifiability no-go for the frozen 14-packet census\\
\bottomrule
\end{longtable}

\section{Conclusion}

Paper VII proves a conditional assembly theorem, not an unrestricted physical
unification.  One ordered, once-counted functional can jointly generate the
declared probability, stress, Abelian current and clock-response terminals;
the same packet yields the Einstein--Maxwell--matter equations, Ward/Bianchi
identities and weak-field limits while avoiding duplicate stress ownership.
ROOT one-particle coherence and TT stress are also shown to be restrictions of
one occupied two-mode preparation on the stated sector.

At the representation level, the pointed exterior-seed theorem makes the
256-vector P3 root orbit exactly isometric, while its full-rank nontracial
control proves that rank and sector weights alone are insufficient.  The
operator-square theorem identifies a positive source functional whose zero
set is exactly the Clifford law, and the three-level closure theorem combines
common-source, algebra/representation and terminal closure inside the
declared enriched MVP.  The support--port theorem now reconstructs its seven
typed source directions from a joint projector flag, two ordered cotangent
ports and one independent unit, while a smooth decoder prevents body
collapse.  The explicit law $\mathcal L_{P3}^{\min}$ combines this source
domain with the once-counted operator-square handoff and pointed cyclic root;
Eq.~\eqref{eq:decisiveparentarrow} therefore closes formal law-to-packet
ownership and occupation inside the proposed enriched physical completion:
the explicit zero and nonnegative action make occupation a variational
consequence rather than an extra premise.

Theorem~\ref{thm:taustandardcomposition} also removes the need to import the
standard two-system tensor shape from the terminal side.  Component-local,
support-preserving closure with no hidden global endpoint constraint makes
the relation-free parent closure a Cartesian product, whose finite
linearization is the tensor-product algebra.  The diagonal and crossing-rule
controls show exactly why disjoint labels, complete marginals or a zero mixed
Hessian would have been insufficient.

The exact boundary is now sharper.  Equation~\eqref{eq:narrowparentnogo}
proves that the earlier narrow reduct cannot derive this law.  Whether Nature
realizes $\mathcal L_{P3}^{\min}$ must be decided by a source-frozen,
parameter-free discriminator, not by another internal selector or another
standard-limit recovery.  The scalar-15 handoff closes ownership inside one
recovery class, fixes one dimensionless mass ratio and supplies its exact
functional RG law.  It still does not predict \(G\), establish physical
type-II occupation, quantize geometry, recover the full particle spectrum or
prove unrestricted ambient-parent exhaustivity.
The public-data frontier is now equally finite: the fourteen incomplete pair
packets cannot be fused into an identifying common-source Gram without
source-frozen transport certificates and a measured cross-role edge, and
even scalar Gram closure would not replace operator-product tomography.


\appendix
\section{Detailed P3 typed-role extraction and port geometry}
\label{app:p3technical}

This appendix preserves the complete P3 extraction argument while keeping the
main proof spine visible.  It contains Propositions~19--31: the squared-Hodge
realization, the complete first-jet quotient, source-rank and typed-extraction
tests, represented-role and common-scale candidates, the O/Q and M/G
stabilizer analysis, and the cotangent-port construction.  None of these
technical reductions proves that the unrestricted physical base--seed parent
occupies the required complete typed incidence.

\begin{proposition}[Full-role squared-Hodge selector and P3 realization]
\label{prop:p3clifford}
Let $E_{P3}$ be the seven-role physical quotient with positive metric $K_*$.
If one source-owned linear self-adjoint operator map
$\mathsf C:E_{P3}\to\mathfrak A_{\mathrm{src}}$ satisfies
\begin{equation}
\mathsf C(v)^2=K_*(v,v)I
\qquad\text{for every }v\in E_{P3},
\label{eq:fullrolesquare}
\end{equation}
then
\begin{equation}
\mathsf C(u)\mathsf C(v)+\mathsf C(v)\mathsf C(u)
=2K_*(u,v)I.
\label{eq:p3clifford}
\end{equation}
Thus an orthonormal seven-role frame generates
$\operatorname{Cl}_7(\mathbb C)\simeq
M_8(\mathbb C)\oplus M_8(\mathbb C)$.  In the canonical eight-qubit
Jordan--Wigner embedding into the common
$\operatorname{Cl}_{16}(\mathbb C)\simeq M_{256}(\mathbb C)$ envelope, the
two simple sectors occur with multiplicities $(16,16)$ and normalized-trace
weights $(1/2,1/2)$.
\end{proposition}

\begin{proof}
Apply Eq.~\eqref{eq:fullrolesquare} to $u+v$ and subtract its $u$ and $v$
instances to obtain Eq.~\eqref{eq:p3clifford}.  The complex Clifford
classification gives the two simple rank-eight sectors.  The canonical first
seven Jordan--Wigner generators have 128 linearly independent Clifford words
and central volume eigenvalues of opposite sign with equal eight-dimensional
multiplicity in the minimal faithful direct sum.  Sixteen copies embed this
faithful representation into dimension 256, giving the stated multiplicities
and trace weights.  These finite checks are regenerated by
\path{scripts/audit_theorem_certificates.py}.
\end{proof}

The operator-square premise is stronger than a normalized trace-Gram or
scalar squared-action identity.  Seven commuting diagonal matrices can have
the same identity trace Gram as seven Clifford generators while their generic
linear combination does not square to a scalar.  Hence the P3 packet is an
explicit coefficient-free completion and Eq.~\eqref{eq:fullrolesquare} is an
exact selector, but the unrestricted base--seed reduct has not selected the
required typed seven-role operator incidence as part of one occupied packet.
The existing rank-two EACER and rank-nine
seed/context squared-Hodge carriers have the wrong domains.  A post-body
Clifford handoff is therefore an admissible minimal working completion, not an
unrestricted base--seed derivation or a modification of the already stabilized body.

There is a further type distinction.  The seven entries in the P3 support
ledger are unital role algebras, not seven canonically one-dimensional source
slots.  Their units are identified in the unital primitive colimit and hence
span only one common line.  Selecting one non-scalar element inside each role
algebra would add noncanonical data and need not be invariant under internal
role automorphisms.  Therefore the support ledger and the current colimit do
not by themselves construct the seven-dimensional domain used in
Proposition~\ref{prop:p3clifford}.

The sufficient source object can be stated without choosing such generators.
Let $(\mathcal P_{P3}^1,G)$ be the complete de-duplicated real physical
first-jet module after the true-null quotient, and suppose the source owns a
surjection
\begin{equation}
q_{P3}:\mathcal P_{P3}^1\longrightarrow E_{P3}^{\rm lbl}\simeq\mathbb R^7
\label{eq:p3rolequotient}
\end{equation}
whose seven typed rows are nonzero, whose kernel is annihilated by every
claimed P3 first-jet terminal, and through which every relevant degree-one
P3 response factors.  If the common action unit normalizes it as
\begin{equation}
q_{P3}G^{-1}q_{P3}^{*}=I_7,
\label{eq:p3coisometry}
\end{equation}
then its metric adjoint
$j_{P3}=G^{-1}q_{P3}^{*}$ is a canonical isometric section:
\begin{equation}
q_{P3}j_{P3}=I_7,
\qquad
j_{P3}^{*}G j_{P3}=I_7.
\label{eq:p3isometricsection}
\end{equation}
The exterior construction on $E_{P3}^{\rm lbl}$ then supplies the normalized
Clifford law of Proposition~\ref{prop:p3clifford}.  This is a conditional
role-quotient theorem, not a derivation of Eq.~\eqref{eq:p3rolequotient} from
the physical base--seed system.  Existing seed/body decoders, primitive-edge
label algebras and downstream separating-record witnesses have different
typed domains and cannot be reused as $q_{P3}$ without an additional source
theorem.

The normalization in Eq.~\eqref{eq:p3coisometry} is not an independent
physical parent-law coefficient if a complete typed role differential is already
source-owned.  For a surjection
$R_{P3}:\mathcal P_{P3}^1\to E_{P3}^{\rm lbl}$, define
\begin{equation}
H_R=R_{P3}G^{-1}R_{P3}^{*}>0,
\qquad
q_{P3}=H_R^{-1/2}R_{P3}.
\label{eq:p3polarquotient}
\end{equation}
The unique positive square root gives
$q_{P3}G^{-1}q_{P3}^{*}=I_7$ and
$\ker q_{P3}=\ker R_{P3}$.

The kernel and no-bypass clauses can be reduced further.  Let
$\{L_\alpha\}_{\alpha\in\mathfrak F_{P3}^{1}}$ be the complete declared
family of degree-one P3 operator and terminal jets on
$\mathcal P_{P3}^{1}$, and set
\begin{equation}
L_\oplus=\bigoplus_\alpha L_\alpha,
\qquad
N_{P3}^{1}=\ker L_\oplus,
\qquad
E_{P3}^{\rm phys}=\mathcal P_{P3}^{1}/N_{P3}^{1}.
\label{eq:p3nullquotient}
\end{equation}

\begin{proposition}[Canonical complete-first-jet quotient]
\label{prop:p3nullquotient}
With the minimum-lift metric induced by $G$, the canonical quotient map
$q_{\rm null}:\mathcal P_{P3}^{1}\to E_{P3}^{\rm phys}$ is surjective,
has kernel $N_{P3}^{1}$, and every $L_\alpha$ factors uniquely through it.
Its horizontal metric adjoint is the minimum-norm isometric section.  Any
other coisometric quotient with the same kernel differs only by an isometry
of $E_{P3}^{\rm phys}$, restricted to a typing-preserving isometry when the
role labels are retained.  Moreover,
\begin{equation}
\dim E_{P3}^{\rm phys}=\operatorname{rank}L_\oplus.
\label{eq:p3nullquotientrank}
\end{equation}
Consequently a complete rank-seven first-jet family supplies the carrier used
in Proposition~\ref{prop:p3clifford} without an additional role-decoder,
normalization or degree-one-bypass choice.
\end{proposition}

\begin{proof}
The first isomorphism theorem gives the kernel, surjectivity and
Eq.~\eqref{eq:p3nullquotientrank}.  Since
$N_{P3}^{1}\subseteq\ker L_\alpha$ for every $\alpha$, each map is constant
on quotient fibres and therefore factors uniquely.  Positivity of $G$ gives
the orthogonal splitting
$\mathcal P_{P3}^{1}=N_{P3}^{1}\oplus(N_{P3}^{1})^{\perp_G}$; restriction of
the quotient map to the horizontal summand defines the minimum-lift metric and
the isometric section.  Two coisometries with the same kernel restrict to
isometries from this same horizontal summand, so their relative map is an
isometry.  QED.
\end{proof}

If a proposed additional degree-one row does not annihilate $N_{P3}^{1}$,
then it does not constitute an independent bypass parameter of the same
complete quotient: it strictly refines the joint null and proves that the
declared family was incomplete.  Thus the remaining upstream premise is only
that the physical base--seed law supplies the complete family in
Eq.~\eqref{eq:p3nullquotient} with rank seven.  The current seed-presentation
candidate has the abstract ingredients from which such a family could be
formed, but the present seed grammar proves neither its completeness nor its
rank in the physically realized packet.

This upstream qualification does not reopen the three downstream levels
already assembled in the enriched MVP: common-source terminal factorization,
the occupied P3 algebra--record--GNS packet, and the conditional terminal
recoveries remain intact under their declared assumptions.  The unresolved
statement is only that the unrestricted physical base--seed parent must select
that joint packet.

The rank premise has an exact finite source certificate.  Let
$K_{\rm src}>0$ be the frozen source Hessian on the true-null-quotiented
first-jet module and let $L_\oplus$ have seven typed rows.  Define
\begin{equation}
G_{P3}^{(1)}=L_\oplus K_{\rm src}^{-1}L_\oplus^*.
\label{eq:p3rankgram}
\end{equation}

\begin{proposition}[Rank-seven source certificate and current-parent no-go]
\label{prop:p3rankcertificate}
The following statements are equivalent:
\begin{equation}
\operatorname{rank}L_\oplus=7
\quad\Longleftrightarrow\quad
G_{P3}^{(1)}>0
\quad\Longleftrightarrow\quad
\lambda_{\min}(G_{P3}^{(1)})>0
\quad\Longleftrightarrow\quad
\det G_{P3}^{(1)}>0.
\label{eq:p3rankequivalence}
\end{equation}
Strict source stability, nonzero activation of all seven typed rows, equal
primitive row costs and a common action unit do not imply these conditions.
Consequently the conditional three-level closure cannot be inverted to prove
rank seven.
\end{proposition}

\begin{proof}
Equation~\eqref{eq:p3rankgram} is the Gram matrix of
$L_\oplus K_{\rm src}^{-1/2}$, and $K_{\rm src}^{-1/2}$ is invertible.
Therefore its rank equals $\operatorname{rank}L_\oplus$; the remaining
equivalences are the finite-dimensional full-rank criteria for a positive
semidefinite $7\times7$ matrix.  For non-entailment, take the same source
space, $K_{\rm src}=I_7$, load and stationary point.  Let $L_7=I_7$, whereas
$L_6$ has rows $e_1,\ldots,e_6$ and
$(e_1+\cdots+e_6)/\sqrt6$.  Both maps have seven nonzero unit rows and total
row-Gram trace seven, but $\operatorname{rank}L_7=7$ and
$\operatorname{rank}L_6=6$.  Hence all the listed upstream properties hold
in both controls while Eq.~\eqref{eq:p3rankequivalence} holds only in the
first.  QED.
\end{proof}

The PS1 triangular dependency action already supplies the untyped rank on its
own codomain.  Let $\pi_{\rm dep}$ delete the root residual and define
\begin{equation}
D_{\rm dep}=\pi_{\rm dep}D_\Sigma:
\bigoplus_{i=0}^{7}E_i\longrightarrow
Z_{\rm dep}:=\bigoplus_{i=1}^{7}E_i.
\label{eq:p3dependencyoperator}
\end{equation}

\begin{proposition}[Onto dependency response and typed-extraction reduction]
\label{prop:p3typedextraction}
The map $D_{\rm dep}$ is surjective and
\begin{equation}
D_{\rm dep}H_{\rm src}^{-1}D_{\rm dep}^*
=\mathcal A_*^{-1}\pi_{\rm dep}K_\Sigma^{-1}\pi_{\rm dep}^*>0.
\label{eq:p3dependencygram}
\end{equation}
If a source-owned typed extraction
$T_{P3}:Z_{\rm dep}\to E_{P3}^{\rm lbl}$ realizes the complete P3 first-jet
family by $L_\oplus=T_{P3}D_{\rm dep}$, then
\begin{equation}
\operatorname{rank}L_\oplus=\operatorname{rank}T_{P3}.
\label{eq:p3typedrank}
\end{equation}
Thus rank seven follows exactly when this complete typed extraction is onto.
The present PS1 action does not select that extraction.
\end{proposition}

\begin{proof}
The restriction of $D_{\rm dep}$ to
$\bigoplus_{i=1}^{7}E_i$ is block lower triangular with identity diagonal,
so it is invertible.  Inverting
$H_{\rm src}=\mathcal A_*D_\Sigma^*K_\Sigma D_\Sigma$ and substituting
$D_{\rm dep}=\pi_{\rm dep}D_\Sigma$ gives
Eq.~\eqref{eq:p3dependencygram}; it is positive because it is a principal
compression of $K_\Sigma^{-1}>0$.  Surjectivity gives
$\operatorname{im}(T_{P3}D_{\rm dep})=\operatorname{im}T_{P3}$, proving
Eq.~\eqref{eq:p3typedrank}.  The rank-six and rank-seven extraction controls
in Proposition~\ref{prop:p3rankcertificate} share the same dependency action,
so that action cannot select between them.  QED.
\end{proof}

The seven dependency blocks $S1,\ldots,S7$ are not the seven generally
non-scalar P3 ownership roles $\{B,O,R,U,Q,M,G\}$.  Consequently this result
does not authorize a label-based identification.  It shows instead that raw
seven-direction source response is already available: the remaining
unrestricted physical datum is one sourced, complete and onto P3-typed
extraction.  No new carrier, decoder gain or terminal normalization is
required, and downstream terminal maps may not be used to choose it.

The common represented trace does not by itself construct that extraction.
On a simple role block $\mathcal A_a=M_n(\mathbb C)$, every real linear
functional on $\mathcal A_a^{\rm sa}$ invariant under all inner
automorphisms is a multiple of the trace and therefore vanishes on the
centered tangent
$\mathcal A_{a,0}^{\rm sa}=\{x=x^*: \omega_a(x)=0\}$.  Thus neither the
trace nor the role ledger canonically chooses a noncentral role direction.

The occupied seed representation nevertheless supplies a canonical
candidate when its role values are noncentral.  Let
$E_a:X_{\rm src}\to\mathcal A_a^{\rm sa}$ be the source-owned represented
role evaluation and define

\begin{equation}
h_a=E_a(x_*)-\omega_a(E_a(x_*))I,
\qquad
\ell_a(\delta x)=
{\operatorname{Re}\omega_a(h_aD E_a\delta x)
 \over \sqrt{\omega_a(h_a^2)}}.
\label{eq:p3radialcovector}
\end{equation}

This is defined for $h_a\ne0$ and is natural under simultaneous unitary
conjugation of the occupied value and its derivative.  Let
$L_{\rm rad}=\bigoplus_a\ell_a:X_{\rm src}\to\mathbb R^7$.

\begin{proposition}[Occupied-role radial extraction criterion]
\label{prop:p3radialextraction}
There is a unique typed map
$T_{P3}^{\rm rad}:Z_{\rm dep}\to E_{P3}^{\rm lbl}$ satisfying
$L_{\rm rad}=T_{P3}^{\rm rad}D_{\rm dep}$ if and only if
\begin{equation}
\ker D_{\rm dep}\subseteq\ker L_{\rm rad}.
\label{eq:p3radialdescent}
\end{equation}
When it exists, it is onto if and only if
\begin{equation}
\operatorname{rank}L_{\rm rad}=7
\quad\Longleftrightarrow\quad
L_{\rm rad}H_{\rm src}^{-1}L_{\rm rad}^{*}>0.
\label{eq:p3radialgram}
\end{equation}
The present source ledger does not entail either condition.
\end{proposition}

\begin{proof}
Surjectivity of $D_{\rm dep}$ identifies $Z_{\rm dep}$ with
$X_{\rm src}/\ker D_{\rm dep}$.  The universal property of this quotient
gives the unique factor precisely under
Eq.~\eqref{eq:p3radialdescent}.  Its image equals the image of
$L_{\rm rad}$, and the final equivalence is the positive Gram criterion of
Proposition~\ref{prop:p3rankcertificate}.  Non-entailment is sharp: a zero
$h_a$ removes a row; two nonzero equal-cost anchors may induce the same
covector and give rank six; and a covector nonzero on $\ker D_{\rm dep}$
does not descend.  QED.
\end{proof}

This replaces an arbitrary $T_{P3}$ by a derivative of occupied role data
whenever the physical seed representation supplies the seven anchors.  It
does not prove that the physical parent law supplies them, that the deleted root is invisible
to all seven radial covectors, or that their Gram is positive definite.  A
non-radial typed differential remains possible only as separately sourced
physics.

The common-unit role requires a typed correction to the all-radial candidate.
If $E_U(x_*)=cI$, then its centered anchor is zero and
Eq.~\eqref{eq:p3radialcovector} supplies no $U$ row.  If the source instead
owns a differentiable positive scale/clock-unit character
$u_*:X_{\rm src}\to\mathbb R_+$, define
\begin{equation}
\kappa_U=d\log u_*|_{x_*},
\qquad
\ell_U^{\rm sc}={\kappa_U\over
\sqrt{\kappa_UH_{\rm src}^{-1}\kappa_U^*}}.
\label{eq:p3unitscalerow}
\end{equation}
Replacing the zero centered $U$ row by $\ell_U^{\rm sc}$ gives a hybrid
family to which Proposition~\ref{prop:p3radialextraction} applies unchanged.
The common action unit normalizes a supplied nonzero unit incidence; it does
not make $\kappa_U$ nonzero.  Hence the current source packet must still
derive that scale/clock differential or another independently typed $U$
incidence.

\begin{proposition}[Current-source metric non-identifiability]
\label{prop:p3currentevaluability}
The current PS1 source ledger does not determine the hybrid seven-role Gram,
even after restricting to full-rank extractions with equal row cost.
\end{proposition}

\begin{proof}
The ledger gives the conditional common tangent $J_*$ and the aggregate
restrictions $e_{P1}J_*$ and $e_OJ_*$, but not the seven represented
evaluation jets $E_a(x_*),DE_a|_{x_*}$ or a selected $d\log u_*$.  This
already prevents evaluation of Eqs.~\eqref{eq:p3radialcovector} and
\eqref{eq:p3unitscalerow}.  Non-identifiability is sharp.  Whiten the fixed
onto dependency response to
$Q_{\rm dep}H_{\rm src}^{-1}Q_{\rm dep}^*=I_7$.  Set $T_0=I_7$ and, for
$0<|c|<1$, let $T_c$ leave all rows fixed except
$e_2^*T_c=ce_1^*+\sqrt{1-c^2}\,e_2^*$.  Both
$L_0=T_0Q_{\rm dep}$ and $L_c=T_cQ_{\rm dep}$ factor through the same
dependency quotient, have rank seven, and have unit diagonal Gram.  However,
the second Gram has off-diagonal entry $c$ and determinant $1-c^2$, whereas
the first is $I_7$.  Therefore the common reduct, rank and row normalization
do not select the physical P3 metric.  QED.
\end{proof}

The finite missing datum is one occupied represented role-evaluation jet
$\{E_a(x_*),DE_a|_{x_*}\}_{a\ne U}$ together with the source-owned unit row
$d\log u_*|_{x_*}$.  Supplying that packet from the physical base--seed
representation would make the descent and positive-Gram tests decisive in
one calculation.  Choosing it from terminal behavior would instead be a
fitted completion.

The common-source terminal assembly cannot be inverted to remove this
missing datum.  If $J_{\rm term}=\mathcal T A_O$ is the stacked terminal
precursor, a factorization $L_{\rm mix}=R J_{\rm term}$ exists exactly when
$\ker J_{\rm term}\subseteq\ker L_{\rm mix}$.  The forward MOPR theorem gives
only $\operatorname{rank}J_{\rm term}\leq\operatorname{rank}A_O$ and does not
establish this kernel inclusion.  Terminal recovery of the source jet would
therefore require an additional source-owned injectivity theorem and cannot
be inferred from successful terminal limits.

The closest existing source module does not remove the gap by dimension
matching.  The conditional HRC role one-jet is
$V_{\rm HRC}=E_S\oplus E_S$, two copies of the three-dimensional standard
$S_4$ representation.  The six non-unit P3 entries, however, are fixed
semantic ownership labels under the current type-preserving ledger; no
tetrahedral action mixing them has been derived.  Therefore
\begin{equation}
\operatorname{Hom}_{S_4}
\left(E_S\oplus E_S,\mathbf1^{\oplus6}\right)=0.
\label{eq:p3hrctypenogo}
\end{equation}
Indeed, each row of an equivariant map would be an invariant covector on a
standard module, and no such nonzero covector exists.  A matched
$E_S\oplus E_S$ target action admits the identity intertwiner, so the
obstruction is typing rather than dimension.  Deriving such a target action
would itself be the missing covariant role bridge and would have to preserve
the distinct typed relation-support classes.  Consequently the count
$6+1=7$, even after adding the independent $U$ scale row, does not construct
the occupied P3 evaluation jet.

The six non-unit maps can nevertheless be reduced to one existing source
object.  Let $\mathcal A_{\rm curr}$ be the de-duplicated primitive colimit,
$j_a:\mathcal A_a\to\mathcal A_{\rm curr}$ its canonical role maps, and let
$x\mapsto\Pi_x$ be a $C^1$ germ of the single coherent physical unital
$*$-homomorphism.  For source-declared self-adjoint anchors
$z_a\in\mathcal A_a$, define
\begin{equation}
E_a(x)=\Pi_x(j_a(z_a)),
\qquad
DE_a|_{x_*}=D\Pi|_{x_*}(j_a(z_a)).
\label{eq:p3coconejet}
\end{equation}
All six non-unit evaluation jets are then restrictions of one representation
one-jet.  Differentiating the homomorphism and dagger identities gives
\begin{equation}
D\Pi(ab)=D\Pi(a)\Pi(b)+\Pi(a)D\Pi(b),
\qquad D\Pi(a^*)=D\Pi(a)^*,
\qquad D\Pi(1)=0.
\label{eq:p3coconelinearized}
\end{equation}
Thus the missing maps are not six independent completion choices once the
physical typed presentation, anchors and cocone germ are supplied.  The last
identity also proves that the algebraic unit cannot generate the nonzero
absolute $U$ scale row.  That row must come from a distinguished positive
$U$-role scale generator or from the separately sourced character $u_*$.  The
remaining finite source datum is therefore one occupied cocone one-jet, its
typed anchor tuple and the unit-scale differential.  The current parent has
not selected this packet, and cocone coherence by itself does not imply the
rank-seven positive-Gram condition.

The unit-scale differential is not wholly structureless on the conditional
protected relation-unit branch already isolated upstream.  If the physical
source occupies the commuting relation-unit square in a source-frozen chart,
then the positive character may be chosen as
$u_*=U_R=a_*/c_R$, and therefore
\begin{equation}
 \boxed{\kappa_U=d\log U_R=d\log a_*-d\log c_R.}
 \label{eq:p3relationscalepullback}
\end{equation}
In the common CHDF Hodge chart, where $U_R=\sqrt2\,\ell_H$, this becomes
$\kappa_U=d\log\ell_H$.  This is a conditional pullback identity, not an
activation theorem.  Equal logarithmic source responses give
$\kappa_U=0$; moreover, the finite common-unit witness fixes ratios at one
packet but does not itself provide a nonzero source tangent.  Thus the
remaining question is physical occupation of the protected square together
with nonzero activation, rather than the form of an otherwise arbitrary
seventh row.

There is also an exact finite novelty certificate.  Let $L_6$ collect the six
descended non-unit rows and assume
$G_6=L_6H_{\rm src}^{-1}L_6^*>0$.  With
$b_U=L_6H_{\rm src}^{-1}\kappa_U^*$, define
\begin{equation}
 \sigma_U=\kappa_UH_{\rm src}^{-1}\kappa_U^*
 -b_U^*G_6^{-1}b_U.
 \label{eq:p3uscaleSchur}
\end{equation}
The block-Gram Schur-complement theorem gives
\begin{equation}
 \operatorname{rank}\!\begin{pmatrix}L_6\\ \kappa_U\end{pmatrix}=7
 \quad\Longleftrightarrow\quad \sigma_U>0.
 \label{eq:p3uscalerank}
\end{equation}
Thus even a nonzero unit differential is insufficient when it lies in the
source-metric row span of the six previous roles.  Physical activation needs
protected-square occupation, quotient descent of $\kappa_U$, and the positive
certificate~\eqref{eq:p3uscaleSchur}.

The unchanged source inventory cannot supply that activation.  The derived
source arrows form a finite $S_4$ relabeling groupoid with zero continuous
tangent dimension, while the occupied CHDF seed orbit is locally isolated
with zero on-shell seed-moduli tangent.  Therefore
\begin{equation}
 T_{x_*}X_{\rm phys}^{\rm curr}=0
 \quad\Longrightarrow\quad
 d\log U_R|_{T_{x_*}X_{\rm phys}^{\rm curr}}=0.
 \label{eq:p3currentsourcenogo}
\end{equation}
The available off-shell seed--body response germ is not a selected family of
physical source arrows and cannot replace this missing tangent.

The adopted MCSCL forward completion provides a conditional reopening.  Its
common calibration tangent $c_U$ satisfies
\begin{equation}
 D\log\ell_H(c_U)=1,\qquad
 D\log\mathcal A_*(c_U)=2,
 \qquad\text{hence}\qquad \kappa_U(c_U)=1
 \label{eq:p3mcsclactivation}
\end{equation}
on the common CHDF/PRUB chart.  Calibration-basicness of the six raw rows is
not an additional physical premise.  For an arbitrary descended six-row
family define its measured common-scale loading and relative part by
\begin{equation}
 v_U:=L_6c_U,\qquad L_6^{\rm rel}:=L_6-v_U\kappa_U.
 \label{eq:p3mcsclseparation}
\end{equation}
Then $L_6^{\rm rel}c_U=0$ and
\begin{equation}
 \begin{pmatrix}L_6^{\rm rel}\\\kappa_U\end{pmatrix}
 =
 \begin{pmatrix}I_6&-v_U\\0&1\end{pmatrix}
 \begin{pmatrix}L_6\\\kappa_U\end{pmatrix}.
 \label{eq:p3scalehortransform}
\end{equation}
The block-triangular row transformation is source-determined, invertible and
has determinant one.  It preserves the complete joint kernel and rank, while
the source-metric Gram changes only by congruence.  Since the relative rows
annihilate $c_U$ but $\kappa_U(c_U)=1$,
\begin{equation}
 \operatorname{rank}\!\begin{pmatrix}L_6\\\kappa_U\end{pmatrix}=7
 \quad\Longleftrightarrow\quad
 \operatorname{rank}L_6^{\rm rel}=6
 \quad\Longleftrightarrow\quad
 L_6^{\rm rel}H_{\rm src}^{-1}(L_6^{\rm rel})^*>0.
 \label{eq:p3scalehorrank}
\end{equation}
Thus the common scale is counted exactly once without discarding it.  This
removes raw calibration-basicness as a completion assumption, but does not
manufacture a missing role direction.  The unrestricted physical parent must
still own the MCSCL character, the protected relation-unit square, the six
occupied evaluation jets, and their positive rank-six relative Gram.

The onto dependency response already supplies the underlying relative
carrier.  Let
$D_{\rm dep}:X_{\rm src}\to Z_{\rm dep}$ be the surjective seven-dimensional
map above and assume quotient descent
\begin{equation}
 \kappa_U=\bar\kappa_UD_{\rm dep},\qquad
 \bar\kappa_U(D_{\rm dep}c_U)=1.
 \label{eq:p3relsplit}
\end{equation}
With $z_U=D_{\rm dep}c_U$, define
\begin{equation}
 E_{\rm dep}^{\rm rel}:=\ker\bar\kappa_U,\qquad
 P_{\rm rel}:=I-z_U\bar\kappa_U.
 \label{eq:p3relprojector}
\end{equation}
Direct multiplication gives
\begin{equation}
 P_{\rm rel}^2=P_{\rm rel},\quad
 \operatorname{im}P_{\rm rel}=E_{\rm dep}^{\rm rel},\quad
 \ker P_{\rm rel}=\operatorname{span}\{z_U\},\quad
 \dim E_{\rm dep}^{\rm rel}=6.
 \label{eq:p3relcarrier}
\end{equation}
Restriction of the positive dependency metric to this kernel is positive
definite.  Hence no additional six-dimensional carrier or metric is needed.
What remains unselected is the occupied typed isomorphism
$T_{\rm rel}:E_{\rm dep}^{\rm rel}\to
E_{P3}^{\ne U}=\operatorname{span}\{B,O,R,Q,M,G\}$.
Orthogonal frames of the same relative carrier have identical upstream
status but different semantic assignments, so positivity alone cannot choose
$T_{\rm rel}$.

The existing typed support ledger reduces, but does not remove, this
ambiguity.  On the non-unit role set retain the named supports
\begin{equation}
 H_{P4}=\{B,M,G\},\qquad
 H_{OI}=\{O,Q\},\qquad
 H_{PM}=\{R,M,G\}.
 \label{eq:p3incidencesupports}
\end{equation}
Here $H_{P4}$ is the non-unit, non-$I$ restriction of the maximal P4
support, while $H_{OI}$ and $H_{PM}$ are the observer--instrument and
photon--matter relation supports.  The incidence fingerprints with respect
to this ordered support family are
\begin{equation}
\begin{array}{c|cccccc}
a&B&O&R&Q&M&G\\ \hline
\chi(a)&100&010&001&010&101&101.
\end{array}
\label{eq:p3incidencefingerprints}
\end{equation}
Therefore the permutation stabilizer of every named support is exactly
\begin{equation}
 \operatorname{Aut}(\mathcal H_{P3}^{\rm rel})
 \simeq S_{\{O,Q\}}\times S_{\{M,G\}},
 \qquad |\operatorname{Aut}(\mathcal H_{P3}^{\rm rel})|=4.
 \label{eq:p3incidenceautomorphism}
\end{equation}
Indeed, $B$ and $R$ have unique fingerprints; the only repeated columns are
the $O/Q$ and $M/G$ pairs, and either pair swap preserves all three supports.
Consequently, if these supports descend to six occupied atomic rays of
$E_{\rm dep}^{\rm rel}$, support preservation reduces the semantic label
freedom to two binary choices.  A source-owned distinction on each unresolved
pair is then necessary and sufficient to make the permutation stabilizer
trivial; one common source object may supply both distinctions.  This is not
yet a construction of $T_{\rm rel}$: the current PS1
action proves neither the atomic support descent nor those two oriented
separators.  It is a finite sharpening of the missing datum, not permission
to select the pair orientations from terminal behavior.

This remainder has an exact quotient formulation.  Put

\begin{equation}
 q_{\rm rel}:=P_{\rm rel}D_{\rm dep}:X_{\rm src}
 \longrightarrow E_{\rm dep}^{\rm rel}.
 \label{eq:p3relativequotient}
\end{equation}
Let $\Xi_{\rm src}:X_{\rm src}\to V_{\rm sep}$ be one source-owned
descriptor.  Since $q_{\rm rel}$ is onto, there is a unique descended map
$\bar\Xi:E_{\rm dep}^{\rm rel}\to V_{\rm sep}$ satisfying
$\Xi_{\rm src}=\bar\Xi q_{\rm rel}$ exactly when
\begin{equation}
 \ker q_{\rm rel}\subseteq\ker\Xi_{\rm src}.
 \label{eq:p3selectordescent}
\end{equation}

\begin{proposition}[Relative pair-selector criterion]
\label{prop:p3pairselector}
Assume that the supports in Eq.~\eqref{eq:p3incidencesupports} descend to
occupied atomic rays with source-owned normalized representatives
$r_a\in E_{\rm dep}^{\rm rel}$.  The subgroup that
preserves both those supports and the values
$\xi_a:=\bar\Xi(r_a)$ is trivial if and only if
\begin{equation}
 \boxed{\xi_O\ne\xi_Q\quad\hbox{and}\quad\xi_M\ne\xi_G.}
 \label{eq:p3pairseparation}
\end{equation}
Thus a single common source descriptor can remove both residual binary
ambiguities without a new carrier or role-dependent gain.
\end{proposition}

\begin{proof}
Equation~\eqref{eq:p3incidenceautomorphism} shows that the support stabilizer
is generated by the independent swaps $(O\ Q)$ and $(M\ G)$.  The first
preserves the descriptor exactly when $\xi_O=\xi_Q$, and the second exactly
when $\xi_M=\xi_G$.  Their common stabilizer is therefore trivial precisely
under Eq.~\eqref{eq:p3pairseparation}.  The descent statement is the universal
property of the onto quotient $q_{\rm rel}$.
\end{proof}

The irreducible observer--source incidence and post-body variational ownership
have the right shape for a two-component candidate: one could distinguish
$O/Q$, the other $M/G$.  Existing rank and terminal-variation theorems do not,
however, prove that they form one source-owned descriptor, satisfy
Eq.~\eqref{eq:p3selectordescent}, or obey
Eq.~\eqref{eq:p3pairseparation}.  Terminal meanings cannot supply those
parent-side facts circularly.  The finite certificate checks the order-four,
one-pair order-two, both-pair trivial and nonfactorizing-null controls; it does
not assert physical base--seed occupation of $\Xi_{\rm src}$.

There is a further exact restriction on the natural candidate. Let
$\Delta_{\rm rec}$ be the present commutative special dagger--Frobenius copy,
$S$ the exchange of its two tensor legs and $P_-=(I-S)/2$. On its copy basis,
\begin{equation}
 \Delta_{\rm rec}(e_k)=e_k\otimes e_k,
 \qquad S\Delta_{\rm rec}=\Delta_{\rm rec},
 \qquad P_-\Delta_{\rm rec}=0.
 \label{eq:p3frobeniusorientation}
\end{equation}

\begin{proposition}[Current-candidate no-go and ordered-port target]
\label{prop:p3currentcandidatenogo}
Relative to the source objects presently constructed, irreducible
observer--source incidence, post-body matter/geometry variational typing and
the commutative Frobenius copy do not define a noncircular
$\Xi_{\rm src}$ satisfying Proposition~\ref{prop:p3pairselector}. A
sufficient finite replacement is one source-owned ordered typed port
incidence
\begin{equation}
 \Pi_{\rm port}:X_{\rm src}\longrightarrow V_{OQ}\oplus V_{MG},
 \qquad
 \Pi_{\rm port}=\bar\Pi_{\rm port}q_{\rm rel},
 \label{eq:p3portincidence}
\end{equation}
whose descended values distinguish both $O/Q$ and $M/G$.
\end{proposition}

\begin{proof}
The observer--source construction supplies a transverse image and rank, but
turning it into an $O/Q$ value requires an additional source-owned codomain
functional or port assignment. Matter/geometry variational ownership is
typed only after the semantic P3 roles have been identified; pulling it back
to $E_{\rm dep}^{\rm rel}$ through the unknown $T_{\rm rel}$ would use the
desired identification in its own construction. Equation
\eqref{eq:p3frobeniusorientation} shows that the current copy has no
exchange-odd component. The current objects therefore coexist with both
residual pair swaps and do not imply Eq.~\eqref{eq:p3pairseparation}.
Conversely, Eq.~\eqref{eq:p3portincidence} obeys the quotient descent
condition by construction, and Proposition~\ref{prop:p3pairselector} removes
the complete residual stabilizer when its two pairwise inequalities hold.
\end{proof}

The previously derived exchange-odd charged-wall operator
$D_{12}=Y_A\otimes I-I\otimes Y_A$ is a type witness for such orientation,
not this selector. It acts on a conditional charged two-leg realization, the
current parent inventory does not own its required $E_{60}$ carrier, and no
factorization through $q_{\rm rel}$ has been established. Proposition
\ref{prop:p3currentcandidatenogo} is therefore a no-go for the current source
inventory, not a theorem that an ordered port incidence cannot exist. Its
physical base--seed occupation remains open.

The existing universal-seed orientation is only one binary sign. This gives
an additional finite no-go.

\begin{proposition}[Single-seed-sign no-go]
\label{prop:p3singleseedbitnogo}
Let
$G_{\rm pair}=S_{\{O,Q\}}\times S_{\{M,G\}}
\simeq\mathbb Z_2\times\mathbb Z_2$ be the residual support stabilizer.
No equivariant action of $G_{\rm pair}$ on one binary seed-orientation line
is faithful. Thus the existing scalar seed sign alone cannot orient both
$O/Q$ and $M/G$.
\end{proposition}

\begin{proof}
Such an action is a character
$\chi:G_{\rm pair}\to\{+1,-1\}\simeq\mathbb Z_2$. Its image has order at
most two, whereas $|G_{\rm pair}|=4$. The first isomorphism theorem therefore
gives $|\ker\chi|\geq2$. At least one nontrivial pair swap remains invisible.
\end{proof}

This result does not multiply the universal morphology seed. One structured
seed or joint base--seed object may carry the faithful bidegree
$(\chi_{OQ},\chi_{MG})$ with trivial common kernel. What is excluded is only
identifying the single global causal/time orientation bit with the complete
relative-role selector. The two typed characters, or an equivalent ordered
port incidence, still require a source derivation and quotient descent.

The required two-character object nevertheless has a canonical finite
representation once its source ownership is supplied. Let $K_{\rm rel}>0$
be the relative metric and let $J_{OQ},J_{MG}$ be commuting
$K_{\rm rel}$-isometric involutions. Assume each has a source-oriented
one-dimensional negative eigenspace, spanned by unit vectors $c_{OQ}$ and
$c_{MG}$, with each anchor even under the other involution. Define
\begin{equation}
 \bar\Pi_{\rm port}(v)=
 \bigl(K_{\rm rel}(c_{OQ},v),K_{\rm rel}(c_{MG},v)\bigr).
 \label{eq:p3portmetricdual}
\end{equation}

\begin{proposition}[Ordered-port representation]
\label{prop:p3portrepresentation}
Suppose $J_{OQ}$ exchanges $r_O,r_Q$ and fixes $r_M,r_G$, whereas
$J_{MG}$ exchanges $r_M,r_G$ and fixes $r_O,r_Q$. If
\begin{equation}
 K_{\rm rel}(c_{OQ},r_O)\ne0,
 \qquad K_{\rm rel}(c_{MG},r_M)\ne0,
 \label{eq:p3portnonzero}
\end{equation}
then Eq.~\eqref{eq:p3portmetricdual} separates both residual pairs and has
trivial support stabilizer. If the raw involutions and anchors preserve
$\ker q_{\rm rel}$, its pullback through $q_{\rm rel}$ is the source
descriptor required by Proposition~\ref{prop:p3pairselector}.
\end{proposition}

\begin{proof}
Isometry and oddness give
$K_{\rm rel}(c_{OQ},r_Q)=
K_{\rm rel}(c_{OQ},J_{OQ}r_O)=
-K_{\rm rel}(c_{OQ},r_O)$, and the first condition in
Eq.~\eqref{eq:p3portnonzero} makes these values distinct. The same argument
applies to $M/G$. Proposition~\ref{prop:p3pairselector} then gives the
trivial stabilizer, while preservation of $\ker q_{\rm rel}$ is precisely
the quotient-descent condition.
\end{proof}

The conditional oriented observer--source double cover already present in
the hub is a candidate for the first odd line, and the once-counted
matter/geometry variational split motivates the second. The current physical
source ledger proves neither joint occupation of those two oriented lines nor
their preservation of $\ker q_{\rm rel}$. Proposition
\ref{prop:p3portrepresentation} is therefore a finite sufficient realization
theorem, not unrestricted parent-law realization.

The two unresolved pairs nevertheless share one variational structure in the
standard leaf. A normalized observer effect evaluates a quantum state,
whereas a metric variation evaluates matter stress:
\begin{equation}
 \mathfrak e_{OQ}(E,\rho)=\operatorname{Tr}(\rho E),
 \qquad
 \mathfrak e_{MG}(h,T)
 =-\frac12\int d^4x\sqrt{-g}\,T^{\mu\nu}h_{\mu\nu}.
 \label{eq:p3dualevaluations}
\end{equation}
These terminal identities do not derive a parent carrier. They do require
every differentiable completion that recovers them to retain two typed
source/response dualities. Let $i\in\{OQ,MG\}$, let
$W_i\subset E_{\rm dep}^{\rm rel}$ be the protected pair subspace, let $h_i>0$
be its inherited metric, let
$\bar\eta_i\in(E_{\rm dep}^{\rm rel})^*$ be a nonzero covector supported on
$W_i$, and admit one conjugate response $v_i\in W_i$. Assume a source-owned
variance representation maps the source and response legs to the two occupied
role rays of the pair. Consider
\begin{equation}
 \Gamma_{\rm port}^{\rm CTC}
 =\sum_i\left[\Psi_i(v_i)+\Psi_i^*(\bar\eta_i)-\bar\eta_i(v_i)\right]
 =\frac12\sum_i\|v_i-\bar\eta_i^\sharp\|_{h_i}^2,
 \qquad \Psi_i(v)=\frac12\|v\|_{h_i}^2.
 \label{eq:p3bicotangent}
\end{equation}

\begin{proposition}[Minimal bicotangent port completion and selection boundary]
\label{prop:p3bicotangent}
In the local positive quadratic one-copy class with inherited metrics and one
common action unit, Eq.~\eqref{eq:p3bicotangent} is the unique completion that
has the exact conjugate graphs $v_i=\bar\eta_i^\sharp$, zero minimized source
remainder and no pair-dependent gain. Its two source/response exchanges are
commuting isometric involutions with one-dimensional odd contrast lines. If
the source pullbacks $\eta_i^{\rm src}=\bar\eta_iq_{\rm rel}$ are nonzero and
the declared variance representation reaches the two role rays, those lines
give the descriptor of Proposition~\ref{prop:p3portrepresentation}.

The frozen current base--seed/body reduct does not entail this completion or
physical occupation of both odd lines.
\end{proposition}

\begin{proof}
Fenchel--Young gives each term in Eq.~\eqref{eq:p3bicotangent} as a
nonnegative square, with equality exactly at $v_i=\bar\eta_i^\sharp$. In a
metric-normalized one-mode chart its joint Hessian is
\begin{equation}
 H_i=\begin{pmatrix}1&-1\\-1&1\end{pmatrix},
\end{equation}
whose response Schur complement is zero. For a general normalized quadratic
candidate with coefficients $(a_i,b_i,c_i)$, exact conjugacy gives
$c_i=a_i$, zero Schur remainder gives $b_i=c_i^2/a_i=a_i$, and the inherited
vertical curvature gives $a_i=1$. Hence the completion is unique in the
declared class. Leg exchange is an isometry of $H_i$; the two exchanges act
on disjoint blocks and commute, and their odd lines are generated by the two
source--response differences. Source basicness is precisely the quotient
descent condition, so Proposition~\ref{prop:p3portrepresentation} applies.

For non-entailment, the hub's cotangent-role control contains two models with
the same base, universal seed, stabilized body, protected covector, metric
and action unit: one omits the conjugate response and one admits it. Their
response ranks differ. Independently, an endpoint-directed source and an
oriented double-cover presentation can share the same local functional, and
the endpoint $\mathbb Z_2^{OS}$ is not the P3 $O/Q$ swap without a typed
intertwiner. Products of these controls have identical frozen reducts but
different ordered-port content. The current reduct therefore cannot select
the bicotangent packet. This is an unrestricted parent-law non-entailment, not a defect
of the sufficient completion.
\end{proof}

Thus standard Born and Hilbert-stress recovery select
Eq.~\eqref{eq:p3bicotangent} as the weakest source-preserving completion in
the declared local class. They do not prove that the earlier physical
base--seed reduct admits and occupies both quotient-basic covectors or owns
their variance representation on the P3 rays. No further port coefficient or
functional-shape search is licensed; the open question is cotangent-role and
typed-port admission by the enriched parent.

The hub's exterior and record results reduce the arbitrariness of that handoff.
Let
\begin{equation}
\mathcal F_{P3}=\Lambda^\bullet E_{P3}\otimes\mathcal R_2,
\qquad \dim\mathcal R_2=2,
\label{eq:p3fockrecord}
\end{equation}
where $\mathcal R_2$ is the already selected two-leg dagger--Frobenius record
carrier.  On $\Lambda^\bullet E_{P3}$ define
\begin{equation}
\mathsf C_{\rm ext}(v)=\varepsilon(v)+\iota_{K_*(v,\cdot)}.
\label{eq:p3extclifford}
\end{equation}
The canonical exterior identities give
$\mathsf C_{\rm ext}(v)^2=K_*(v,v)I$, and tensoring with
$I_{\mathcal R_2}$ leaves this law unchanged.  The exterior carrier has
dimension $2^7=128$ with equal central sectors of dimension $64$; the record
leg doubles them to $(128,128)$.  Since each simple $M_8(\mathbb C)$ module has
dimension eight, the resulting multiplicities are $(16,16)$ and the normalized
trace weights are $(1/2,1/2)$.

This is a functorial, coefficient-free realization of the premise of
Proposition~\ref{prop:p3clifford}; it does not identify the seven Paper-IV
dependency hyperedges with the seven P3 roles.  Its physical premise is instead
joint source ownership and occupation of the already typed $E_{P3}$ role
quotient, its positive metric and the two-leg record carrier.  Those ingredients
exist inside the declared source-irredundant enriched MVP.  More precisely,
if $r_*=\Delta_\tau e_*$ is its normalized occupied Frobenius record state,
then
\begin{equation}
J_{P3,R}\psi=\psi\otimes r_*,\qquad
J_{P3,R}^\dagger J_{P3,R}=I,
\label{eq:p3recordlift}
\end{equation}
and
$(\mathsf C_{\rm ext}(v)\otimes I)J_{P3,R}
=J_{P3,R}\mathsf C_{\rm ext}(v)$.
Thus the common-source lift cannot annihilate occupied P3 support and adds no
gain: joint assembly is closed inside that declared MVP carrier class.  This
conclusion is not a consequence of marginal occupation alone.  On
$\mathbb C^2$, the orthogonal projectors
$P=\lvert0\rangle\!\langle0\rvert$ and
$R=\lvert1\rangle\!\langle1\rvert$ both act nontrivially on
$(\lvert0\rangle+\lvert1\rangle)/\sqrt2$, while $PR=0$.
Consequently the unrestricted base--seed reduct still does not imply that its
physical parent law realizes and occupies the common enriched MVP carrier, but no additional internal
joint-occupation gate remains after that carrier is fixed.

The same packet removes the variance-representation premise left explicit in
Proposition~\ref{prop:p3bicotangent}.  Define
\begin{equation}
 \Phi_{P3}(v)=\mathsf C_{\rm ext}(v)\otimes I_{\mathcal R_2},
 \qquad
 \omega_{P3,R}(A)={1\over256}\operatorname{Tr}(A).
 \label{eq:p3variancerepresentation}
\end{equation}

\begin{proposition}[Canonical port variance and operational occupation]
\label{prop:p3portoccupation}
On the occupied typed P3 exterior--record packet,
\begin{equation}
 \omega_{P3,R}(\Phi_{P3}(v))=0,
 \qquad
 \operatorname{Re}\omega_{P3,R}
 \left(\Phi_{P3}(u)^*\Phi_{P3}(v)\right)=K_*(u,v).
 \label{eq:p3variancegram}
\end{equation}
Hence $\Phi_{P3}$ is a faithful metric-isometric variance representation and
its restrictions to the $O/Q$ and $M/G$ pair subspaces supply the
representation assumed in Proposition~\ref{prop:p3bicotangent} without a new
gain or carrier.

Assume in addition complete degree-one factorization through the typed P3
quotient.  A normalized POVM outcome family induces a nonzero $O/Q$ first
jet.  On every test region where the renormalized matter stress is a nonzero
distribution, metric variation induces a nonzero $M/G$ first jet.  Their
quotient covectors are basic, and the positive Riesz maps give unique nonzero
bicotangent responses.  Thus no further internal port-occupation bit is
required on such support.
\end{proposition}

\begin{proof}
The operators exchange the even and odd exterior grades, so their normalized
trace is zero.  The exterior Clifford relation gives
\begin{equation}
 \Phi_{P3}(u)\Phi_{P3}(v)+\Phi_{P3}(v)\Phi_{P3}(u)
 =2K_*(u,v)I.
\end{equation}
The operators are self-adjoint; normalized-trace cyclicity proves
Eq.~\eqref{eq:p3variancegram}.  Setting $u=v$ gives
$\|\Phi_{P3}(v)\|_{2,\omega}^2=K_*(v,v)$, so positivity of $K_*$ proves
injectivity and nonzero images for all typed role rays.

For the state-dependent claim, POVM normalization gives
$\sum_a\operatorname{Tr}(\rho_OE_a)=1$, so at least one typed effect-source
derivative is nonzero.  A nonzero stress distribution is exactly a nonzero
linear functional on compactly supported metric test variations, hence some
$h$ has
$-\frac12\int\sqrt{-g}\,T^{\mu\nu}h_{\mu\nu}\ne0$.
Complete typed factorization descends both jets to nonzero quotient
covectors; their source pullbacks annihilate the quotient kernel.  Since a
positive metric has an invertible Riesz map, both conjugate responses are
nonzero. QED.
\end{proof}

The nonzero-stress hypothesis is necessary.  An identically zero
renormalized stress on a region gives a zero $M/G$ first jet even though the
abstract $M$ and $G$ role rays and their faithful representation remain
present.  Proposition~\ref{prop:p3portoccupation} therefore closes an
internal representation/activation issue of the already occupied enriched
MVP; it does not prove that the unrestricted physical base--seed parent
selects the typed quotient, the Clifford--record packet or a nonzero-stress
observer support.

The same construction fixes the representation class rather than merely its
dimension.  Since
$\operatorname{Cl}_7(\mathbb C)\simeq M_8(\mathbb C)\oplus M_8(\mathbb C)$
has dimension 128, its left regular module contains eight copies of each
simple eight-dimensional module.  These are exactly the multiplicities of
$\Lambda^\bullet E_{P3}$.  The exterior vacuum has full Clifford-word orbit
rank 128; after the two central record projectors are included, the joint
orbit rank is 256.  Hence the exterior--record packet is a faithful cyclic
regular representation of its P3--record algebra.  Its represented normalized
trace fixes the equal Clifford-sector weights, so no separate cyclicity, GNS
carrier or central-weight choice remains on this subpacket.

\section{Complete-incidence characterization}
\label{app:p3certificate}

This does not
establish minimal-GNS faithfulness or differential zero cokernel for the full
$S0$--$S7$ source signature.  The positive connected extension
$\Gamma_{\rm ext}(z,w)=\Gamma_{P3,R}(z)+h\lVert w-bz\rVert^2/2$ has a
generically nonzero raw coordinate while exact elimination gives
$\operatorname{Schur}_wH_{\rm ext}=H_{P3,R}$.  It therefore disproves raw
coordinate exhaustivity.  It is not, however, a distinct physical-source
completion under the adopted source ontology: in the defect coordinate
$u=w-bz$ it is the centered positive auxiliary
$\Gamma_{P3,R}(z)+h\lVert u\rVert^2/2$ with $u_*=0$, zero independent load and
Schur-neutral elimination.  Physical exhaustivity modulo presentation
auxiliaries is instead reduced to uniqueness of the occupied stable branch
and zero cokernel of the complete typed source incidence.  Those two premises
remain open for the unrestricted physical base--seed parent.  On the declared
source-frozen finite MVP action, strict positivity already supplies the unique
stationary branch, so its sole remaining internal exhaustivity test is the
typed zero-cokernel condition.  Physical base--seed realization and occupation
of that MVP action remain
an upstream physical assumption.

The preceding construction can nevertheless be compressed into one exact
source certificate.  Let $P_{P3}^1$ be the physical degree-one source
quotient, with its declared seven typed roles, and let
\begin{equation}
 \mathcal C_{P3}^{\rm phys}:P_{P3}^1\longrightarrow
 \mathcal A_{\rm body}^{\rm sa}
 \label{eq:p3physicalcertificate}
\end{equation}
be a source-owned incidence into the self-adjoint part of the occupied body
algebra.  Call Eq.~\eqref{eq:p3physicalcertificate} \emph{complete} when
(i) its quotient rank is seven and its kernel is exactly the complete P3
physical null, with no terminal-visible bypass; (ii) after polar
normalization it obeys
\begin{equation}
 \mathcal C(v)^2=K_*(v,v)I;
 \label{eq:p3physicalsquare}
\end{equation}
and (iii) the canonical exterior extension and the two-leg
dagger--Frobenius record algebra used in Eq.~\eqref{eq:p3recordlift} are
jointly occupied on the same source carrier by a cyclic vector with full
algebra orbit (rank 256 in the present finite realization).
Here \emph{minimal} means generated by the physical degree-one source
quotient with no additional P3 bypass; it does not mean smallest Hilbert-space
dimension.  The declared recovery class uses the canonical cyclic regular
exterior--record realization.

\begin{theorem}[Minimal physical P3 certificate]
\label{thm:p3physicalcertificate}
Within the finite source-irredundant Clifford--recovery class, the physical
parent contains the minimal faithful typed P3 Clifford--record subpacket,
uniquely up to typed unitary equivalence, if and only if it owns a complete
incidence \eqref{eq:p3physicalcertificate}.
\end{theorem}

\begin{proof}
For sufficiency, the exact-null condition descends the incidence faithfully
to the seven-dimensional typed quotient.  Polar normalization fixes the
positive metric and Eq.~\eqref{eq:p3physicalsquare}; polarization then gives
the Clifford anticommutator.  The exterior construction
\eqref{eq:p3extclifford}, the occupied record lift
\eqref{eq:p3recordlift}, and Proposition~\ref{prop:p3portoccupation} provide
the faithful cyclic Clifford--record representation and its normalized
trace.  Source irredundancy removes invisible degree-one copies, while
complete-null/no-bypass removes a second terminal-visible realization.
The remaining freedom is therefore typed unitary equivalence.

Conversely, restrict a minimal faithful typed P3 Clifford--record subpacket
to its grade-one self-adjoint Clifford generators and pull them back along
the physical source first jet.  Faithfulness gives rank seven, the physical
quotient gives the exact kernel and no-bypass property, the Clifford law
gives Eq.~\eqref{eq:p3physicalsquare}, and regular cyclic occupation supplies
the same full exterior--record orbit.  Hence the pulled-back incidence is
complete. QED.
\end{proof}
\section{Enriched variational ownership and ambient-extension control}
\label{app:p3ownership}

The ownership question can be closed inside an explicit ordered enriched
completion without fixing the Clifford map in advance.  After the
morphological body has been solved and frozen, vary
$C\in\operatorname{Hom}_{\mathbb R}(E_{P3}^{\rm phys},
\mathcal A_{P3}^{\rm sa})$ and one operator response $S_{P3}$.  For the
source-frozen occupied descriptor $e_{P3}$ define
\begin{align}
 \Gamma_{P3}^{\rm enr}[C,S_{P3};e_{P3}]
 ={}&{\mathcal A_*\over2}\tau\!\left[
 (S_{P3}-C(e_{P3}))^\dagger(S_{P3}-C(e_{P3}))\right]
 \nonumber\\
 &+{\mathcal A_*\over2}\int_{S(E_{P3},K_*)}
 \tau\!\left[\mathfrak R_C(v)^\dagger\mathfrak R_C(v)\right]d\mu(v),
 \qquad
 \mathfrak R_C(v)=C(v)^2-K_*(v,v)I.
 \label{eq:p3enrichedparentaction}
\end{align}
This is a staged parent action: body variables are not varied in the
post-body P3 stage.  The two rows use the same faithful trace and action unit,
and the second counts the primitive operator relation once.

\begin{theorem}[Enriched-parent variational ownership]
\label{thm:p3enrichedownership}
In the declared source-irredundant enriched completion, supplemented by the
complete physical P3 null, cyclic exterior--record occupation and no-bypass
conditions, Eq.~\eqref{eq:p3enrichedparentaction} owns and selects the minimal
faithful P3 Clifford--record subpacket.  Its zero locus is
\begin{equation}
 S_{P3}=C(e_{P3}),
 \qquad
 C(u)C(v)+C(v)C(u)=2K_*(u,v)I,
 \label{eq:p3enrichedzerolocus}
\end{equation}
and exact elimination of $S_{P3}$ leaves
$\mathcal V_{\rm OS}[C]$.  Hence the handoff and source selector are not
double-counted.
\end{theorem}

\begin{proof}
Both terms in Eq.~\eqref{eq:p3enrichedparentaction} are nonnegative.  The
canonical exterior Clifford incidence together with
$S_{P3}=C(e_{P3})$ reaches zero.  Conversely, every global zero has vanishing
graph defect and vanishing operator-square integral.  The first gives the
response equation in Eq.~\eqref{eq:p3enrichedzerolocus}; faithfulness,
continuity and homogeneity give
$C(v)^2=K_*(v,v)I$ for every $v$, and polarization gives the displayed
anticommutator.  Minimization over $S_{P3}$ removes exactly the graph square
and leaves Eq.~\eqref{eq:p3operatorsquarefunctional}.  Theorem
\ref{thm:p3physicalcertificate} then supplies typed-unitary uniqueness of the
minimal occupied subpacket. QED.
\end{proof}

The construction is not an arbitrary terminal repair.  Its typed domain is
the already fixed P3 quotient; its target is the canonical generated exterior
algebra; its metric, trace and unit are inherited; body-first staging forbids
readout-to-body backreaction; and it replaces rather than duplicates the
fixed graph-normal handoff.  Numerically, the graph term minimizes to zero
for both the Clifford and commuting controls, while the reduced enriched
costs remain $0$ and $4/3$, respectively; the zero incidence has normalized
cost one.

This is a conditional source-law closure, not an unrestricted
parent-realization theorem.
The current reduct has no source variable typed as a complete rank-seven
operator incidence.  It admits the equal-Gram controls in
Eq.~\eqref{eq:p3operatorsquarecontrol}, so it cannot entail the selector.
Equation~\eqref{eq:p3operatorsquarefunctional} is also not claimed unique
among all possible positive parent functionals; it is the coefficient-free
representative of the declared one-copy orthogonally averaged squared-defect
class.  The theorem proves ownership by the explicitly declared enriched
parent, not that the earlier reduct physically realizes and occupies that
enrichment.

Theorem~\ref{thm:p3physicalcertificate} is a characterization theorem, not a
derivation of certificate existence from the current base--seed reduct.  Nor
can finite readouts establish that this selected subpacket exhausts the
ambient parent.  Indeed, for any visible action $\Gamma(x)$ and finite
terminal stack $F(x)$, set
\begin{equation}
 \widetilde\Gamma(x,w)=\Gamma(x)
 +{q\over2}(w-bx)^2-j(w-bx),\qquad
 \widetilde F(x,w)=F(x),
 \quad q>0,\ j\ne0.
 \label{eq:p3loadedhiddenextension}
\end{equation}
At every visible stationary point $x_*$, the extension has
$w_*=bx_*+j/q$, preserves $F(x_*)$, and satisfies
\begin{equation}
 \operatorname{Schur}_w H_{\widetilde\Gamma}=H_\Gamma.
 \label{eq:p3loadedhiddenschur}
\end{equation}
Thus the finite terminal family cannot distinguish the original parent from
one with an additional independently loaded source component.  In the
invariant defect coordinate $u=w-bx$, Eq.~\eqref{eq:p3loadedhiddenextension}
is $\Gamma(x)+qu^2/2-ju$: it is not an unloaded presentation auxiliary, but
it is also not a second interacting realization of the same connected
universe packet.  The selected connected packet is exhaustive only after the
independent parent-side condition
$\operatorname{coker}J_{\rm phys}=0$ (or an equivalent source law) is proved.

\section{Finite mixed-Gram operational witness}
\label{app:p3rwwitness}

The remaining cokernel is not distributed over all named terminals.  Existing
hub theorems make the P1 gravity--time--radiation/matter differential core and
the finite quantum/P3 restriction separately onto on their declared working
packets.  The unresolved datum is their typed gluing: either
$\ker J_Q\subseteq\ker J_{P1}$ and a source-owned intertwiner
$J_{P1}=\iota_{P1}J_Q$ places P1 inside the complete quantum descriptor, or an
independent P1 factor requires a nonseparable joint incidence.  The relevant
finite calculation is therefore the ROOT-to-QFT preparation/stress
cross-Jacobian, not an additional terminal rank counted by name.
Equation~\eqref{eq:rootp1factor} closes its genuinely quantum stress component
on the declared same-state two-mode sector, and
Eq.~\eqref{eq:rootp1tangentfactor} closes the full P1 restriction on the pure
post-body ROOT tangent.  Equations~\eqref{eq:stagedglue}--
\eqref{eq:stagedgluerank} eliminate that tangent exactly and reduce the
residual kernel test to one operator on the complementary source domain.
The controls in Eq.~\eqref{eq:restcontrols} prove that present marginal data
do not select its rank. They do not decide whether those remaining directions are basic on the
complete quantum descriptor or belong to an additional jointly occupied P1
factor.

The preceding ranks belong to the source-minimal primitive carrier.  They
must not be identified with the row rank of an observer instrument evaluated
on a larger morphology tangent.  The latter can contain several endpoint
functionals of the same primitive role without creating additional parent
source directions.  The following finite witness makes that distinction
explicit and also proves that the mixed-Gram increment can be nonzero.

\subsubsection{Finite operational witness}

\begin{proposition}[Finite protected Robertson--Walker witness]
\label{prop:rwwitness}
On the endpoint-normalized eight-mode homogeneous tangent
\begin{equation}
e_n(t)=t^n-1,\qquad n=1,\ldots,8,
\label{eq:rwbasis}
\end{equation}
take three comoving source times
$t_s\in\{0.21,0.47,0.79\}$ and the metric P1 rows
\begin{align}
Z_s[f]&=-f(t_s),\nonumber\\
J_s[f]&=(1-t_s)f(t_s)-\int_{t_s}^{1}f(t)\,dt,\nonumber\\
E[f]&=-\ddot f(1).
\label{eq:rwmetricrows}
\end{align}
Let the two real ROOT rows be the real and imaginary parts of
\begin{equation}
q_\omega[f]=\int_0^1 e^{-2i\omega t}\dot f(t)\,dt,
\qquad \omega=2.35.
\label{eq:rwquantumrow}
\end{equation}
On the protected quadratic body branch the direct Green-clock third-jet row
vanishes.  The resulting finite matrices satisfy
\begin{equation}
\operatorname{rank}A_{P1}^{RW}=7,
\qquad
\operatorname{rank}\begin{pmatrix}A_{P1}^{RW}\\
\Re q_\omega\\ \Im q_\omega\end{pmatrix}=8,
\qquad r_{Q\perp}^{RW}=1.
\label{eq:rwranks}
\end{equation}
Hence a common-source quantum row can detect one morphology direction erased
by this finite geometric observer packet.
\end{proposition}

\begin{proof}
Substitution of Eq.~\eqref{eq:rwbasis} into
Eq.~\eqref{eq:rwmetricrows} gives three redshift rows, three isotropic Jacobi
rows and one tidal row.  Direct singular-value evaluation gives rank seven
and a one-dimensional null.  Numerical quadrature of
Eq.~\eqref{eq:rwquantumrow} gives a nonzero norm
$5.673092\times10^{-4}$ after projection to this null, so stacking its real
and imaginary parts raises the rank to eight.  The calculation is regenerated
by \path{scripts/audit_theorem_certificates.py}.
\end{proof}

This is a source-matched positive operational witness, not unrestricted
parent-law realization.  It is resolution dependent: sufficiently complete homogeneous
geometric endpoint data can already become injective, in which case the same
quantum functional adds zero rank.  The one-dimensional unresolved fiber is
also not by itself a Hilbert space or a proof of Born structure.  Finally,
Eq.~\eqref{eq:rwranks} does not contradict Eq.~\eqref{eq:deduprank}: the
former is an observer-operator rank on an eight-mode morphology tangent,
whereas the latter is the de-duplicated primitive source-carrier rank.

\section{Class-minimal support--port enriched-law certificate}
\label{app:p3minimalparentlaw}

The support--port reconstruction separates the data that define the proposed
physical law from the conclusions generated by that law.  The support
projectors and common-unit character specify the admissible source domain;
the cotangent, operator-square and record terms specify one staged
once-counted functional; and the decoder specifies noncollapsing body
realization.  They are not independent terminal fits.

\begin{proposition}[Variational occupation, class-relative componentwise minimality and narrow-reduct no-go]
\label{prop:p3minimalparentcertificate}
In the finite source-irredundant one-copy class used in
Theorem~\ref{thm:p3supportportparent}, the normalized P3 functional is
nonnegative and has an explicit simultaneous zero.  Consequently its global
minimum is zero and every global minimizer
lies on the complete packet zero locus; no separate P3 common-zero occupation
assumption is required.  Moreover, every clause of the proposed P3 law is
required for that theorem's conclusion in the following precise sense:
\begin{enumerate}[label=(\roman*),nosep]
\item deleting the support flag leaves the B/R singleton blocks and O/Q,
  M/G planes unidentified;
\item deleting either ordered cotangent port leaves the corresponding
  nontrivial pair swap in the typed stabilizer;
\item placing $z_U$ in the relative span lowers the source rank to six even
  at fixed total row cost;
\item deleting the left-decoder premise permits a rank-six body realization;
\item deleting the operator-square term admits a commuting packet with the
  same normalized marginal trace Gram;
\item deleting the pointed root admits a faithful full-rank nontracial orbit;
  and
\item deleting cyclicity or no bypass permits unoccupied or additional record
  support.
\end{enumerate}
The two cotangent response legs and the graph-normal P3 response can be
eliminated exactly with zero additional source stiffness, so their inclusion
does not duplicate the retained operator-square law.  The unchanged narrow
base--seed reduct admits the first, third and fifth controls with the same
declared marginal upstream data and therefore does not entail
$\mathcal L_{P3}^{\min}$.
\end{proposition}

\begin{proof}
The shifted stabilized-body term is nonnegative on its declared basin.  Each
cotangent contribution is a Fenchel--Young square; the P3 graph and
operator-square contributions are squared faithful-trace norms; and the
record contribution is the normalized squared record defect.  Choose the
canonical exterior Clifford incidence, set its graph response equal to its
source evaluation, use the two metric-dual ordered ports, the stabilized body
and its exact left decoder, and take the occupied empty exterior word with
the normalized two-leg record root.  Every displayed defect then vanishes
simultaneously.  Hence the action floor is zero.  A sum of nonnegative terms
can attain that floor only when every term vanishes, proving variational
occupation of the packet.

The support stabilizer is
$S_{\{O,Q\}}\times S_{\{M,G\}}$.  One port removes one factor and two ports
remove both.  The independent/dependent unit controls have ranks seven and
six with equal total squared row norm.  The redundant body realization has a
left inverse, whereas identifying its unit column with another column lowers
the rank and makes an exact left inverse impossible.  Clifford generators and
the commuting control have the same identity trace Gram, but only the former
have a state-neutral square and vanishing mixed anticommutators.  The
positive noncentral root-density control has full orbit rank and equal central
weights but a nonidentity Gram.  The cyclic/no-bypass controls are the
complete-incidence and pointed-orbit countermodels proved above.  Finally,
each normalized Fenchel--Young port has zero source Schur remainder, and
minimizing the graph square in Eq.~\eqref{eq:p3enrichedparentaction} leaves
the operator-square functional exactly.  These facts prove all clauses and
the same-reduct non-entailment.  QED.
\end{proof}

The proposition is a class-relative minimality certificate for the stated
theorem, not a claim that no richer physical parent law can realize the same
packet.  The superscript ``min'' denotes only componentwise minimality in the
declared finite, source-irredundant one-copy class.  Under the four-dimensional
recovery-completion policy and within that class, the O/Q and M/G dualities,
their de-duplicated relation supports, the common unit and body noncollapse
are the weakest preimages of the already declared standard recovery targets.
The state-neutral operator-square norm is the genuinely new Tau-specific
cross-role clause.  The pointed cyclic root is instead the minimal
record-side response, within the class, to the faithful full-rank nontracial
countermodel.  Hence the paper no longer leaves an unspecified
``complete rank-seven incidence'' as a free input: it proposes the explicit
law in Eq.~\eqref{eq:p3minimalparentlaw}, derives the packet from it, and
derives occupation from global minimization rather than assuming its common
zero separately.  Its
empirical truth and unrestricted ambient-source exhaustivity remain separate
questions.


\section{Self-contained dependency certificate}
\label{app:dependencycertificate}

This appendix separates logical premises from provenance.  The conclusions of
Paper VII use only the following statements.  A reader may treat each imported
statement as an explicit assumption and verify every theorem in this paper
without consulting another Tau manuscript.

\begin{description}[style=nextline,leftmargin=0pt]
\item[\textbf{D1: body-first order.}]
There is a $C^2$ pre-body action $\mathcal A_{\rm pre}[y;B,s]$ with an occupied
stationary point $y_*$ and positive physical Hessian $H_B$.  The post-body
functional is evaluated on the frozen parameter $y_*$; post-body source
variation is not varied back into the body equation.  This is the only
body-first premise used in Eqs.~\eqref{eq:jointfunctional} and
\eqref{eq:triangular}.

\item[\textbf{D2: common typed source domain.}]
On the selected body and observer support there is one finite source domain
$\mathcal E_{*,O}$, one access map $A_O$, typed linear role maps
$R_a:\mathcal E_{*,O}\to\mathcal Y_a$, and stable resolution maps
$Q_O^{(a)}$.  No direct-sum decomposition of $\mathcal E_{*,O}$ is assumed.
Proposition~\ref{prop:moprassembly} is proved here from exactly this premise.

\item[\textbf{D3: occupied operation data.}]
An operation algebra $\mathfrak A_{\rm src}$, occupied vector $\Omega_s$ and
positive root covariance $X_0$ are given.  Whenever the full-occupation branch
is used, $\Omega_s$ is cyclic and every required $R_a$ is surjective.
Whenever the rooted-path branch is used,
$\operatorname{supp}X_0\subseteq V_s$ and every admitted $L_e$ preserves
$V_s=\operatorname{span}(\mathfrak A_{\rm src}\Omega_s)$.
Proposition~\ref{prop:occupation} proves the resulting full-occupation or
activation statement; physical realization of either premise set is not
imported.

\item[\textbf{D4: fixed quantum preparation.}]
The CTP boundary datum satisfies $\rho_{\rm in,O}\geq0$ and
$\operatorname{Tr}\rho_{\rm in,O}=1$.  It is held fixed while $g$, $A$ and
$\Phi$ are varied.  The CTP evolution and observer restriction produce a
normalized output state $\rho_O$.  CTP response is obtained by varying the
difference source before taking the physical equal-branch limit.  No Euler
variation with respect to a density matrix is used.

\item[\textbf{D5: standard Lorentzian leaf.}]
On the selected leaf the geometric and Abelian functionals satisfy the
classification hypotheses of Proposition~\ref{prop:standardforms}, with a
common unit packet.  Their Einstein--Hilbert and Maxwell shapes then follow
inside that class.  The matter functional, measure and renormalization
prescription are diffeomorphism and $U(1)$ gauge invariant and anomaly-free
for the Ward identities invoked here.  This is an assumed standard-limit
leaf, not a Tau derivation of those actions.

\item[\textbf{D6: operational effect source.}]
Each declared outcome family is a POVM, $E_a\geq0$ and $\sum_aE_a=I$, and its
source coupling is Eq.~\eqref{eq:effectcoupling}.  This premise, together with
D4, is sufficient for the normalized Born derivative in
Eq.~\eqref{eq:sourceowners1}; no stronger quantum reconstruction theorem is
used.

\item[\textbf{D7: clock-source boundary.}]
The clock character $\vartheta$ is an admitted external source coordinate of
the post-body functional.  Paper VII uses only the formal response derivative
$\delta W_O^{\rm ren}/\delta\vartheta$.  It neither imports nor proves a
duration law, global clock integrability or a Tau-specific clock distortion.

\item[\textbf{D8: once-counted calibration.}]
Every displayed sector uses one declared action unit, each occupied
contribution and counterterm has one ledger owner, and the overlap
$K_{m,m+2}$ is calculated from Eq.~\eqref{eq:overlapformula} and frozen before
the ROOT--TT comparison.  This bookkeeping premise forbids an additional
fitted gain but does not select the common unit, occupied mode pair, screen or
subtraction prescription.

\item[\textbf{D9: component-local closure scope.}]
Whenever Theorem~\ref{thm:taustandardcomposition} is invoked, the two
occupied supports are parent relation-free; every primitive closure rule has
connected support, preserves its component support, and the primitive rules
exhaust joint admissibility without a hidden global endpoint predicate.  The
unrescaled metric conclusion additionally assumes the source-rigid monoidal
normalization LC4.  These premises are stated explicitly in the theorem and
are not inferred from a vanishing terminal residual or from complete local
marginals.

\item[\textbf{D10: joint locality/P3 completion scope.}]
Theorem~\ref{thm:jointlocalp3} uses one action that is additive only across
connected primitive-incidence components.  It does not factorize the seven
internal P3 roles.  Its P3 component is the already declared
$\mathcal L_{P3}^{\min}$, and all component zeros are assumed attainable.
The theorem proves class-relative law-to-occupation and exact non-entailment
from the earlier narrow reduct; it does not prove Nature-level realization or
ambient-parent exhaustivity.
\end{description}

\subsection*{Provenance and verification locators}

\begin{center}
\small
\begin{tabularx}{0.96\textwidth}{@{}p{0.16\textwidth}p{0.32\textwidth}X@{}}
\toprule
Companion source & Detailed locator & Logical use in Paper VII\\
\midrule
Paper I & \emph{Body-first variational order} and \emph{One common post-body
functional} & Provenance for D1; D1 is fully restated above\\
Paper II & \emph{Finite typed source signatures} and \emph{Cyclic
generated-support criterion} & Provenance for D2--D3; the required
consequences are reproved in
Propositions~\ref{prop:moprassembly}--\ref{prop:occupation}\\
Paper III & \emph{Recovery domains} and \emph{Recovery sandwich} &
Architectural context for stable resolution; no recovery theorem is required
here\\
Paper IV & \emph{Common-source primitive-readout assembly} and
\emph{Current-reduct selection no-go} & Motivation; the relevant assembly and
counterpair are independently stated and proved here\\
Paper V & \emph{Observer clock form} and \emph{Pathwise VEP clock} &
Provenance for the symbol $\vartheta$; no Paper-V temporal theorem is used\\
Paper VI & \emph{Quantum source packet} and \emph{Born and interference
terminal} & Provenance for D4 and D6; the effect derivative needed here is
restated explicitly\\
P1/P3 hub packets & gravity--time--radiation ledger and Abelian handoff &
Provenance for D5 and $Z_F=p_*K_Y$; both enter here as explicit assumptions\\
\bottomrule
\end{tabularx}
\end{center}

Thus rejection or later revision of a companion result changes the
interpretation of the corresponding D-item, but creates no hidden step in the
conditional mathematics of this manuscript.  Unrestricted parent-law
realization and occupation of D1--D8 jointly remain outside the theorem set.

\begingroup
\small
\bibliographystyle{unsrtnat}
\bibliography{refs}
\endgroup
\end{document}
